% =====================================================================
% hiemalis.tex -- LaTeX conversion of hiemalis.md (Breviarium Romanum,
% Pars Hiemalis), BODY CONTENT ONLY.
%
% No preamble, \documentclass, or \newcommand definitions are included
% here, per instruction -- the command vocabulary below is to be
% defined locally. See the command legend delivered alongside this
% file for what each command means and how it should be styled
% (one-column vs. two-column, red vs. black, etc.).
%
% BATCH 1 of N: Ordinarium divini Officii (hiemalis.md lines 1-1394)
% =====================================================================

% [ornamental woodcut headpiece]
\partheading{Ordinarium divini Officii}
\partsubtitle{Juxta Ritum Romanum Persolvendi}

\rubric{Clericis in sacris constitutis, qui integrum divinum Officium, quamvis in partes distributum, coram Ssmo Sacramento, sive publicæ adorationi exposito sive in tabernaculo adservato, devote recitaverint, conceditur: Indulgentia plenaria, si præterea sacramentalem confessionem instituerint, ad eucharisticum Convivium accesserint et ad mentem Summi Pontificis oraverint. \textit{Pius Pp. XI, 23 Octobris 1930.}}

\rubric{Iis vero, qui divinum Officium coram Ssmo Sacramento, ut supra, ex parte tantum recitaverint, conceditur: Indulgentia quingentorum dierum pro unaquaque hora canonica. \textit{Pius Pp. XI, 18 Maji 1933.}}

\rubric{Iisdem autem clericis, qui divinum Officium in alias preces commutatum rite obtinuerint, si preces ejusmodi coram Ssmo Sacramento, uti supra, devote recitaverint, conceditur: Indulgentia plenaria, additis sacramentali confessione, sacra Communione et oratione ad mentem Summi Pontificis. \textit{Pius Pp. XI, 7 Novembris 1932.}}

\sectionheading{Ante divinum Officium}

\rubric{Antequam inchoëtur Officium laudabiliter dicitur, sub singulari semper numero, sequens Oratio; pro qua Summus Pontifex Pius Pp. XI sub die 17 Novembris 1933 Indulgentiam trium annorum concessit.}

\oratioheading{Oratio}

\lettrine[lines=2]{A}{PERI}, Dómine, os meum ad benedicéndum nomen sanctum tuum: munda quoque cor meum ab ómnibus vanis, perversis et aliénis cogitatiónibus; intelléctum illúmina, afféctum inflámma, ut digne, atténte ac devóte hoc Offícium recitáre váleam, et exaudíri mérear ante conspéctum divínæ Majestátis tuæ. Per Christum, Dóminum nostrum. \R\ Amen.

Dómine, in unióne illíus divínæ intentiónis, qua ipse in terris laudes Deo persolvísti, has tibi Horas (\textit{vel} hanc tibi Horam) persólvo.

% N° 28. -- H. 1 [printer's signature mark; omitted from running text]

\sectionornament
% <!-- p. 2 -->
\label{p:2}

\sectionheading{Ante singulas Horas}

\rubric{Ante Matutinum, sicut et ante Laudes a Matutino extra Chorum separatas, et ante omnes Horas, præterquam ad Completorium, dicitur secreto:}

\lettrine[lines=2]{P}{ATER} noster, qui es in cælis, sanctificétur nomen tuum. Advéniat regnum tuum. Fiat volúntas tua, sicut in cælo et in terra. Panem nostrum quotidiánum da nobis hódie. Et dimítte nobis débita nostra, sicut et nos dimíttimus debitóribus nostris. Et ne nos indúcas in tentatiónem: sed líbera nos a malo. Amen.

\lettrine[lines=2]{A}{VE}, María, grátia plena; Dóminus tecum: benedícta tu in muliéribus, et benedíctus fructus ventris tui Jesus. Sancta María, Mater Dei, ora pro nobis peccatóribus, nunc et in hora mortis nostræ. Amen.

\rubric{In principio autem Matutini ac Primæ additur secreto Symbolum Apostolorum:}

\lettrine[lines=2]{C}{REDO} in Deum, Patrem omnipoténtem, Creatórem cæli et terræ. Et in Jesum Christum Fílium ejus únicum, Dóminum nostrum: qui concéptus est de Spíritu Sancto, natus ex María Vírgine, passus sub Póntio Piláto, crucifíxus, mórtuus, et sepúltus: descéndit ad ínferos; tértia die resurréxit a mórtuis; ascéndit ad cælos; sedet ad déxteram Dei Patris omnipoténtis: inde ventúrus est judicáre vivos et mórtuos. Credo in Spíritum Sanctum, sanctam Ecclésiam cathólicam, Sanctórum communiónem, remissiónem peccatórum, carnis resurrectiónem, vitam ætérnam. Amen.

\rubric{Officium ad Matutinum et ad reliquas Horas, nisi aliter suo loco indicetur, per totum annum persolvitur juxta formam, quæ in Rubricis hujus Ordinarii exhibetur.}

\sectionornament
% <!-- p. 3 -->
\label{p:3}

\hourheading{Ad Matutinum}

Pater noster, Ave María et Credo.

\rubric{Deinde, clara voce, dicitur Versus:}

\V\ \lettrine[lines=2]{D}{OMINE}, lábia mea apéries. \R\ Et os meum annuntiábit laudem tuam. \V\ Deus, in adjutórium meum inténde. \R\ Dómine, ad adjuvándum me festína.

Glória Patri, et Fílio, et Spirítui Sancto. Sicut erat in princípio, et nunc, et semper, et in sæcula sæculórum. Amen. Allelúja.

\rubric{Sic dicitur Allelúja ad omnes Horas, præter quam a Completorio Sabbati ante Septuagesimam; tunc enim dicitur:}

Laus tibi, Dómine, Rex ætérnæ glóriæ.

\rubric{Postea dicitur conveniens Invitatorium, quod ante Psalmum bis integrum recitatur, et ad singulos ejusdem Psalmi versus vel integrum vel dimidiatum ab Asterisco * alternis vicibus, ut infra, repetitur.}

\psalmheading{Psalmus 94}

\lettrine[lines=2]{V}{ENITE}, exsultémus Dómino, jubilémus Deo, salutári nostro: præoccupémus fáciem ejus in confessióne, et in psalmis jubilémus ei.

\rubric{Repetitur integrum Invitatorium.}

Quóniam Deus magnus Dóminus, et Rex magnus super omnes deos: quóniam non repéllet Dóminus plebem suam: quia in manu ejus sunt omnes fines terræ, et altitúdines móntium ipse cónspicit.

\rubric{Repetitur altera pars Invitatorii.}

\rubric{In sequenti Psalmi versu, ad verba veníte, adorémus, et procidámus ante Deum, genuflectitur.}

Quóniam ipsíus est mare, et ipse fecit illud, et áridam fundavérunt manus ejus: veníte, adorémus, et procidámus ante Deum: plorémus coram Dómino, qui fecit nos, quia ipse est Dóminus, Deus noster; nos autem pópulus ejus, et oves páscuæ ejus.

\rubric{Repetitur integrum Invitatorium.}

Hódie, si vocem ejus audíeritis, nolíte obduráre corda vestra, sicut in exacerbatióne secúndum diem tentatiónis in desérto: ubi tentavérunt me patres vestri, probavérunt et vidérunt ópera mea.

\rubric{Repetitur altera pars Invitatorii.}

Quadragínta annis próximus fui generatióni huic, et dixi: Semper hi errant corde; ipsi vero non cognovérunt vias meas: quibus jurávi in ira mea: Si introíbunt in réquiem meam.

\rubric{Repetitur integrum Invitatorium.}

Glória Patri, et Fílio, et Spirítui Sancto. Sicut erat in princípio, et nunc, et semper, et in sæcula sæculórum. Amen.

\rubric{Repetitur altera pars Invitatorii.}

\rubric{Denique repetitur integrum Invitatorium.}

\rubric{Expleto Psalmo, dicitur Hymnus Invitatorio respondens.}

\variantheading{Per Annum}

\rubric{In Officio dominicali: In omnibus Dominicis, extra Octavam Nativitatis, Invitatorium et Hymnus ut in Psalterio.}

\rubric{In Officio feriali: In omnibus per Annum Feriis, et a Feria IV Cinerum usque ad Sabbatum sequens inclusive, Invitatorium et Hymnus de Feria currenti, ut in Psalterio.}

\variantheading{Tempore Adventus}

\rubric{A Dominica I usque ad Sabbatum ante Dominicam III inclusive, tam in dominicali quam in feriali Officio, dicitur quotidie sequens}

\invitatlabel{Invitat.} Regem ventúrum Dóminum, * Veníte, adorémus. \psalmlabel{Ps. 94} Veníte.

\rubric{A Dominica vero III usque ad ultimam diem ante Vigiliam Nativitatis Domini inclusive, tam in dominicali quam in feriali Officio, dicitur quotidie sequens}

\invitatlabel{Invitat.} Prope est jam Dóminus: * Veníte, adorémus. \psalmlabel{Ps. 94} Veníte.

\rubric{In Vigilia Nativitatis Domini, etiam in Dominicam IV incidenti, Invitatorium ut in Proprio de Tempore.}

\variantheading{Hymnus}

\lettrine[lines=2]{V}{ERBUM} supérnum, pródiens\\
E Patris ætérni sinu,\\
Qui natus orbi súbvenis,\\
Labénte cursu témporis:

Illúmina nunc péctora,\\
Tuóque amóre cóncrema;\\
Ut cor cadúca déserens\\
Cæli volúptas ímpleat.

Ut, cum tribúnal Júdicis\\
Damnábit igni nóxios,\\
Et vox amíca débitum\\
Vocábit ad cælum pios;

Non esca flammárum nigros\\
Volvámur inter túrbines,\\
Vultu Dei sed cómpotes\\
Cæli fruámur gáudiis.

Patri, simúlque Fílio,\\
Tibíque, Sancte Spíritus,\\
Sicut fuit, sit júgiter\\
Sæclum per omne glória.\\
Amen.

\rubric{Conclusio communis, in Hymno præcedenti, et in aliis idem metrum habentibus, semper omittitur, quando specialis in omnibus Horis adhibenda præscribitur; et, si plures Conclusiones propriæ occurrant, sumitur Conclusio Officii currentis, aut secus Officii ipsa die et primo quidem loco inter cetera propriam Conclusionem habentia commemorandi, aut demum de occurrenti Octava communi vel de Tempore. Conclusio tamen \rubricblack{Jesu, tibi sit glória, Qui natus es de Vírgine,} in Officiis de Tempore Adventus numquam adhibetur.}

\variantheading{In Festis}

\rubric{In Officio cujuslibet Festi vel Octavæ, ac sanctæ Mariæ in Sabbato, Invitatorium et Hymnus ut in Proprio vel Communi; in Vigilia autem Epiphaniæ et in Dominica, etiam translata, infra Octavam Nativitatis, si de eis fiat Officium, Invitatorium et Hymnus ut in Proprio de Tempore. Quæ norma ad reliquas item Horas erit observanda pro illis Officii partibus, quæ juxta Rubricas ex Ordinario vel Psalterio non fuerint in Festis desumendæ.}

\variantheading{In Officio novem Lectionum}

\rubric{Expleto Hymno, dicuntur Antiphonæ convenientes, quæ in Officiis ritus duplicis ante et post Psalmos integræ recitantur; in Officiis autem ritus semiduplicis initio Psalmi inchoantur tantum et usque ad Asteriscum * perducuntur, atque in fine integræ pronuntiantur; servata alterutra norma etiam in omnibus Antiphonis, quæ ad Laudes ac Vesperas adhibentur.}

\nocturnoheading{In I Nocturno}

\rubric{Sub congruentibus Antiphonis dicuntur tres Psalmi, ac deinde subjungitur Versus, prouti Officium occurrens requirit.}

\rubric{Post Versum cujuslibet Nocturni dicitur}

Pater noster \rubric{secreto usque ad} \V\ Et ne nos indúcas in tentatiónem. \R\ Sed líbera nos a malo.

\absolutioheading{Absolutio}

\lettrine[lines=2]{E}{XÁUDI}, Dómine Jesu Christe, preces servórum tuórum, et miserére nobis: Qui cum Patre et Spíritu Sancto vivis et regnas in sæcula sæculórum. \R\ Amen.

\V\ Jube, domne, benedícere.

\rubric{Extra Chorum, quando ab uno tantum recitatur Officium, ante singulas Lectiones Matutini atque ad Lectionem brevem Primæ et Completorii, dicitur: Jube, Dómine, benedícere; et subjungitur congruens Benedictio. Ab Episcopo autem, ultimam Matutini Lectionem cantaturo, item dicitur: Jube, Dómine, benedícere; et respondetur a Choro: Amen.}

\variantheading{Pro I Lectione}

\benedlabel{Bened.} Benedictióne perpétua benedícat nos Pater ætérnus. \R\ Amen.

\rubric{Deinde dicuntur in unoquoque Nocturno Lectiones, prouti Officium occurrens requirit, et in fine cujuslibet Lectionis additur:}

\V\ Tu autem, Dómine, miserére nobis. \R\ Deo grátias.

\rubric{Post quamlibet vero Lectionem, quæ Hymnum Te Deum immediate non præcedat, congruens dicitur Responsorium, et in fine ultimi Responsorii cujusque Nocturni additur Versus: Glória Patri, et Fílio, et Spirítui Sancto, et Responsorium ab Asterisco *, et quidem a secundo, si duo fuerint, repetitur.}

\variantheading{Pro ii Lectione}

\benedlabel{Bened.} Unigénitus Dei Fílius nos benedícere et adjuváre dignétur. \R\ Amen.

\variantheading{Pro iii Lectione}

\benedlabel{Bened.} Spíritus Sancti grátia illúminet sensus et corda nostra. \R\ Amen.

\nocturnoheading{In II Nocturno}

\rubric{Sub congruentibus item Antiphonis dicuntur tres Psalmi et Versus, sicut in I Nocturno.}

\rubric{Post Versum dicitur.}

Pater noster \rubric{secreto usque ad} \V\ Et ne nos indúcas in tentatiónem. \R\ Sed líbera nos a malo.

\absolutioheading{Absolutio}

\lettrine[lines=2]{I}{PSIUS} píetas et misericórdia nos ádjuvet, qui cum Patre et Spíritu Sancto vivit et regnat in sæcula sæculórum. \R\ Amen.

\V\ Jube, domne, benedícere.

\variantheading{Pro iv Lectione}

\benedlabel{Bened.} Deus Pater omnípotens sit nobis propítius et clemens. \R\ Amen.

\variantheading{Pro v Lectione}

\benedlabel{Bened.} Christus perpétuæ det nobis gáudia vitæ. \R\ Amen.

\variantheading{Pro vi Lectione}

\benedlabel{Bened.} Ignem sui amóris accéndat Deus in córdibus nostris. \R\ Amen.

\nocturnoheading{In III Nocturno}

\rubric{Sub congruentibus denique Antiphonis dicuntur tres Psalmi et Versus, sicut in I et II Nocturno.}

\rubric{Post Versum dicitur}

Pater noster \rubric{secreto usque ad} \V\ Et ne nos indúcas in tentatiónem. \R\ Sed líbera nos a malo.

\absolutioheading{Absolutio}

\lettrine[lines=2]{A}{} \textsc{vinculis} peccatórum nostrórum absólvat nos omnípotens et miséricors Dóminus. \R\ Amen.

\V\ Jube, domne, benedícere.

\variantheading{Pro vii Lectione}

\benedlabel{Bened.} Evangélica léctio sit nobis salus et protéctio. \R\ Amen.

\variantheading{Pro viii Lectione}

\rubric{In Festis Domini et in Dominicis:}

\benedlabel{Bened.} Divínum auxílium máneat semper nobíscum. \R\ Amen.

\rubric{In Festis beatæ Mariæ Virginis:}

\benedlabel{Bened.} Cujus festum cólimus, ipsa Virgo virgínum intercédat pro nobis ad Dóminum. \R\ Amen.

\rubric{In Festis Sanctorum:}

\benedlabel{Bened.} Cujus (\textit{vel} Quorum \textit{aut} Quarum) festum cólimus, ipse (\textit{vel} ipsa \textit{aut} ipsi \textit{vel} ipsæ) intercédat (\textit{vel} intercédant) pro nobis ad Dóminum. \R\ Amen.

\variantheading{Pro ix Lectione}

\benedlabel{Bened.} Ad societátem cívium supernórum perdúcat nos Rex Angelórum. \R\ Amen.

\rubric{Si autem legenda sit Lectio ix de Homilia cum Evangelio Dominicæ, vel Feriæ, aut Vigiliæ:}

\benedlabel{Bened.} Per evangélica dicta deleántur nostra delícta. \R\ Amen.

\variantheading{In Officio trium Lectionum}

\nocturnoheading{In Nocturno}

\rubric{In Officiis trium Lectionum, post Hymnum dicuntur Antiphonæ convenientes, quæ initio Psalmi inchoantur tantum et usque ad Asteriscum * perducuntur, ac deinde in fine integræ pronuntiantur.}

\rubric{Sub eisdem vero Antiphonis dicuntur novem Psalmi Feriæ currentis, quibus subjungitur Versus in III Nocturno positus, omissis Versibus pro I et II Nocturno assignatis.}

\rubric{Post Versum dicitur}

Pater noster \rubric{secreto usque ad} \V\ Et ne nos indúcas in tentatiónem. \R\ Sed líbera nos a malo.

\absolutioheading{Absolutio}

\rubric{Feria II et V}

Exáudi, Dómine, \rubric{ut in I Nocturno.}

\rubric{Feria III et VI}

Ipsíus píetas, \rubric{ut in II Nocturno.}

\rubric{Feria IV et Sabbato}

A vínculis, \rubric{ut in III Nocturno Officii novem Lectionum.}

\rubric{Deinde leguntur Lectiones cum Responsoriis, prouti Officium occurrens requirit; et ante eas dicuntur sequentes}

\variantheading{Benedictiones}

\rubric{In Feriis, quando legitur Homilia cum Evangelio, erit:}

\benedlabel{1 Bened.} Evangélica léctio sit nobis salus et protéctio. \R\ Amen.

\benedlabel{2 Bened.} Divínum auxílium máneat semper nobíscum. \R\ Amen.

\benedlabel{3 Bened.} Ad societátem cívium supernórum perdúcat nos Rex Angelórum. \R\ Amen.

\rubric{In Feriis, quando non legitur Homilia cum Evangelio, Feria II et V Benedictiones ut in I Nocturno, Feria III et VI ut in II Nocturno Officii novem Lectionum; Feria autem IV et Sabbato erit:}

\benedlabel{1 Bened.} Ille nos benedícat, qui sine fine vivit et regnat. \R\ Amen.

\benedlabel{2 Bened.} Divínum auxílium máneat semper nobíscum. \R\ Amen.

\benedlabel{3 Bened.} Ad societátem cívium supernórum perdúcat nos Rex Angelórum. \R\ Amen.

\rubric{In Festis Sanctorum erit:}

\benedlabel{1 Bened.} Ille nos benedícat, qui sine fine vivit et regnat. \R\ Amen.

\benedlabel{2 Bened.} Cujus (\textit{vel} Quorum \textit{aut} Quarum) festum cólimus, ipse (\textit{vel} ipsa \textit{aut} ipsi \textit{vel} ipsæ) intercédat (\textit{vel} intercédant) pro nobis ad Dóminum. \R\ Amen.

\benedlabel{3 Bened.} Ad societátem cívium supernórum perdúcat nos Rex Angelórum. \R\ Amen.

\rubric{In Officio de sancta María in Sabbato Absolutio et Benedictiones, quæ ibidem ponuntur.}

\rubric{Post ultimam Lectionem, in omnibus Dominicis per Annum minoribus, etiam repositis vel anticipatis, in Vigilia Epiphaniæ, in Festis cujusvis ritus, excepto tamen sanctorum Innocentium Festo, nisi hoc in Dominicam incidat, aut ritu gaudeat duplici I classis, per omnes Octavas, et in Officio sanctæ Mariæ in Sabbato, dicitur Hymnus Ambrosianus. In Adventu autem, et a Dominica Septuagesimæ usque ad Sabbatum post Dominicam Quinquagesimæ inclusive, non dicitur nisi in Festis. Alio Tempore numquam dicitur in feriali Officio. Quando vero Hymnus prædictus omittitur, ejus loco dicitur IX aut III Responsorium.}

\hymnusheading{Hymnus Ambrosianus}

\lettrine[lines=2]{T}{E} Deum laudámus: * te Dóminum confitémur.

Te ætérnum Patrem * omnis terra venerátur.

Tibi omnes Angeli, * tibi Cæli, et univérsæ Potestátes:

Tibi Chérubim et Séraphim * incessábili voce proclámant:

Sanctus, * Sanctus, * Sanctus * Dóminus, Deus Sábaoth.

Pleni sunt cæli et terra * majestátis glóriæ tuæ.

Te gloriósus * Apostolórum chorus,

Te Prophetárum * laudábilis númerus,

Te Mártyrum candidátus * laudat exércitus.

Te per orbem terrárum * sancta confitétur Ecclésia,

Patrem * imménsæ majestátis;

Venerándum tuum verum * et únicum Fílium;

Sanctum quoque * Paráclitum Spíritum.

Tu Rex glóriæ, * Christe.

Tu Patris * sempitérnus es Fílius.

Tu, ad liberándum susceptúrus hóminem, * non horruísti Vírginis úterum.

Tu, devícto mortis acúleo, * aperuísti credéntibus regna cælórum.

Tu ad déxteram Dei sedes, * in glória Patris.

Judex créderis * esse ventúrus.

\rubric{Sequens versus dicitur flexis genibus.}

Te ergo quæsumus, tuis fámulis súbveni, * quos pretióso sánguine redemísti.

Ætérna fac cum Sanctis tuis * in glória numerári.

Salvum fac pópulum tuum, Dómine, * et bénedic hereditáti tuæ.

Et rege eos, * et extólle illos usque in ætérnum.

Per síngulos dies * benedícimus te;

Et laudámus nomen tuum in séculum, * et in séculum séculi.

Dignáre, Dómine, die isto * sine peccáto nos custodíre.

Miserére nostri, Dómine, * miserére nostri.

Fiat misericórdia tua, Dómine, super nos, * quemádmodum sperávimus in te.

In te, Dómine, sperávi: * non confúndar in ætérnum.

\rubric{Dicto Te Deum, aut ultimo Responsorio, statim incipiuntur Laudes a Versu Deus, in adjutórium, ut infra.}

\rubric{In privata vero recitatione Matutinum a Laudibus separari potest; quo in casu, post Hymnum Te Deum, vel post ultimum Responsorium, dicitur:}

\V\ Dóminus vobíscum. \R\ Et cum spíritu tuo.

\rubric{Hic Versus Dóminus vobíscum non dicitur ab eo, qui non est saltem in ordine Diaconatus; sed ejus loco substituitur, etiam in aliis Horis, quoties in antecedentibus Precibus dictus non fuerit:}

\V\ Dómine, exáudi oratiónem meam. \R\ Et clamor meus ad te véniat.

Orémus.

\rubric{Et dicitur Oratio ut ad Laudes.}

\rubric{Post Orationem additur:}

\V\ Dóminus vobíscum. \R\ Et cum spíritu tuo.

\V\ Benedicámus Dómino. \R\ Deo grátias.

\V\ Fidélium ánimæ per misericórdiam Dei requiéscant in pace. \R\ Amen.

\rubric{Deinde dicitur Pater noster totum secreto.}

\rubric{In publica vero recitatione Matutinum in Choro a Laudibus separari non potest.}

\sectionornament
% <!-- p. 10 -->
\label{p:10}

\hourheading{Ad Laudes}

\rubric{Si Laudes extra Chorum separentur a Matutino, ante eas dicitur secreto Pater noster et Ave María; secus absolute incipiuntur, ut sequitur:}

\V\ Deus in adjutórium.

\rubric{Deinde, sub congruentibus Antiphonis, dicuntur Psalmi et Cánticum, prouti Officium occurrens requirit.}

\rubric{Repetita vero post ultimum Psalmum Antiphona, dicuntur Capitulum, Hymnus et Versus, prouti Officium occurrens exigit.}

\variantheading{Per Annum}

\rubric{In Officio dominicali: In Dominicis Septuagesimæ, Sexagesimæ et Quinquagesimæ, Capitulum ut in Proprio de Tempore, Hymnus et Versus ut in Psalterio.}

\rubric{In Dominicis per Annum minoribus, extra Octavam Nativitatis, Capitulum, Hymnus et Versus ut in Psalterio.}

\rubric{In Officio feriali: In omnibus per Annum Feriis, et a Feria IV Cinerum usque ad Sabbatum sequens inclusive, Capitulum, Hymnus et Versus de Feria currenti, ut in Psalterio.}

\variantheading{Tempore Adventus}

\rubric{In Officio dominicali, atque in Vigilia Nativitatis Domini, Capitulum ut in Proprio de Tempore.}

\rubric{In Officio autem feriali, usque ad ultimam Feriam ante Vigiliam Nativitatis Domini inclusive, dicitur quotidie sequens}

\capitulumlabel{Capitulum}{Is. 2, 3}

\lettrine[lines=2]{V}{ENITE}, et ascendámus ad montem Dómini, et ad domum Dei Jacob, et docébit nos vias suas, et ambulábimus in sémitis ejus: quia de Sion exíbit lex, et verbum Dómini de Jerúsalem. \R\ Deo grátias.

\rubric{Sic respondetur in fine omnium Capitulorum.}

\rubric{In Officio tamen dominicali et feriali dicitur sequens Hymnus et Versus:}

\hymnusheading{Hymnus}

\lettrine[lines=2]{E}{N}, clara vox redárguit\\
Obscúra quæque, pérsonans:\\
Procul fugéntur sómnia:\\
Ab alto Jesus prómicat.

Mens jam resúrgat, tórpida\\
Non ámplius jacens humi:\\
Sidus refúlget jam novum,\\
Ut tollat omne nóxium.

En, Agnus ad nos míttitur\\
Laxáre gratis débitum:\\
Omnes simul cum lácrimis\\
Precémur indulgéntiam:

Ut, cum secúndo fúlserit\\
Metúque mundum cínxerit,\\
Non pro reátu púniat,\\
Sed nos pius tunc prótegat.

Virtus, honor, laus, glória\\
Deo Patri cum Fílio,\\
Sancto simul Paráclito,\\
In sæculórum sæcula.\\
Amen.

\sectionornament
% <!-- p. 11 -->
\label{p:11}

\V\ Vox clamántis in desérto: Paráte viam Dómini. \R\ Rectas fácite sémitas ejus.

\rubric{Recitato Versu, dicitur cum Antiphona convenienti sequens}

\canticumheading{Canticum Zachariæ}{Luc. 1, 68-69}

\lettrine[lines=2]{B}{ENEDICTUS} \textdagger\ Dóminus, Deus Israël, * quia visitávit, et fecit redemptiónem plebis suæ:

Et eréxit cornu salútis nobis: * in domo David, púeri sui.

Sicut locútus est per os sanctórum, * qui a século sunt, Prophetárum ejus:

Salútem ex inimícis nostris, * et de manu ómnium, qui odérunt nos:

Ad faciéndam misericórdiam cum pátribus nostris: * et memorári testaménti sui sancti.

Jusjurándum, quod jurávit ad Abraham, patrem nostrum, * datúrum se nobis:

Ut sine timóre, de manu inimicórum nostrórum liberáti, * serviámus illi.

In sanctitáte, et justítia coram ipso, * ómnibus diébus nostris.

Et tu, puer, Prophéta Altíssimi vocáberis: * præíbis enim ante fáciem Dómini paráre vias ejus:

Ad dandam sciéntiam salútis plebi ejus: * in remissiónem peccatórum eórum:

Per víscera misericórdiæ Dei nostri: * in quibus visitávit nos, óriens ex alto:

Illumináre his, qui in ténebris, et in umbra mortis sedent: * ad dirigéndos pedes nostros in viam pacis.

Glória Patri, et Fílio, * et Spirítui Sancto.

Sicut erat in princípio, et nunc, et semper, * et in sécula seculórum. Amen.

\rubric{Hic Versus Glória Patri dicitur in fine omnium Canticorum, nisi aliter adnotetur.}

\rubric{Expleto Cantico ac repetita Antiphona, si Preces non fuerint recitandæ, statim dicitur Versus Dóminus vobíscum cum Oratione et reliquis omnibus usque ad finem Horæ, ut infra.}

\rubric{In Feriis vero Adventus, Quadragesimæ a die Cinerum, atque in omnibus Vigiliis præter quam Nativitatis Domini et Epiphaniæ, si fiat Officium de Feria aut Vigilia currenti, et quamvis in eo peragenda sit Commemoratio Duplicis aut Semiduplicis ad instar Simplicis redacti, post repetitam ad Benedictus Antiphonam dicuntur flexis genibus sequentes Preces quæ aliis temporibus omittuntur:}

KYRIE, eléison. Christe, eléison. Kýrie, eléison.

\rubric{Pater noster, quod dicitur a solo Hebdomadario totum clara voce, usque ad} \V\ Et ne nos indúcas in tentatiónem. \R\ Sed líbera nos a malo.

\V\ Ego dixi: Dómine, miserére mei. \R\ Sana ánimam meam, quia peccávi tibi. \V\ Convértere, Dómine, úsquequo? \R\ Et deprecábilis esto super servos tuos. \V\ Fiat misericórdia tua, Dómine, super nos. \R\ Quemádmodum sperávimus in te. \V\ Sacerdótes tui induántur justítiam. \R\ Et sancti tui exsúltent.

\V\ Orémus pro beatíssimo Papa nostro N. \R\ Dóminus consérvet eum, et vivíficet eum, et beátum fáciat eum in terra, et non tradat eum in ánimam inimicórum ejus.

\V\ Orémus et pro Antístite nostro N. \R\ Stet et pascat in fortitúdine tua, Dómine, in sublimitáte nóminis tui.

\rubric{Romæ præcedens Versus cum suo Responsorio omittitur; alibi vero, ad litteram N., ab omnibus nomen diœcesani Episcopi exprimatur. Vacante Apostolica vel Episcopali Sede, alteruter vel uterque respectivus Versus cum suo Responsorio prætéritur.}

\V\ Dómine, salvum fac regem. \R\ Et exáudi nos in die, qua invocavérimus te. \V\ Salvum fac pópulum tuum, Dómine, et bénedic hereditáti tuæ. \R\ Et rege eos, et extólle illos usque in ætérnum. \V\ Meménto Congregatiónis tuæ. \R\ Quam possedísti ab inítio. \V\ Fiat pax in virtúte tua. \R\ Et abundántia in túrribus tuis. \V\ Orémus pro benefactóribus nostris. \R\ Retribúere dignáre, Dómine, ómnibus, nobis bona faciéntibus propter nomen tuum, vitam ætérnam. Amen. \V\ Orémus pro fidélibus defúnctis. \R\ Réquiem ætérnam dona eis, Dómine, et lux perpétua lúceat eis. \V\ Requiéscant in pace. \R\ Amen. \V\ Pro frátribus nostris abséntibus. \R\ Salvos fac servos tuos, Deus meus, sperántes in te. \V\ Pro afflíctis et captívis. \R\ Líbera eos, Deus Israël, ex ómnibus tribulatiónibus suis. \V\ Mitte eis, Dómine, auxílium de sancto. \R\ Et de Sion tuére eos. \V\ Dómine, Deus virtútum, convérte nos. \R\ Et osténde fáciem tuam, et salvi érimus. \V\ Exsúrge, Christe, ádjuva nos. \R\ Et líbera nos propter nomen tuum. \V\ Dómine, exáudi oratiónem meam. \R\ Et clamor meus ad te véniat.

\rubric{Absolutis Precibus, vel, si Preces locum non habeant, repetita Antiphona post Benedictus, statim dicitur:}

\V\ Dóminus vobíscum. \R\ Et cum spíritu tuo.

Orémus.

\rubric{Et dicitur Oratio conveniens. Postea fiunt Commemorationes, si occurrant.}

\rubric{Deinde in Officio semiduplici aut simplici tam de Tempore quam de Sanctis, præter quam a Dominica I Adventus usque ad Octavam Epiphaniæ, ac præter omnes dies, in quibus habita sit Commemoratio Duplicis vel factum fuerit Officium aut Commemoratio cujusvis Octavæ, fit sequens}

\suffragiumheading{Suffragium de omnibus Sanctis}

\antlabel{Ant.} Beáta Dei Génitrix Virgo María, Sanctíque omnes intercédant pro nobis ad Dóminum.

\V\ Mirificávit Dóminus Sanctos suos. \R\ Et exaudívit eos clamántes ad se.

Orémus. \oratiolabel{Oratio}

\lettrine[lines=2]{A}{} \textsc{cunctis} nos, quæsumus, Dómine, mentis et corpóris defénde perículis: et, intercedénte beáta et gloriósa semper Vírgine Dei Genitríce María, cum beáto Joseph, beátis Apóstolis tuis Petro et Paulo, atque beáto N. et ómnibus Sanctis, salútem nobis tríbue benígnus et pacem; ut, destrúctis adversitátibus et erróribus univérsis, Ecclésia tua secúra tibi sérviat libertáte. Per eúmdem Dóminum nostrum. \R\ Amen.

\rubric{In hac Oratione, ad litteram N., exprimitur nomen Titularis propriæ Ecclesiæ, dummodo Titulus non sit Persona divina vel Mysterium Domini, aut de ipso ratione Vigiliæ vel Festi factum non fuerit Officium vel Commemoratio, aut illius nomen non sit in eadem Suffragii Oratione expressum; ac nomina sanctorum Angelorum et sancti Joannis Baptistæ, si Titulares fuerint, præponuntur nomini sancti Joseph. In hisce autem omnibus casibus omittuntur verba atque beáto N.}

\rubric{Si vero Officium vel Commemoratio fuerit de beata María Virgine, tunc dicitur ut infra; atque insuper, sicubi Titularis juxta superiorem Rubricam ante omnes fuerit nominandus, ad Orationem particula cum omittitur ante nomen sanctorum Petri et Pauli ac præponitur nomini sancti Joseph.}

\antlabel{Ant.} Sancti omnes intercédant pro nobis ad Dóminum.

\V\ Mirificávit Dóminus Sanctos suos. \R\ Et exaudívit eos clamántes ad se.

Orémus. \oratiolabel{Oratio}

\lettrine[lines=2]{A}{} \textsc{cunctis} nos, quæsumus, Dómine, mentis et corpóris defénde perículis: et, intercedénte beáto Joseph, cum beátis Apóstolis tuis Petro et Paulo, atque beáto N. et ómnibus Sanctis, salútem nobis tríbue benígnus et pacem; ut, destrúctis adversitátibus et erróribus univérsis, Ecclésia tua secúra tibi sérviat libertáte. Per Dóminum. \R\ Amen.

\rubric{Post ultimam Orationem additur:}

\V\ Dóminus vobíscum. \R\ Et cum spíritu tuo. \V\ Benedicámus Dómino. \R\ Deo grátias. \V\ Fidélium ánimæ per misericórdiam Dei requiéscant in pace. \R\ Amen.

\rubric{Deinde, in publica Officii recitatione, si discedendum sit a Choro, itemque in recitatione privata, si tunc terminetur Officium; alioquin, in fine ultimæ Horæ dicitur Pater noster, totum secreto; eoque recitato, subjungitur:}

\V\ Dóminus det nobis suam pacem. \R\ Et vitam ætérnam. Amen.

\rubric{Et immediate dicitur, cum suis Versu et Oratione, una ex finalibus beatæ Mariæ Virginis Antiphonis, quæ inferius post Completorium pro diversitate Temporis assignantur.}

\rubric{Postea concluditur:}

\V\ Divínum auxílium máneat semper nobíscum. \R\ Amen.

\sectionornament
% <!-- p. 14 -->
\label{p:14}

\hourheading{Ad Primam}

Pater, Ave et Credo.

\V\ Deus, in adjutórium.

\hymnusheading{Hymnus}

\lettrine[lines=2]{J}{AM} lucis orto sídere,\\
Deum precémur súpplices,\\
Ut in diúrnis áctibus\\
Nos servet a nocéntibus.

Linguam refrénans témperet,\\
Ne litis horror ínsonet:\\
Visum fovéndo cóntegat,\\
Ne vanitátes háuriat.

Sint pura cordis íntima,\\
Absístat et vecórdia:\\
Carnis terat supérbiam\\
Potus cibíque párcitas.

Ut, cum dies abscésserit\\
Noctémque sors redúxerit,\\
Mundi per abstinéntiam\\
Ipsi canámus glóriam.

Deo Patri sit glória,\\
Ejúsque soli Fílio,\\
Cum Spíritu Paráclito,\\
Nunc et per omne séculum.\\
Amen.

\rubric{Expleto Hymno, convéniens dicitur Antiphona usque ad Asteriscum *, prouti occurrens Officium requirit.}

\variantheading{Per Annum}

\rubric{In Officio dominicali: In Dominicis Septuagesimæ, Sexagesimæ et Quinquagesimæ Antiphona ut in Proprio de Tempore.}

\rubric{In Dominicis per Annum minoribus, extra Octavam Nativitatis, Antíphona ut in Psalterio.}

\rubric{In Officio feriali: In omnibus per Annum Feriis, et a Feria IV Cinerum usque ad Sabbatum sequens inclusive Antiphona de Feria currenti ut in Psalterio.}

\variantheading{Tempore Adventus}

\rubric{In Officio dominicali, Antiphona prima de Laudibus Dominicæ currentis, ut in Proprio de Tempore.}

\rubric{In Officio feriali, Antiphona prima de Laudibus antecedentis Dominicæ, et in sex Feriis ante Vigiliam Nativitatis Domini ex Laudibus Feriæ currentis, ut in Proprio de Tempore.}

\rubric{Atque in Officiis tam de Tempore quam de Sanctis usque ad Vigiliam Nativitatis Domini inclusive, dicitur quotidie in Responsorio brevi sequens Versus, nisi alius in Festis juxta Rubricas fuerit assumendus:}

\V\ Qui ventúrus es in mundum.

\rubric{Post enuntiatam Antiphonam dicuntur tres Psalmi Officio diei congruentes.}

\rubric{Quando autem adhibetur posterius Laudum schema, Psalmus, in priori Laudum ferialium schemate omissus, ad Primam Feriæ currentis post alios Psalmos resumitur.}

\rubric{Repetita post ultimum Psalmum integre Antiphona, convéniens dicitur Capitulum; nimirum:}

\rubric{In omnibus Dominicis, etiam repositis aut anticipatis, in Vigiliis privilegiatis, in Officio cujuslibet Festi vel Octavæ ac sanctæ Mariæ in Sabbato:}

\capitulumlabel{Capitulum}{1 Tim. 1, 17}

\lettrine[lines=2]{R}{EGI} séculórum immortáli et invisíbili, soli Deo honor et glória in sécula seculórum. Amen. \R\ Deo grátias.

\rubric{In omnibus autem Feriis et in Vigiliis communibus:}

\capitulumlabel{Capitulum}{Zach. 8, 19}

\lettrine[lines=2]{P}{ACEM} et veritátem dilígite, ait Dóminus omnípotens.

\rubric{Dicto Capitulo, subjungitur Responsorium breve, ut sequitur. Versus tamen Qui sedes ad déxteram Patris semper omittitur, quando specialis Versus adhibendus præscribitur, Officii nempe currentis, aut secus Officii ad Laudes et primo quidem loco inter cetera specialem Versum habentia commemorati, aut demum de occurrenti Octava communi vel de Tempore proprius.}

\R\ br. Christe, Fili Dei vivi, * Miserére nobis. \rubric{Et repetitur:} Christe, Fili Dei vivi, miserére nobis. \V\ Qui sedes ad déxteram Patris. \R\ Miserére nobis. \V\ Glória Patri, et Fílio, et Spirítui Sancto. \R\ Christe, Fili Dei vivi, miserére nobis. \V\ Exsúrge, Christe, ádjuva nos. \R\ Et líbera nos propter nomen tuum.

\rubric{Ita dicuntur semper Responsoria brevia, etiam ad Tertiam, Sextam, Nonam et Completorium.}

\rubric{Expleto Responsorio brevi, si Preces non fuerint recitandæ, statim dicitur Versus Dóminus vobíscum cum subsequenti Oratione et reliquis omnibus usque ad finem Horæ ut infra.}

\rubric{In quolibet Officio semiduplici, in Officio sanctæ Mariæ in Sabbato, in Festis simplicibus, atque in Feriis per Annum communibus, post Responsorium breve dicuntur sequentes Preces dominicales; quæ tamen omittuntur in Vigilia Epiphaniæ, et omnibus insuper diebus, in quibus ad Laudes habita sit Commemoratio Duplicis, vel factum fuerit Officium aut Commemoratio cujusvis Octavæ:}

KYRIE, eléison. Christe, eléison. Kýrie, eléison.

Pater noster \rubric{secreto usque ad} \V\ Et ne nos indúcas in tentatiónem. \R\ Sed líbera nos a malo.

Credo in Deum \rubric{secreto usque ad} \V\ Carnis resurrectiónem. \R\ Vitam ætérnam. Amen.

\V\ Et ego ad te, Dómine, clamávi. \R\ Et mane orátio mea prævéniet te. \V\ Repleátur os meum laude. \R\ Ut cantem glóriam tuam, tota die magnitúdinem tuam. \V\ Dómine, avérte fáciem tuam a peccátis meis. \R\ Et omnes iniquitátes meas dele. \V\ Cor mundum crea in me, Deus. \R\ Et spíritum rectum ínnova in viscéribus meis. \V\ Ne projícias me a fácie tua. \R\ Et spíritum sanctum tuum ne áuferas a me. \V\ Redde mihi lætítiam salutáris tui. \R\ Et spíritu principáli confírma me. \V\ Adjutórium nostrum in nómine Dómini. \R\ Qui fecit cælum et terram.

\rubric{Deinde Hebdomadarius facit Confessionem cum Absolutione, eamque Chorus repetit, ut infra post feriales Preces.}

\rubric{In Feriis vero Adventus, et Quadragesimæ, a die Cinerum, atque in omnibus Vigiliis, si ad Laudes recitatæ fuerint Preces, post Responsorium breve dicuntur flexis genibus sequentes Preces feriales:}

KYRIE, eléison. Christe, eléison. Kýrie, eléison.

Pater noster \rubric{secreto usque ad} \V\ Et ne nos indúcas in tentatiónem. \R\ Sed líbera nos a malo.

Credo in Deum \rubric{secreto usque ad} \V\ Carnis resurrectiónem. \R\ Vitam ætérnam. Amen.

\V\ Et ego ad te, Dómine, clamávi. \R\ Et mane orátio mea prævéniet te. \V\ Repleátur os meum laude. \R\ Ut cantem glóriam tuam, tota die magnitúdinem tuam. \V\ Dómine, avérte fáciem tuam a peccátis meis. \R\ Et omnes iniquitátes meas dele. \V\ Cor mundum crea in me, Deus. \R\ Et spíritum rectum ínnova in viscéribus meis. \V\ Ne projícias me a fácie tua. \R\ Et spíritum sanctum tuum ne áuferas a me. \V\ Redde mihi lætítiam salutáris tui. \R\ Et spíritu principáli confírma me. \V\ Eripe me, Dómine, ab hómine malo. \R\ A viro iníquo éripe me. \V\ Eripe me de inimícis meis, Deus meus. \R\ Et ab insurgéntibus in me líbera me. \V\ Eripe me de operántibus iniquitátem. \R\ Et de viris sánguinum salva me. \V\ Sic psalmum dicam nómini tuo in séculum séculi. \R\ Ut reddam vota mea de die in diem. \V\ Exáudi nos, Deus, salutáris noster. \R\ Spes ómnium fínium terræ, et in mari longe. \V\ Deus, in adjutórium meum inténde. \R\ Dómine, ad adjuvándum me festína. \V\ Sanctus Deus, Sanctus fortis, Sanctus immortális. \R\ Miserére nobis. \V\ Bénedic, ánima mea, Dómino. \R\ Et ómnia, quæ intra me sunt, nómini sancto ejus. \V\ Bénedic, ánima mea, Dómino. \R\ Et noli oblivísci omnes retribútiónes ejus. \V\ Qui propitiátur ómnibus iniquitátibus tuis. \R\ Qui sanat omnes infirmitátes tuas. \V\ Qui rédimit de intéritu vitam tuam. \R\ Qui corónat te in misericórdia et miseratiónibus. \V\ Qui replet in bonis desidérium tuum. \R\ Renovábitur ut áquilæ juvéntus tua. \V\ Adjutórium nostrum in nómine Dómini. \R\ Qui fecit cælum et terram.

\rubric{Deinde Hebdomadarius facit Confessionem:}

\lettrine[lines=2]{C}{ONFITEOR} Deo omnipoténti, beátæ Maríæ semper Vírgini, beáto Michaéli Archángelo, beáto Joánni Baptístæ, sanctis Apóstolis Petro et Paulo, ómnibus Sanctis, et vobis, fratres, quia peccávi nimis cogitatióne, verbo et ópere: mea culpa, mea culpa, mea máxima culpa. Ideo precor beátam Maríam semper Vírginem, beátum Michaélem Archángelum, beátum Joánnem Baptístam, sanctos Apóstolos Petrum et Paulum, omnes Sanctos, et vos, fratres, oráre pro me ad Dóminum, Deum nostrum.

\rubric{Chorus respondet:}

\lettrine[lines=2]{M}{ISEREATUR} tui omnípotens Deus, et, dimíssis peccátis tuis, perdúcat te ad vitam ætérnam. \R\ Amen.

\rubric{Deinde repetit Confessionem, et ubi dicitur: vobis, fratres, et vos, fratres, dicat: tibi, pater, et te, pater.}

\rubric{Facta Confessione a Choro, Hebdomadarius dicit:}

\lettrine[lines=2]{M}{ISEREATUR} vestri omnípotens Deus, et, dimíssis peccátis vestris, perdúcat vos ad vitam ætérnam. \R\ Amen.

\lettrine[lines=2]{I}{NDULGENTIAM}, absolutiónem, et remissiónem peccatórum nostrórum tríbuat nobis omnípotens et miséricors Dóminus. \R\ Amen.

\rubric{Extra Chorum, si unus vel duo tantum recitent Officium, et in Choro Monialium, semel tantum ac simul ab omnibus ita fit Confessio:}

\lettrine[lines=2]{C}{ONFITEOR} Deo omnipoténti, beátæ Maríæ semper Vírgini, beáto Michaéli Archángelo, beáto Joánni Baptístæ, sanctis Apóstolis Petro et Paulo, et ómnibus Sanctis, quia peccávi nimis cogitatióne, verbo et ópere: mea culpa, mea culpa, mea máxima culpa. Ideo precor beátam Maríam semper Vírginem, beátum Michaélem Archángelum, beátum Joánnem Baptístam, sanctos Apóstolos Petrum et Paulum, et omnes Sanctos, oráre pro me ad Dóminum, Deum nostrum.

\rubric{Deinde dicitur:}

\lettrine[lines=2]{M}{ISEREATUR} nostri omnípotens Deus, et, dimíssis peccátis nostris, perdúcat nos ad vitam ætérnam. Amen.

\lettrine[lines=2]{I}{NDULGENTIAM}, absolutiónem, et remissiónem peccatórum nostrórum tríbuat nobis omnípotens et miséricors Dóminus. Amen.

\rubric{Facta Absolutione, subjungitur:}

\V\ Dignáre, Dómine, die isto. \R\ Sine peccáto nos custodíre. \V\ Miserére nostri, Dómine. \R\ Miserére nostri. \V\ Fiat misericórdia tua, Dómine, super nos. \R\ Quemádmodum sperávimus in te. \V\ Dómine, exáudi oratiónem meam. \R\ Et clamor meus ad te véniat.

\rubric{Absolutis Precibus, vel, si Preces locum non habeant, post Responsorium breve statim dicitur:}

\V\ Dóminus vobíscum. \R\ Et cum spíritu tuo.

Orémus. \oratiolabel{Oratio}

\lettrine[lines=2]{D}{OMINE}, Deus omnípotens, qui ad princípium hujus diéi nos perveníre fecísti: tua nos hódie salva virtúte; ut in hac die ad nullum declinémus peccátum, sed semper ad tuam justítiam faciéndam nostra procédant elóquia, dirigántur cogitatiónes et ópera. Per Dóminum. \R\ Amen.

\V\ Dóminus vobíscum. \R\ Et cum spíritu tuo. \V\ Benedicámus Dómino. \R\ Deo grátias.

\rubric{Deinde in Choro legitur Martyrologium, quod laudabiliter fit etiam extra Chorum. Postea Hebdomadarius dicit:}

\V\ Pretiósa in conspéctu Dómini. \R\ Mors Sanctórum ejus.

\rubric{Et absolute, sine Orémus:}

\oratiolabel{Oratio}

\lettrine[lines=2]{S}{ANCTA} María, et omnes Sancti intercédant pro nobis ad Dóminum, ut nos mereámur ab eo adjuvári et salvári, qui vivit et regnat in sécula sæculórum. \R\ Amen.

\V\ Deus, in adjutórium meum inténde. \R\ Dómine, ad adjuvándum me festína.

\rubric{Et dicitur ter; ultimo additur:} Glória Patri, et Fílio, et Spirítui Sancto. Sicut erat in princípio, et nunc, et semper, et in sécula seculórum. Amen.

Kýrie, eléison. Christe, eléison. Kýrie, eléison.

Pater noster \rubric{secreto usque ad} \V\ Et ne nos indúcas in tentatiónem. \R\ Sed líbera nos a malo.

\V\ Réspice in servos tuos, Dómine, et in ópera tua, et dírige fílios eórum. \R\ Et sit splendor Dómini, Dei nostri, super nos, et ópera mánuum nostrárum dírige super nos, et opus mánuum nostrárum dírige. \V\ Glória Patri. \R\ Sicut erat.

Orémus. \oratiolabel{Oratio}

\lettrine[lines=2]{D}{IRIGERE} et sanctificáre, régere et gubernáre dignáre, Dómine Deus, Rex cæli et terræ, hódie corda et córpora nostra, sensus, sermónes et actus nostros in lege tua, et in opéribus mandatórum tuórum: ut hic, et in ætérnum, te auxiliánte, salvi et líberi esse mereámur, Salvátor mundi: Qui vivis et regnas in sécula sæculórum. \R\ Amen.

\rubric{Deinde dicitur Lectio brevis, ad absolutionem Capituli, prouti Officium occurrens requirit.}

\V\ Jube, domne, benedícere. \benedlabel{Bened.} Dies et actus nostros in sua pace dispónat Dóminus omnípotens. \R\ Amen.

\variantheading{In Officio dominicali et feriali}

\rubric{A die 14 Januarii usque ad Sabbatum ante Dominicam I Quadragesimæ inclusive:}

\lectiolabel{Lectio br.}{2 Thess. 3, 5}

\lettrine[lines=2]{D}{OMINUS} autem dírigat corda et córpora nostra in caritáte Dei, et patiéntia Christi. Tu autem, Dómine, miserére nobis. \R\ Deo grátias.

\rubric{Ita semper concluduntur Lectiones breves ad Primam, et ita in earum fine respondetur.}

\rubric{A Dominica I Adventus usque ad ultimam diem ante Vigiliam Nativitatis Domini inclusive:}

\lectiolabel{Lectio brevis}{Is. 33, 2}

\lettrine[lines=2]{D}{OMINE}, miserére nostri: te enim exspectávimus: esto bráchium nostrum in mane, et salus nostra in témpore tribulatiónis. Tu autem.

\rubric{In Officio autem cujuslibet Festi vel Octavæ, atque sanctæ Mariæ in Sabbato; itemque in Vigiliis Nativitatis Domini et Epiphaniæ, in Dominica, etiam translata, infra Octavam Nativitatis, si de eis fiat Officium:}

\rubric{Pro Lectione brevi sumitur Capitulum, quod ad Nonam ponitur ipsius Officii currentis, et in fine dicitur: Tu autem, Dómine, ut supra.}

\rubric{Post Lectionem brevem dicitur:}

\V\ Adjutórium nostrum in nómine Dómini. \R\ Qui fecit cælum et terram.

\V\ Benedícite. \R\ Deus.

\benedlabel{Bened.} Dóminus nos benedícat, et ab omni malo deféndat, et ad vitam perdúcat ætérnam. Et fidélium ánimæ per misericórdiam Dei requiéscant in pace. \R\ Amen.

\rubric{Deinde dicitur tantum Pater noster secreto, nisi sequatur alia Hora.}

\sectionornament
% <!-- p. 20 -->
\label{p:20}

\hourheading{Ad Tertiam}

Pater noster et Ave.

\V\ Deus, in adjutórium.

\hymnusheading{Hymnus}

\lettrine[lines=2]{N}{UNC}, Sancte, nobis, Spíritus,\\
Unum Patri cum Fílio,\\
Dignáre promptus íngeri\\
Nostro refúsus péctori.

Os, lingua, mens, sensus, vigor,\\
Confessiónem pérsonent,\\
Flamméscat igne cáritas,\\
Accéndat ardor próximos.

Præsta, Pater piíssime,\\
Patríque compar Unice,\\
Cum Spíritu Paráclito\\
Regnans per omne séculum.\\
Amen.

\rubric{Expleto Hymno, convéniens dicitur Antiphona usque ad Asteriscum *, et tres Psalmi, prouti Officium occurrens requirit; atque Psalmis absolutis, integra repetitur Antiphona.}

\variantheading{Per Annum}

\rubric{In Officio dominicali: In Dominicis Septuagesimæ, Sexagesimæ et Quinquagesimæ, Antiphona et Capitulum ut in Proprio de Tempore.}

\rubric{In Dominicis per Annum minoribus, extra Octavam Nativitatis, Antiphona ut in Psalterio, et dicitur sequens}

\capitulumlabel{Capitulum}{1 Joann. 4, 16}

\lettrine[lines=2]{D}{EUS} cáritas est: et qui manet in caritáte, in Deo manet, et Deus in eo.

\rubric{In omnibus vero Dominicis superius indicatis, dicitur sequens}

\R\ br. Inclína cor meum, Deus, * In testimónia tua. Inclína. \V\ Avérte óculos meos, ne vídeant vanitátem: in via tua vivífica me. In testimónia tua. Glória Patri. Inclína. \V\ Ego dixi: Dómine, miserére mei. \R\ Sana ánimam meam, quia peccávi tibi.

\rubric{In Officio feriali: In omnibus per Annum Feriis, et a Feria IV Cinerum usque ad Sabbatum sequens inclusive, Antiphona de Feria currenti, ut in Psalterio; et dicitur quotidie sequens}

\capitulumlabel{Capitulum}{Jer. 17, 14}

\lettrine[lines=2]{S}{ANA} me, Dómine, et sanábor: salvum me fac, et salvus ero: quóniam laus mea tu es.

\R\ br. Sana ánimam meam, * Quia peccávi tibi. Sana. \V\ Ego dixi: Dómine, miserére mei. Quia. Glória Patri. Sana. \V\ Adjútor meus esto, ne derelínquas me. \R\ Neque despícias me, Deus, salutáris meus.

\variantheading{Tempore Adventus}

\rubric{In Officio dominicali, Antiphona secunda de Laudibus Dominicæ currentis, et Capitulum ut in Proprio de Tempore.}

\rubric{In Officio feriali, Antiphona secunda de Laudibus antecedentis Dominicæ, et in sex Feriis ante Vigiliam Nativitatis Domini ex Laudibus Feriæ currentis, ut in Proprio de Tempore; et dicitur quotidie sequens}

\capitulumlabel{Capitulum}{Jer. 23, 5}

\lettrine[lines=2]{E}{CCE}, dies véniunt, dicit Dóminus, et suscitábo David germen justum: et regnábit rex, et sápiens erit: et fáciet judícium et justítiam in terra.

\rubric{In Officio tamen dominicali et feriali dicitur sequens}

\R\ br. Veni ad liberándum nos, * Dómine, Deus virtútum. Veni. \V\ Osténde fáciem tuam, et salvi érimus. Dómine. Glória Patri. Veni. \V\ Timébunt gentes nomen tuum, Dómine. \R\ Et omnes reges terræ glóriam tuam.

\rubric{Expleto Responsorio brevi, si Preces non fuerint recitandæ, statim dicitur Versus Dóminus vobíscum cum Oratione et reliquis omnibus usque ad finem Horæ, ut infra.}

\rubric{Quando vero ad Laudes recitatæ sunt Preces, post Responsorium breve dicuntur sequentes flexis genibus:}

KYRIE, eléison. Christe, eléison. Kýrie, eléison.

Pater noster \rubric{secreto usque ad} \V\ Et ne nos indúcas in tentatiónem. \R\ Sed líbera nos a malo. \V\ Dómine, Deus virtútum, convérte nos. \R\ Et osténde fáciem tuam, et salvi érimus. \V\ Exsúrge, Christe, ádjuva nos. \R\ Et líbera nos propter nomen tuum. \V\ Dómine, exáudi oratiónem meam. \R\ Et clamor meus ad te véniat.

\rubric{Absolutis Precibus, vel, si Preces locum non habeant, post Responsorium breve statim dicitur:}

\V\ Dóminus vobíscum. \R\ Et cum spíritu tuo.

Orémus.

\rubric{Et dicitur Oratio convéniens. Post Orationem additur:}

\V\ Dóminus vobíscum. \R\ Et cum spíritu tuo. \V\ Benedicámus Dómino. \R\ Deo grátias. \V\ Fidélium ánimæ per misericórdiam Dei requiéscant in pace. \R\ Amen.

\rubric{Deinde dicitur tantum Pater noster secreto, nisi sequatur alia Hora.}

\sectionornament

\hourheading{Ad Sextam}

Pater noster et Ave.

\V\ Deus, in adjutórium.

\hymnusheading{Hymnus}

\lettrine[lines=2]{R}{ECTOR} potens, verax Deus,\\
Qui témperas rerum vices,\\
Splendóre mane illúminas,\\
Et ígnibus merídiem:

Exstíngue flammas lítium,\\
Aufer calórem nóxium,\\
Confer salútem córporum,\\
Verámque pacem córdium.

Præsta, Pater piíssime,\\
Patríque compar Unice,\\
Cum Spíritu Paráclito\\
Regnans per omne séculum.\\
Amen.

\rubric{Expleto Hymno, convéniens dicitur Antiphona usque ad Asteriscum *, et tres Psalmi, prouti Officium occurrens requirit; atque, Psalmis absolutis, integra repetitur Antiphona.}

\variantheading{Per Annum}

\rubric{In Officio dominicali: In Dominicis Septuagesimæ, Sexagesimæ et Quinquagesimæ, Antiphona et Capitulum ut in Proprio de Tempore.}

\rubric{In Dominicis per Annum minoribus, extra Octavam Nativitatis, Antiphona ut in Psalterio, et dicitur sequens}

\capitulumlabel{Capitulum}{Gal. 6, 2}

\lettrine[lines=2]{A}{LTER} altérius ónera portáte, et sic adimplébitis legem Christi.

\rubric{In omnibus vero Dominicis superius indicatis, dicitur sequens}

\R\ br. In ætérnum, Dómine, * Pérmanet verbum tuum. In ætérnum. \V\ In séculum séculi véritas tua. Pérmanet. Glória Patri. In ætérnum. \V\ Dóminus regit me, et nihil mihi déerit. \R\ In loco páscuæ ibi me collocávit.

\rubric{In Officio feriali: In omnibus per Annum Feriis, et a Feria IV Cinerum usque ad Sabbatum sequens inclusive, Antiphona de Feria currenti, ut in Psalterio; et dicitur quotidie sequens}

\capitulumlabel{Capitulum}{Rom. 13, 8}

\lettrine[lines=2]{N}{EMINI} quidquam debeátis, nisi ut ínvicem diligátis: qui enim díligit próximum, legem implévit.

\R\ br. Benedícam Dóminum * In omni témpore. Benedícam. \V\ Semper laus ejus in ore meo. In. Glória Patri. Benedícam. \V\ Dóminus regit me, et nihil mihi déerit. \R\ In loco páscuæ ibi me collocávit.

\variantheading{Tempore Adventus}

\rubric{In Officio dominicali, Antiphona tertia de Laudibus Dominicæ currentis, et Capitulum ut in Proprio de Tempore.}

\rubric{In Officio feriali, Antiphona tertia de Laudibus antecedentis Dominicæ, et in sex Feriis ante Vigiliam Nativitatis Domini ex Laudibus Feriæ currentis, ut in Proprio de Tempore; et dicitur quotidie sequens}

\capitulumlabel{Capitulum}{Jer. 23, 6}

\lettrine[lines=2]{I}{N} diébus illis salvábitur Juda, et Israël habitábit confidénter: et hoc est nomen, quod vocábunt eum, Dóminus justus noster.

\rubric{In Officio tamen dominicali et feriali dicitur sequens}

\R\ br. Osténde nobis, Dómine, * Misericórdiam tuam. Osténde. \V\ Et salutáre tuum da nobis. Misericórdiam tuam. Glória Patri. Osténde. \V\ Meménto nostri, Dómine, in beneplácito pópuli tui. \R\ Vísita nos in salutári tuo.

\rubric{Expleto Responsorio brevi, si Preces non fuerint recitandæ, statim dicitur Versus Dóminus vobíscum cum Oratione et reliquis omnibus usque ad finem Horæ, ut infra.}

\rubric{Quando vero ad Laudes recitatæ sunt Preces, post Responsorium breve dicuntur sequentes flexis genibus:}

KYRIE, eléison. Christe, eléison. Kýrie, eléison.

Pater noster \rubric{secreto usque ad} \V\ Et ne nos indúcas in tentatiónem. \R\ Sed líbera nos a malo. \V\ Dómine, Deus virtútum, convérte nos. \R\ Et osténde fáciem tuam, et salvi érimus. \V\ Exsúrge, Christe, ádjuva nos. \R\ Et líbera nos propter nomen tuum. \V\ Dómine, exáudi oratiónem meam. \R\ Et clamor meus ad te véniat.

\rubric{Absolutis Precibus, vel, si Preces locum non habeant, post Responsorium breve statim dicitur:}

\V\ Dóminus vobíscum. \R\ Et cum spíritu tuo.

Orémus.

\rubric{Et dicitur Oratio convéniens. Post Orationem additur:}

\V\ Dóminus vobíscum. \R\ Et cum spíritu tuo. \V\ Benedicámus Dómino. \R\ Deo grátias. \V\ Fidélium ánimæ per misericórdiam Dei requiéscant in pace. \R\ Amen.

\rubric{Deinde dicitur tantum Pater noster secreto, nisi sequatur alia Hora.}

\sectionornament

\hourheading{Ad Nonam}

Pater noster et Ave.

\V\ Deus, in adjutórium.

\hymnusheading{Hymnus}

\lettrine[lines=2]{R}{ERUM}, Deus, tenax vigor,\\
Immótus in te pérmanens,\\
Lucis diúrnæ témpora\\
Successíbus detérminans:

Largíre lumen véspere,\\
Quo vita nusquam décidat,\\
Sed prémium mortis sacræ\\
Perénnis instet glória.

Præsta, Pater piíssime,\\
Patríque compar Unice,\\
Cum Spíritu Paráclito\\
Regnans per omne séculum.\\
Amen.

\rubric{Expleto Hymno, convéniens dicitur Antiphona usque ad Asteriscum *, et tres Psalmi, prouti Officium occurrens requirit; atque, Psalmis absolutis, integra repetitur Antiphona.}

\variantheading{Per Annum}

\rubric{In Officio dominicali: In Dominicis Septuagesimæ, Sexagesimæ et Quinquagesimæ, Antiphona et Capitulum ut in Proprio de Tempore.}

\rubric{In Dominicis per Annum minoribus, extra Octavam Nativitatis, Antiphona ut in Psalterio, et dicitur sequens}

\capitulumlabel{Capitulum}{1 Cor. 6, 20}

\lettrine[lines=2]{E}{MPTI} enim estis prétio magno. Glorificáte et portáte Deum in córpore vestro.

\rubric{In omnibus vero Dominicis superius indicatis, dicitur sequens}

\R\ br. Clamávi in toto corde meo: * Exáudi me, Dómine. Clamávi. \V\ Justificatiónes tuas requíram. Exáudi. Glória Patri. Clamávi. \V\ Ab occúltis meis munda me, Dómine. \R\ Et ab aliénis parce servo tuo.

\rubric{In Officio feriali: In omnibus per Annum Feriis, et a Feria IV Cinerum usque ad Sabbatum sequens inclusive, Antiphona de Feria currenti, ut in Psalterio; et dicitur quotidie sequens}

\capitulumlabel{Capit.}{1 Petr. 1, 17-19}

\lettrine[lines=2]{I}{N} timóre incolátus vestri témpore conversámini: sciéntes quod non corruptibílibus auro vel argénto redémpti estis, sed pretióso sánguine quasi Agni immaculáti Christi.

\R\ br. Rédime me, Dómine, * Et miserére mei. Rédime. \V\ Pes enim meus stetit in via recta. Et. Glória Patri. Rédime. \V\ Ab occúltis meis munda me, Dómine. \R\ Et ab aliénis parce servo tuo.

\variantheading{Tempore Adventus}

\rubric{In Officio dominicali, Antiphona quinta de Laudibus Dominicæ currentis, et Capitulum ut in Proprio de Tempore.}

\rubric{In Officio feriali: Antiphona quinta de Laudibus antecedentis Dominicæ, et in sex Feriis ante Vigiliam Nativitatis Domini ex Laudibus Feriæ currentis, ut in Proprio de Tempore; et dicitur quotidie sequens}

\capitulumlabel{Capitulum}{Is. 14, 1}

\lettrine[lines=2]{P}{ROPE} est ut véniat tempus ejus, et dies ejus non elongabúntur: miserébitur enim Dóminus Jacob, et Israël salvábitur.

\rubric{In Officio tamen dominicali et feriali dicitur sequens}

\R\ br. Super te, Jerúsalem, * Oriétur Dóminus. Super te. \V\ Et glória ejus in te vidébitur. Oriétur Dóminus. Glória Patri. Super te. \V\ Veni, Dómine, et noli tardáre. \R\ Reláxa facínora plebi tuæ.

\rubric{Expleto Responsorio brevi, si Preces non fuerint recitandæ, statim dicitur Versus Dóminus vobíscum cum Oratione et reliquis omnibus usque ad finem Horæ, ut infra.}

\rubric{Quando vero ad Laudes recitatæ sunt Preces post Responsorium breve dicuntur sequentes flexis genibus:}

KYRIE, eléison. Christe, eléison. Kýrie, eléison.

Pater noster \rubric{secreto usque ad} \V\ Et ne nos indúcas in tentatiónem. \R\ Sed líbera nos a malo. \V\ Dómine, Deus virtútum, convérte nos. \R\ Et osténde fáciem tuam, et salvi érimus. \V\ Exsúrge, Christe, ádjuva nos. \R\ Et líbera nos propter nomen tuum. \V\ Dómine, exáudi oratiónem meam. \R\ Et clamor meus ad te véniat.

\rubric{Absolutis Precibus, vel, si Preces locum non habeant, post Responsorium breve statim dicitur:}

\V\ Dóminus vobíscum. \R\ Et cum spíritu tuo.

Orémus.

\rubric{Et dicitur Oratio convéniens.}

\rubric{Post Orationem additur:}

\V\ Dóminus vobíscum. \R\ Et cum spíritu tuo. \V\ Benedicámus Dómino. \R\ Deo grátias. \V\ Fidélium ánimæ per misericórdiam Dei requiéscant in pace. \R\ Amen.

\rubric{Deinde dicitur tantum Pater noster secreto, nisi sequatur alia Hora.}

\sectionornament

\hourheading{Ad Vesperas}

Pater noster et Ave.

\V\ Deus, in adjutórium.

\rubric{Deinde, sub congruentibus Antiphonis, dicuntur quinque Psalmi, prouti Officium occurrens requirit.}

\rubric{Repetita vero post ultimum Psalmum Antiphona, dicuntur Capitulum, Hymnus et Versus, prouti Officium occurrens exigit.}

\variantheading{Per Annum}

\rubric{In Officio dominicali: In Dominicis Septuagesimæ, Sexagesimæ et Quinquagesimæ, Capitulum ut in Proprio de Tempore; Hymnus et Versus ad I Vesperas e Sabbato, ad II Vesperas e Dominica, ut in Psalterio.}

\rubric{In Dominicis per Annum minoribus, extra Octavam Nativitatis, Capitulum, Hymnus et Versus ad I Vesperas e Sabbato, ad II Vesperas e Dominica, ut in Psalterio.}

\rubric{In Officio feriali: In omnibus per Annum Feriis, et a Feria IV Cinerum usque ad Feriam VI sequentem inclusive, Capitulum, Hymnus et Versus de Feria currenti, ut in Psalterio.}

\variantheading{Tempore Adventus}

\rubric{In Officio dominicali, etiam in I Vesperis Dominicæ IV occurrentis in Vigilia Nativitatis Domini, Capitulum ut in Proprio de Tempore.}

\rubric{In Officio feriali, usque ad ultimam Feriam ante Vigiliam Nativitatis Domini inclusive, dicitur quotidie sequens}

\capitulumlabel{Capitulum}{Gen. 49, 10}

\lettrine[lines=2]{N}{ON} auferétur sceptrum de Juda, et dux de fémore ejus, donec véniat qui mitténdus est: et ipse erit exspectátio géntium.

\rubric{In Officio tamen dominicali et feriali dicitur sequens Hymnus et Versus:}

\hymnusheading{Hymnus}

\lettrine[lines=2]{C}{REATOR} alme síderum,\\
Ætérna lux credéntium,\\
Jesu, Redémptor ómnium,\\
Inténde votis súpplicum.

Qui dæmonis ne fráudibus\\
Períret orbis, ímpetu\\
Amóris actus, lánguidi\\
Mundi medéla factus es.

Commúne qui mundi nefas\\
Ut expiáres, ad Crucem\\
E Vírginis sacrário\\
Intácta prodis víctima.

Cujus potéstas glóriæ,\\
Nomínque cum primum sonat,\\
Et Célites et ínferi\\
Treménte curvántur genu.

Te deprecámur últimæ\\
Magnum diéi Júdicem,\\
Armis supérnæ grátiæ\\
Defénde nos ab hóstibus.

Virtus, honor, laus, glória\\
Deo Patri cum Fílio,\\
Sancto simul Paráclito,\\
In séculórum sécula.\\
Amen.

\V\ Roráte, cæli, désuper, et nubes pluant justum. \R\ Aperiátur terra, et gérminet Salvatórem.

\rubric{Recitato Versu, dicitur cum Antiphona convenienti sequens}

\canticumheading{Canticum beatæ Mariæ Virginis}{Luc. 1, 46-55}

\lettrine[lines=2]{M}{AGNIFICAT} \textdagger\ ánima mea Dóminum:

Et exsultávit spíritus meus * in Deo, salutári meo.

Quia respéxit humilitátem ancíllæ suæ: * ecce enim, ex hoc beátam me dicent omnes generatiónes.

Quia fecit mihi magna qui potens est: * et sanctum nomen ejus.

Et misericórdia ejus a progénie in progénies * timéntibus eum.

Fecit poténtiam in bráchio suo: * dispérsit supérbos mente cordis sui.

Depósuit poténtes de sede, * et exaltávit húmiles.

Esuriéntes implévit bonis: * et dívites dimísit inánes.

Suscépit Israël, púerum suum, * recordátus misericórdiæ suæ.

Sicut locútus est ad patres nostros, * Abraham, et sémini ejus in sécula.

Glória Patri.

\rubric{Expleto Cantico ac repetita Antiphona, si Preces non fuerint recitandæ, statim dicitur Versus Dóminus vobíscum cum Oratione et reliquis omnibus usque ad finem Horæ, ut infra.}

\rubric{In Feriis vero Adventus, in Feria III post Dominicam Quinquagesimæ, si factum fuerit Officium de aliqua Vigilia, et in Feriis Quadragesimæ a die Cinerum, si fiat Officium de Feria currenti, et quamvis in eo peragenda sit Commemoratio Duplicis aut Semiduplicis ad instar Simplicis redacti, post repetitam ad Magníficat Antíphonam dicuntur flexis genibus sequentes Preces, quæ alias omittuntur:}

KYRIE, eléison. Christe, eléison. Kýrie, eléison.

\rubric{Pater noster, quod dicitur a solo Hebdomadario totum clara voce, usque ad} \V\ Et ne nos indúcas in tentatiónem. \R\ Sed líbera nos a malo. \V\ Ego dixi: Dómine, miserére mei. \R\ Sana ánimam meam, quia peccávi tibi. \V\ Convértere, Dómine, úsquequo? \R\ Et deprecábilis esto super servos tuos. \V\ Fiat misericórdia tua, Dómine, super nos. \R\ Quemádmodum sperávimus in te. \V\ Sacerdótes tui induántur justítiam. \R\ Et sancti tui exsúltent.

\V\ Orémus pro beatíssimo Papa nostro N. \R\ Dóminus consérvet eum, et vivíficet eum, et beátum fáciat eum in terra, et non tradat eum in ánimam inimicórum ejus. \V\ Orémus et pro Antístite nostro N. \R\ Stet et pascat in fortitúdine tua, Dómine, in sublimitáte nóminis tui.

\rubric{Romæ præcedens Versus cum suo Responsorio omittitur; alibi vero, ad litteram N., ab omnibus nomen diœcesani Episcopi exprimatur. Vacante Apostolica vel Episcopali Sede, alteruter vel uterque respectivus Versus cum suo Responsorio prætéritur.}

\V\ Dómine, salvum fac regem. \R\ Et exáudi nos in die, qua invocavérimus te. \V\ Salvum fac pópulum tuum, Dómine, et bénedic hereditáti tuæ. \R\ Et rege eos, et extólle illos usque in ætérnum. \V\ Meménto Congregatiónis tuæ. \R\ Quam possedísti ab inítio. \V\ Fiat pax in virtúte tua. \R\ Et abundántia in túrribus tuis. \V\ Orémus pro benefactóribus nostris. \R\ Retribúere dignáre, Dómine, ómnibus, nobis bona faciéntibus propter nomen tuum, vitam ætérnam. Amen. \V\ Orémus pro fidélibus defúnctis. \R\ Réquiem ætérnam dona eis, Dómine, et lux perpétua lúceat eis. \V\ Requiéscant in pace. \R\ Amen. \V\ Pro frátribus nostris abséntibus. \R\ Salvos fac servos tuos, Deus meus, sperántes in te. \V\ Pro afflíctis et captívis. \R\ Líbera eos, Deus Israël, ex ómnibus tribulatiónibus suis. \V\ Mitte eis, Dómine, auxílium de sancto. \R\ Et de Sion tuére eos. \V\ Dómine, Deus virtútum, convérte nos. \R\ Et osténde fáciem tuam, et salvi érimus. \V\ Exsúrge, Christe, ádjuva nos. \R\ Et líbera nos propter nomen tuum. \V\ Dómine, exáudi oratiónem meam. \R\ Et clamor meus ad te véniat.

\rubric{Absolutis Precibus, vel, si Preces locum non habeant, repetita Antiphona post Magníficat, statim dicitur:}

\V\ Dóminus vobíscum. \R\ Et cum spíritu tuo.

Orémus.

\rubric{Et dicitur Oratio convéniens.}

\rubric{Postea fiunt Commemorationes, si occurrant.}

\rubric{Deinde, in Officio semiduplici aut simplici tam de Tempore quam de Sanctis, præter quam a Sabbato ante Dominicam I Adventus usque ad diem Octavam Epiphaniæ, ac præter omnes dies, in quibus habita sit Commemoratio Duplicis, etiam concurrentis, vel factum fuerit Officium aut Commemoratio cujusvis Octavæ, fit sequens}

\suffragiumheading{Suffragium de omnibus Sanctis}

\antlabel{Ant.} Beáta Dei Génitrix Virgo María, Sanctíque omnes intercédant pro nobis ad Dóminum.

\V\ Mirificávit Dóminus Sanctos suos. \R\ Et exaudívit eos clamántes ad se.

Orémus. \oratiolabel{Oratio}

\lettrine[lines=2]{A}{} \textsc{cunctis} nos, quæsumus, Dómine, mentis et corpóris defénde perículis: et, intercedénte beáta et gloriósa semper Vírgine Dei Genitríce María, cum beáto Joseph, beátis Apóstolis tuis Petro et Paulo, atque beáto N. et ómnibus Sanctis, salútem nobis tríbue benígnus et pacem; ut, destrúctis adversitátibus et erróribus univérsis, Ecclésia tua secúra tibi sérviat libertáte. Per eúmdem Dóminum. \R\ Amen.

\rubric{In hac Oratione, ad litteram N., exprimitur nomen Titularis propriæ Ecclesiæ, dummodo Titulus non sit Persona divina vel Mysterium Domini, aut de ipso ratione Festi factum non fuerit Officium vel Commemoratio, aut illius nomen non sit in eadem Suffragii Oratione expressum; ac nomina sanctorum Angelorum et sancti Joannis Baptistæ, si Titulares fuerint, præponuntur nomini sancti Joseph. In hisce autem omnibus casibus omittuntur verba atque beáto N.}

\rubric{Si vero Officium vel Commemoratio fuerit de beata María Virgine, tunc dicitur ut infra; atque insuper, sicubi Titularis juxta superiorem Rubricam ante omnes fuerit nominandus, ad Orationem particula cum omittitur ante nomen sanctorum Petri et Pauli ac præponitur nomini sancti Joseph.}

\antlabel{Ant.} Sancti omnes intercédant pro nobis ad Dóminum.

\V\ Mirificávit Dóminus Sanctos suos. \R\ Et exaudívit eos clamántes ad se.

Orémus. \oratiolabel{Oratio}

\lettrine[lines=2]{A}{} \textsc{cunctis} nos, quæsumus, Dómine, mentis et corpóris defénde perículis: et, intercedénte beáto Joseph, cum beátis Apóstolis tuis Petro et Paulo, atque beáto N. et ómnibus Sanctis, salútem nobis tríbue benígnus et pacem; ut, destrúctis adversitátibus et erróribus univérsis, Ecclésia tua secúra tibi sérviat libertáte. Per Dóminum. \R\ Amen.

\rubric{Post ultimam Orationem additur:}

\V\ Dóminus vobíscum. \R\ Et cum spíritu tuo. \V\ Benedicámus Dómino. \R\ Deo grátias. \V\ Fidélium ánimæ per misericórdiam Dei requiéscant in pace. \R\ Amen.

\rubric{Si post Vesperas immediate sequatur Completorium, dicto Versu Fidélium ánimæ, statim incipitur Versus Jube, domne, benedícere, ut infra ad Completorium; secus autem, si tunc terminetur Officium, dicitur tantum Pater noster secreto.}

\sectionornament

\hourheading{Ad Completorium}

\rubric{Ante Completorium non dicitur Pater noster et Ave María, sicut in aliis Horis, sed absolute Lector incipit:}

\V\ Jube, domne, benedícere. \benedlabel{Bened.} Noctem quiétam, et finem perféctum concédat nobis Dóminus omnípotens. \R\ Amen.

\lectioheading{Lectio brevis}{1 Petr. 5, 8-9}

\lettrine[lines=2]{F}{RATRES}: Sóbrii estóte, et vigiláte: quia adversárius vester diábolus tamquam leo rúgiens círcuit, quærens quem dévoret: cui resístite fortes in fide. Tu autem, Dómine, miserére nobis. \R\ Deo grátias.

\V\ Adjutórium nostrum in nómine Dómini. \R\ Qui fecit cælum et terram.

Pater noster, \rubric{quod dicitur totum secreto.}

\rubric{Deinde Hebdomadarius facit Confessionem:}

\lettrine[lines=2]{C}{ONFITEOR} Deo omnipoténti, beátæ Maríæ semper Vírgini, beáto Michaéli Archángelo, beáto Joánni Baptístæ, sanctis Apóstolis Petro et Paulo, ómnibus Sanctis, et vobis, fratres, quia peccávi nimis cogitatióne, verbo et ópere: mea culpa, mea culpa, mea máxima culpa. Ideo precor beátam Maríam semper Vírginem, beátum Michaélem Archángelum, beátum Joánnem Baptístam, sanctos Apóstolos Petrum et Paulum, omnes Sanctos, et vos, fratres, oráre pro me ad Dóminum, Deum nostrum.

\rubric{Chorus respondet:}

\lettrine[lines=2]{M}{ISEREATUR} tui omnípotens Deus, et, dimíssis peccátis tuis, perdúcat te ad vitam ætérnam. \R\ Amen.

\rubric{Deinde repetit Confessionem, et ubi dicitur: vobis, fratres, et vos, fratres, dicat: tibi, pater, et te, pater.}

\rubric{Facta Confessione a Choro, Hebdomadarius dicit:}

\lettrine[lines=2]{M}{ISEREATUR} vestri omnípotens Deus, et, dimíssis peccátis vestris, perdúcat vos ad vitam ætérnam. \R\ Amen.

\lettrine[lines=2]{I}{NDULGENTIAM}, absolutiónem, et remissiónem peccatórum nostrórum tríbuat nobis omnípotens et miséricors Dóminus. \R\ Amen.

\rubric{Extra Chorum, si unus vel duo tantum recitent Officium, et in Choro Monialium, semel tantum ac simul ab omnibus ita fit Confessio:}

\lettrine[lines=2]{C}{ONFITEOR} Deo omnipoténti, beátæ Maríæ semper Vírgini, beáto Michaéli Archángelo, beáto Joánni Baptístæ, sanctis Apóstolis Petro et Paulo, et ómnibus Sanctis, quia peccávi nimis cogitatióne, verbo et ópere: mea culpa, mea culpa, mea máxima culpa. Ideo precor beátam Maríam semper Vírginem, beátum Michaélem Archángelum, beátum Joánnem Baptístam, sanctos Apóstolos Petrum et Paulum, et omnes Sanctos, oráre pro me ad Dóminum, Deum nostrum.

\rubric{Deinde dicitur:}

\lettrine[lines=2]{M}{ISEREATUR} nostri omnípotens Deus, et, dimíssis peccátis nostris, perdúcat nos ad vitam ætérnam. Amen.

\lettrine[lines=2]{I}{NDULGENTIAM}, absolutiónem, et remissiónem peccatórum nostrórum tríbuat nobis omnípotens et miséricors Dóminus. Amen.

\rubric{Facta Absolutione, dicitur:}

\V\ Convérte nos, Deus, salutáris noster. \R\ Et avérte iram tuam a nobis. \V\ Deus, in adjutórium.

\rubric{Deinde, inchoata convenienti Antiphona, dicuntur tres Psalmi, prouti Officium occurrens requirit.}

\rubric{Repetita integre post ultimum Psalmum Antiphona, dicitur sequens}

\hymnusheading{Hymnus}

\lettrine[lines=2]{T}{E} lucis ante términum,\\
Rerum Creátor, póscimus,\\
Ut pro tua cleméntia\\
Sis præsul et custódia.

Procul recédant sómnia,\\
Et noctium phantásmata;\\
Hostémque nostrum cómprime,\\
Ne polluántur córpora.

Præsta, Pater piíssime,\\
Patríque compar Unice,\\
Cum Spíritu Paráclito\\
Regnans per omne séculum.\\
Amen.

\rubric{Absoluto Hymno, dicitur sequens Capitulum et Responsorium breve:}

\capitulumlabel{Capitulum}{Jer. 14, 9}

\lettrine[lines=2]{T}{U} autem in nobis es, Dómine, et nomen sanctum tuum invocátum est super nos: ne derelínquas nos, Dómine, Deus noster.

\R\ br. In manus tuas, Dómine, * Comméndo spíritum meum. In manus tuas. \V\ Redemísti nos, Dómine, Deus veritátis. Comméndo spíritum meum. Glória Patri, et Fílio, et Spirítui Sancto. In manus tuas, Dómine, comméndo spíritum meum.

\V\ Custódi nos, Dómine, ut pupíllam óculi. \R\ Sub umbra alárum tuárum prótege nos.

\rubric{Postea ad Canticum dicitur sequens}

\antlabel{Ant.} Salva nos.

\canticumheading{Canticum Simeonis}{Luc. 2, 29-32}

\lettrine[lines=2]{N}{UNC} dimíttis servum tuum, Dómine, * secúndum verbum tuum in pace:

Quia vidérunt óculi mei * salutáre tuum.

Quod parásti * ante fáciem ómnium populórum,

Lumen ad revelatiónem géntium, * et glóriam plebis tuæ Israël.

Glória Patri.

\antlabel{Ant.} Salva nos, Dómine, vigilántes, custódi nos dormiéntes; ut vigilémus cum Christo, et requiescámus in pace.

\rubric{Deinde, si Preces non fuerint recitandæ, statim dicitur Versus Dóminus vobíscum cum subsequenti Oratione et reliquis omnibus usque ad finem Completorii, ut infra.}

\rubric{In quolibet Officio semiduplici, in Officio sanctæ Mariæ in Sabbato, in Festis simplicibus, atque in Feriis per Annum communibus, post repetitam ad Nunc dimíttis Antiphonam dicuntur sequentes Preces, quæ tamen omittuntur in Vigilia Epiphaniæ, et omnibus insuper diebus, in quibus ad Vesperas dictæ sint Preces feriales.}

\rubric{Semper autem recitari debent, et flexis quidem genibus, in iis Feriis, in quibus ad Vesperas dictæ sint Preces feriales.}

KYRIE, eléison. Christe, eléison. Kýrie, eléison.

Pater noster \rubric{secreto usque ad} \V\ Et ne nos indúcas in tentatiónem. \R\ Sed líbera nos a malo.

Credo in Deum \rubric{secreto usque ad} \V\ Carnis resurrectiónem. \R\ Vitam ætérnam. Amen.

\V\ Benedíctus es, Dómine, Deus patrum nostrórum. \R\ Et laudábilis et gloriósus in sécula. \V\ Benedicámus Patrem et Fílium cum Sancto Spíritu. \R\ Laudémus, et superexaltémus eum in sécula. \V\ Benedíctus es, Dómine, in firmaménto cæli. \R\ Et laudábilis, et gloriósus, et superexaltátus in sécula. \V\ Benedícat et custódiat nos omnípotens et miséricors Dóminus. \R\ Amen. \V\ Dignáre, Dómine, nocte ista. \R\ Sine peccáto nos custodíre. \V\ Miserére nostri, Dómine. \R\ Miserére nostri. \V\ Fiat misericórdia tua, Dómine, super nos. \R\ Quemádmodum sperávimus in te. \V\ Dómine, exáudi oratiónem meam. \R\ Et clamor meus ad te véniat.

\rubric{Absolutis Precibus, vel, si Preces locum non habeant, repetita Antiphona post Nunc dimíttis, statim dicitur:}

\V\ Dóminus vobíscum. \R\ Et cum spíritu tuo.

Orémus. \oratiolabel{Oratio}

\lettrine[lines=2]{V}{ISITA}, quæsumus, Dómine, habitatiónem istam, et omnes insídias inimíci ab ea longe repélle: Angeli tui sancti hábitent in ea, qui nos in pace custódiant; et benedíctio tua sit super nos semper. Per Dóminum nostrum. \R\ Amen.

\rubric{Post Orationem additur:}

\V\ Dóminus vobíscum. \R\ Et cum spíritu tuo. \V\ Benedicámus Dómino. \R\ Deo grátias.

\benedictiolabel{Benedictio} Benedícat et custódiat nos omnípotens et miséricors Dóminus, Pater, et Fílius, et Spíritus Sanctus. \R\ Amen.

\rubric{Et non dicitur Versus Fidélium ánimæ, sed immediate subjungitur una ex finalibus beátæ Mariæ Virginis Antiphonis, quæ infra pro diversitate Temporis assignantur.}

\rubric{Postea subjungitur:}

\V\ Divínum auxílium máneat semper nobíscum. \R\ Amen.

\rubric{Deinde dicitur secreto Pater noster, Ave María et Credo.}

\rubric{Si post Completorium immediate sequatur Matutinum, tunc bis dicitur Pater noster, Ave María et Credo, pro fine videlicet Officii diei præcedentis, et pro initio Officii diei sequentis.}

\sectionornament

\sectionheading{Post singulas Horas}

\rubric{In fine Laudum, vel, si post Laudes immediate subsequatur alia Hora, in fine ultimæ Horæ, sicuti et semper in fine Completorii, quamvis immediate sequatur Matutinum diei sequentis, tam in publica quam in privata divini Officii recitatione; atque in publica insuper recitatione post alias quascumque Horas, si tunc terminetur Officium et discedatur a Choro, dicitur sequens Antiphona finalis beátæ Mariæ Virginis.}

\rubric{I Post Laudes ac ceteras Horas, præter Completorium, dicto Versu Fidélium ánimæ et Oratione Dominica secreto recitata, subjungitur:}

\V\ Dóminus det nobis suam pacem. \R\ Et vitam ætérnam. Amen.

\rubric{Atque statim dicitur Antiphona finalis, ut infra.}

\rubric{II Post Completorium vero, recitata Benedictione Benedícat et custódiat nos, immediate subjungitur Antiphona finalis, ut infra.}

\sectionheading{Antiphonæ finales Beátæ Mariæ Virginis}

\rubric{1 A Vesperis Sabbati ante Dominicam I Adventus usque ad II Vesperas Purificationis inclusive.}

\antiphonaheading{Antiphona}

\lettrine[lines=2]{A}{LMA} Redemptóris Mater, quæ pérvia cæli\\
Porta manes, et stella maris, succúrre cadénti,\\
Súrgere qui curat, pópulo: tu quæ genuísti,\\
Natúra miránte, tuum sanctum Genitórem,\\
Virgo prius ac postérius, Gabriélis ab ore\\
Sumens illud Ave, peccatórum miserére.

\variantheading{In Adventu}

\V\ Angelus Dómini nuntiávit Maríæ. \R\ Et concépit de Spíritu Sancto.

Orémus. \oratiolabel{Oratio}

\lettrine[lines=2]{G}{RATIAM} tuam, quæsumus, Dómine, méntibus nostris infúnde: ut, qui, Angelo nuntiánte, Christi Fílii tui incarnatiónem cognóvimus; per passiónem ejus et crucem, ad resurrectiónis glóriam perducámur. Per eúmdem Christum Dóminum nostrum. \R\ Amen.

\variantheading{A I Vesperis Nativitatis Domini et deinceps}

\V\ Post partum, Virgo, invioláta permansísti. \R\ Dei Génitrix, intercéde pro nobis.

Orémus. \oratiolabel{Oratio}

\lettrine[lines=2]{D}{EUS}, qui salútis ætérnæ, beátæ Maríæ virginitáte fecúnda, humáno géneri prǽmia præstitísti: tríbue, quæsumus; ut ipsam pro nobis intercédere sentiámus, per quam merúimus auctórem vitæ suscípere, Dóminum nostrum Jesum Christum Fílium tuum. \R\ Amen.

\rubric{2 Post Purificationem, id est a Completorio diei 2 Februarii, etiam quando transfertur Festum Purificationis beátæ Mariæ Virginis, usque ad Nonam Sabbati ante Dominicam I Quadragesimæ inclusive:}

\antiphonaheading{Antiphona}

\lettrine[lines=2]{A}{VE}, Regína cælórum,\\
Ave, Dómina Angelórum:\\
Salve, radix, salve, porta,\\
Ex qua mundo lux est orta:\\
Gaude, Virgo gloriósa,\\
Super omnes speciósa,\\
Vale, o valde decóra,\\
Et pro nobis Christum exóra.

\V\ Dignáre me laudáre te, Virgo sacráta. \R\ Da mihi virtútem contra hostes tuos.

Orémus. \oratiolabel{Oratio}

\lettrine[lines=2]{C}{ONCEDE}, miséricors Deus, fragilitáti nostræ præsídium: ut, qui sanctæ Dei Genitrícis memóriam ágimus; intercessiónis ejus auxílio, a nostris iniquitátibus resurgámus. Per eúmdem Christum Dóminum nostrum. \R\ Amen.

\rubric{Absoluta vero Antiphona finali cum suis Versu et Oratione, concluditur:}

\V\ Divínum auxílium máneat semper nobíscum. \R\ Amen.

\sectionornament

\sectionheading{Post divinum Officium}

\rubric{Cum expletum fuerit Officium, laudabiliter dicitur sequens Oratio; pro qua Summus Pontifex Pius Papa XI sub die 1 Decembris 1933 indulgéntiam trium annórum concessit. Præterea Clericis in sacris constitútis et Sacerdótibus eam devóte post Officium recitántibus Summus Pontifex Leo X defectus et culpas in eo persolvéndo ex humána fragilitáte contráctas, indúlsit. Dicitur autem flexis semper génibus in privata étiam recitatióne, præter quam ab iis, qui ob certam infirmitátem vel graviόris impediménti causam nequeant genufléctere.}

\oratioheading{Oratio}

\lettrine[lines=2]{S}{ACROSANCTAE} et indivíduæ Trinitáti, crucifíxi Dómini nostri Jesu Christi humanitáti, beatíssimæ et gloriosíssimæ semperque Vírginis Maríæ fecúndæ integritáti, et ómnium Sanctórum universitáti sit sempitérna laus, honor, virtus et glória ab omni creatúra, nobísque remíssio ómnium peccatórum, per infiníta sécula sæculórum. \R\ Amen.

\V\ Beáta víscera Maríæ Vírginis, quæ portavérunt ætérni Patris Fílium. \R\ Et beáta úbera, quæ lactavérunt Christum Dóminum.

\rubric{Et dicitur secreto Pater noster et Ave María.}

% <!-- p. [unnumbered; ornamental section-opening page] -->
% [ornamental image beginning next part: Psalterium Breviarii Romani]

% ===================================================================
% BATCH 2: Psalterium Breviarii Romani, Dominica (hiemalis.md lines
% 1395 onward)
% ===================================================================

\partheading{Psalterium Breviarii Romani}
\partsubtitle{Per omnes ac singulos hebdomadæ dies dispositum}

\feriaheading{Dominica}

\hourheading{Ad Matutinum}

Pater, Ave et Credo.

\V\ Dómine, lábia mea.

\V\ Deus, in adjutórium.

\variantheading{In Adventu}

\rubric{In Dominica I et II:}

\invitatlabel{Invit.} Regem ventúrum Dóminum, * Veníte, adorémus.

\rubric{In Dominica III et IV:}

\invitatlabel{Invit.} Prope est jam Dóminus: * Veníte, adorémus.

\rubric{In Vigilia Nativitatis Dom.:}

\invitatlabel{Invit.} Hódie sciétis quia véniet Dóminus: * Et mane vidébitis glóriam ejus.

\hymnusheading{Hymnus}

\lettrine[lines=2]{V}{ERBUM} supérnum, pródiens\\
E Patris ætérni sinu,\\
Qui natus orbi súbvenis,\\
Labénte cursu témporis:

Illúmina nunc péctora,\\
Tuóque amóre cóncrema;\\
Ut cor cadúca déserens\\
Cæli volúptas ímpleat.

Ut, cum tribúnal Júdicis\\
Damnábit igni nóxios,\\
Et vox amíca débitum\\
Vocábit ad cælum pios;

Non esca flammárum nigros\\
Volvámur inter túrbines,\\
Vultu Dei sed cómpotes\\
Cæli fruámur gáudiis.

Deo Patri sit glória,\\
Ejúsque soli Fílio,\\
Cum Spíritu Paráclito,\\
Nunc et per omne séculum.\\
Amen.

\variantheading{Per Annum}

\rubric{In Dominicis post Epiphaniam a die 14 Januarii et deinceps ante Septuagesimam occurrentibus dicitur}

\invitatlabel{Invitat.} Adorémus Dóminum, * Quóniam ipse fecit nos. \psalmlabel{Ps. 94} Veníte. \textdagger

\rubric{In Dominicis Septuagesimæ, Sexagesimæ et Quinquagesimæ dicitur}

\invitatlabel{Invit.} Præoccupémus fáciem Dómini: * Et in psalmis jubilémus ei. \psalmlabel{Ps. 94} Veníte. \textdagger

\rubric{In primo autem Psalmi versu omittitur: præoccupémus fáciem ejus in confessióne, et in psalmis jubilémus ei; et repetitur a Choro Invitatórium.}

\rubric{In omnibus vero Dominicis per Annum superius indicatis dicitur}

\hymnusheading{Hymnus}

\lettrine[lines=2]{P}{RIMO} die, quo Trínitas\\
Beáta mundum cóndidit,\\
Vel quo resúrgens Cónditor\\
Nos, morte victa, líberat:

Pulsis procul torpóribus,\\
Surgámus omnes ócius,\\
Et nocte quærámus Deum,\\
Prophéta sicut præcípit:

Nostras preces ut áudiat,\\
Suámque dextram pórrigat,\\
Et expiátos sórdibus\\
Reddat polórum sédibus:

Ut, quique sacratíssimo\\
Hujus diéi témpore\\
Horis quiétis psállimus,\\
Donis beátis múneret.

Jam nunc, Patérna cláritas,\\
Te postulámus áffatim:\\
Absint faces libídinis,\\
Et omnis actus nóxius.

Ne fœda sit, vel lúbrica\\
Compágo nostri córporis,\\
Ob cujus ignes ígnibus\\
Avérnus urat ácrius.

Mundi Redémptor, quæsumus,\\
Tu probra nostra díluas;\\
Nobísque largus cómmoda\\
Vitæ perénnis cónferas.

Præsta, Pater piíssime,\\
Patríque compar Unice,\\
Cum Spíritu Paráclito\\
Regnans per omne séculum.\\
Amen.

\nocturnoheading{In I Nocturno}

\rubric{In Adventu:} \antlabel{Ant.} Véniet ecce Rex.

\rubric{Per Annum:} \antlabel{Ant.} Beátus vir. \textdagger

\rubric{Et non repetitur in Psalmo; quod semper fit, quando Antiphona incipit a primis Psalmi vel Cantici verbis. Et Psalmus vel Canticum tunc incipitur ab eo verbo, ante quod vel inchoata vel integra desinit Antiphona; si tamen eadem sint verba, et continuetur Antiphona cum Psalmo vel Cantico, neque additur in fine ipsius Antiphonæ Allelúja.}

\rubric{Antiphonæ autem cum Psalmis et Versibus, in Adventu, necnon per Annum assignatæ, adhibentur semper, extra Octavam Nativitatis, quando fit Officium de Dominica seu minori seu majori.}

\rubric{Asteriscus *, in Psalmorum versibus et in aliis Officii partibus insertus, denotat pausam cantus et recitationis in Choro; præter quam in Antiphonis, in quibus nempe significat tum præintonationem Antiphonæ, tum Antiphonam ipsam ante Psalmum seu Canticum eousque producendam esse ad omnes et singulas Horas in Officio ritus semiduplicis ac simplicis, atque ad Horas tantum minores in Officio ritus duplicis.}

\psalmheading{Psalmus 1}

\lettrine[lines=2]{B}{EATUS} vir, \textdagger\ qui non ábiit in consílio impiórum, et in via peccatórum non stetit, * et in cáthedra pestiléntiæ non sedit:

Sed in lege Dómini volúntas ejus, * et in lege ejus meditábitur die ac nocte.

Et erit tamquam lignum, quod plantátum est secus decúrsus aquárum, * quod fructum suum dabit in témpore suo:

Et fólium ejus non défluet: * et ómnia quæcúmque fáciet, prosperabúntur.

Non sic ímpii, non sic: * sed tamquam pulvis, quem prójicit ventus a fácie terræ.

Ideo non resúrgent ímpii in judício: * neque peccatóres in concílio justórum.

Quóniam novit Dóminus viam justórum: * et iter impiórum períbit.

Glória Patri, et Fílio, * et Spirítui Sancto.

Sicut erat in princípio, et nunc, et semper, * et in sécula sæculórum. Amen.

\rubric{Hic versus Glória Patri semper dicitur in fine omnium Psalmorum et Canticorum, nisi aliter notetur.}

\rubric{In Adventu:} \antlabel{Ant.} Véniet ecce Rex excélsus cum potestáte magna ad salvándas gentes, allelúja.

\antlabel{Ant.} Confortáte.

\rubric{Per Annum:} \antlabel{Ant.} Beátus vir, qui in lege Dómini meditátur.

\antlabel{Ant.} Servíte Dómino.

\psalmheading{Psalmus 2}

\lettrine[lines=2]{Q}{UARE} fremuérunt gentes, * et pópuli meditáti sunt inánia?

Astitérunt reges terræ, et príncipes convenérunt in unum * advérsus Dóminum, et advérsus Christum ejus.

Dirumpámus víncula eórum: * et projiciámus a nobis jugum ipsórum.

Qui hábitat in cælis, irridébit eos: * et Dóminus subsannábit eos.

Tunc loquétur ad eos in ira sua, * et in furóre suo conturbábit eos.

Ego autem constitútus sum Rex ab eo super Sion montem sanctum ejus, * prǽdicans præcéptum ejus.

Dóminus dixit ad me: * Fílius meus es tu, ego hódie génui te.

Póstula a me, et dabo tibi gentes hereditátem tuam, * et possessiónem tuam términos terræ.

Reges eos in virga férrea, * et tamquam vas fíguli confrínges eos.

Et nunc, reges, intellígite: * erudímini, qui judicátis terram.

Servíte Dómino in timóre: * et exsultáte ei cum tremóre.

Apprehéndite disciplínam, nequándo irascátur Dóminus, * et pereátis de via justa.

Cum exárserit in brevi ira ejus, * beáti omnes qui confídunt in eo.

\rubric{In Adventu:} \antlabel{Ant.} Confortáte manus dissolútas: confortámini, dícite: Ecce, Deus noster véniet, et salvábit nos, allelúja.

\antlabel{Ant.} Gaudéte omnes.

\rubric{Per Annum:} \antlabel{Ant.} Servíte Dómino in timóre: et exsultáte ei cum tremóre.

\antlabel{Ant.} Exsúrge.

\psalmheading{Psalmus 3}

\lettrine[lines=2]{D}{OMINE}, quid multiplicáti sunt qui tríbulant me? * multi insúrgunt advérsum me.

Multi dicunt ánimæ meæ: * Non est salus ipsi in Deo ejus.

Tu autem, Dómine, suscéptor meus es, * glória mea, et exáltans caput meum.

Voce mea ad Dóminum clamávi: * et exaudívit me de monte sancto suo.

Ego dormívi, et soporátus sum: * et exsurréxi, quia Dóminus suscépit me.

Non timébo míllia pópuli circumdántis me: * exsúrge, Dómine, salvum me fac, Deus meus.

Quóniam tu percussísti omnes adversántes mihi sine causa: * dentes peccatórum contrivísti.

Dómini est salus: * et super pópulum tuum benedíctio tua.

\rubric{In Adventu:} \antlabel{Ant.} Gaudéte omnes, et lætámini: ecce, quia véniet Dóminus ultiónis, addúcet retributiónem: ipse véniet, et salvábit nos.

\V\ Ex Sion spécies decóris ejus. \R\ Deus noster maniféste véniet.

\rubric{Per Annum:} \antlabel{Ant.} Exsúrge, Dómine, salvum me fac, Deus meus.

\V\ Memor fui nocte nóminis tui, Dómine. \R\ Et custodívi legem tuam.

\absolutiolabel{Absol.} Exáudi, Dómine.

\benedlabel{Bened.} 1 Benedictióne perpétua. 2 Unigénitus Dei Fílius. 3 Spíritus Sancti grátia.

\nocturnoheading{In II Nocturno}

\rubric{In Adventu:} \antlabel{Ant.} Gaude et lætáre.

\rubric{Per Annum:} \antlabel{Ant.} Quam admirábile.

\psalmheading{Psalmus 8}

\lettrine[lines=2]{D}{OMINE}, Dóminus noster, * quam admirábile est nomen tuum in univérsa terra!

Quóniam eleváta est magnificéntia tua, * super cælos.

Ex ore infántium et lacténtium perfecísti laudem propter inimícos tuos, * ut déstruas inimícum et ultórem.

Quóniam vidébo cælos tuos, ópera digitórum tuórum: * lunam et stellas, quæ tu fundásti.

Quid est homo, quod memor es ejus? * aut fílius hóminis, quóniam vísitas eum?

Minuísti eum paulo minus ab Angelis, glória et honóre coronásti eum: * et constituísti eum super ópera mánuum tuárum.

Omnia subjecísti sub pédibus ejus, * oves et boves univérsas: ínsuper et pécora campi.

Vólucres cæli, et pisces maris, * qui perámbulant sémitas maris.

Dómine, Dóminus noster, * quam admirábile est nomen tuum in univérsa terra!

\rubric{In Adventu:} \antlabel{Ant.} Gaude et lætáre, fília Jerúsalem: ecce, Rex tuus véniet tibi: Sion, noli timére, quia cito véniet salus tua.

\antlabel{Ant.} Rex noster.

\rubric{Per Annum:} \antlabel{Ant.} Quam admirábile est nomen tuum, Dómine, in univérsa terra!

\antlabel{Ant.} Sedísti super thronum.

\psalmheading{Psalmus 9, i}

\lettrine[lines=2]{C}{ONFITEBOR} tibi, Dómine, in toto corde meo: * narrábo ómnia mirabília tua.

Lætábor et exsultábo in te: * psallam nómini tuo, Altíssime,

In converténdo inimícum meum retrórsum: * infirmabúntur, et períbunt a fácie tua.

Quóniam fecísti judícium meum et causam meam: * sedísti super thronum, qui júdicas justítiam.

Increpásti gentes, et périit ímpius: * nomen eórum delésti in ætérnum, et in séculum séculi.

Inimíci defecérunt frámeæ in finem: * et civitátes eórum destruxísti.

Périit memória eórum cum sónitu: * et Dóminus in ætérnum pérmanet.

Parávit in judício thronum suum: * et ipse judicábit orbem terræ in æquitáte, judicábit pópulos in justítia.

Et factus est Dóminus refúgium páuperi: * adjútor in opportunitátibus, in tribulatióne.

Et sperent in te qui novérunt nomen tuum: * quóniam non dereliquísti quæréntes te, Dómine.

\rubric{In Adventu:} \antlabel{Ant.} Rex noster advéniet Christus, quem Joánnes prædicávit Agnum esse ventúrum.

\antlabel{Ant.} Ecce, vénio.

\rubric{Per Annum:} \antlabel{Ant.} Sedísti super thronum, qui júdicas justítiam.

\antlabel{Ant.} Exsúrge, Dómine.

\psalmheading{Psalmus 9, ii}

\lettrine[lines=2]{P}{SALLITE} Dómino, qui hábitat in Sion: * annuntiáte inter gentes stúdia ejus:

Quóniam requírens sánguinem eórum recordátus est: * non est oblítus clamórem páuperum.

Miserére mei, Dómine: * vide humilitátem meam de inimícis meis.

Qui exáltas me de portis mortis, * ut annúntiem omnes laudatiónes tuas in portis fíliæ Sion.

Exsultábo in salutári tuo: * infíxæ sunt gentes in intéritu, quem fecérunt.

In láqueo isto, quem abscondérunt, * comprehénsus est pes eórum.

Cognoscétur Dóminus judícia fáciens: * in opéribus mánuum suárum comprehénsus est peccátor.

Convertántur peccatóres in inférnum: * omnes gentes quæ obliviscúntur Deum.

Quóniam non in finem oblívio erit páuperis: * patiéntia páuperum non períbit in finem.

Exsúrge, Dómine, non confortétur homo: * judicéntur gentes in conspéctu tuo.

Constítue, Dómine, legislatórem super eos: * ut sciant gentes quóniam hómines sunt.

\rubric{In Adventu:} \antlabel{Ant.} Ecce, vénio cito et merces mea mecum est, dicit Dóminus: dare unicuíque secúndum ópera sua.

\V\ Emítte Agnum, Dómine, Dominatórem terræ. \R\ De Petra desérti ad montem fíliæ Sion.

\rubric{Per Annum:} \antlabel{Ant.} Exsúrge, Dómine, non prævéleat homo.

\V\ Média nocte surgébam ad confiténdum tibi. \R\ Super judícia justificatiónis tuæ.

\absolutiolabel{Absol.} Ipsíus píetas et.

\benedlabel{Bened.} 4 Deus Pater omnípotens. 5 Christus perpétuæ. 6 Ignem sui amóris.

\nocturnoheading{In III Nocturno}

\rubric{In Adventu:} \antlabel{Ant.} Gábriel Angelus.

\rubric{Per Annum:} \antlabel{Ant.} Ut quid, Dómine. \textdagger

\psalmheading{Psalmus 9, iii}

\lettrine[lines=2]{U}{T} quid, Dómine, \textdagger\ recessísti longe, * déspicis in opportunitátibus, in tribulatióne?

Dum supérbit ímpius, incénditur pauper: * comprehendúntur in consíliis quibus cógitant.

Quóniam laudátur peccátor in desidériis ánimæ suæ: * et iníquus benedícitur.

Exacerbávit Dóminum peccátor, * secúndum multitúdinem iræ suæ non quæret.

Non est Deus in conspéctu ejus: * inquinátæ sunt viæ illíus in omni témpore.

Auferúntur judícia tua a fácie ejus: * ómnium inimicórum suórum dominábitur.

Dixit enim in corde suo: * Non movébor a generatióne in generatiónem sine malo.

Cujus maledictióne os plenum est, et amaritúdine, et dolo: * sub lingua ejus labor et dolor.

Sedet in insídiis cum divítibus in occúltis, * ut interfíciat innocéntem.

Óculi ejus in páuperem respíciunt: * insidiátur in abscóndito, quasi leo in spelúnca sua.

Insidiátur ut rápiat páuperem: * rápere páuperem, dum áttrahit eum.

In láqueo suo humiliábit eum: * inclinábit se, et cadet, cum dominátus fúerit páuperum.

Dixit enim in corde suo: Oblítus est Deus, * avértit fáciem suam ne vídeat in finem.

\rubric{In Adventu:} \antlabel{Ant.} Gábriel Angelus locútus est Maríæ, dicens: Ave, grátia plena; Dóminus tecum; benedícta tu inter mulíeres.

\antlabel{Ant.} María dixit.

\rubric{Per Annum:} \antlabel{Ant.} Ut quid, Dómine, recessísti longe?

\antlabel{Ant.} Exsúrge. \textdagger

\psalmheading{Psalmus 9, iv}

\lettrine[lines=2]{E}{XSURGE}, \textdagger\ Dómine Deus, exaltétur manus tua: * ne obliviscáris páuperum.

Propter quid irritávit ímpius Deum? * dixit enim in corde suo: Non requíret.

Vides quóniam tu labórem et dolórem consíderas: * ut tradas eos in manus tuas.

Tibi derelíctus est pauper: * órphano tu eris adjútor.

Cóntere bráchium peccatóris et malígni: * quærétur peccátum illíus, et non inveniétur.

Dóminus regnábit in ætérnum, et in séculum séculi: * períbitis, gentes, de terra illíus.

Desidérium páuperum exaudívit Dóminus: * præparatiónem cordis eórum audívit auris tua.

Judicáre pupíllo et húmili, * ut non appónat ultra magnificáre se homo super terram.

\rubric{In Adventu:} \antlabel{Ant.} María dixit: Putas quális est ista salutátio, quia conturbáta est ánima mea, et quia paritúra sum Regem, qui claustrum virginitátis meæ non violábit?

\antlabel{Ant.} In advéntu.

\rubric{Per Annum:} \antlabel{Ant.} Exsúrge, Dómine Deus, exaltétur manus tua.

\antlabel{Ant.} Justus Dóminus.

\psalmheading{Psalmus 10}

\lettrine[lines=2]{I}{N} Dómino confído: quómodo dícitis ánimæ meæ: * Tránsmigra in montem sicut passer?

Quóniam ecce peccatóres intendérunt arcum, paravérunt sagíttas suas in pháretra, * ut sagíttent in obscúro rectos corde.

Quóniam quæ perfecísti, destruxérunt: * justus autem quid fecit?

Dóminus in templo sancto suo, * Dóminus in cælo sedes ejus:

Óculi ejus in páuperem respíciunt: * pálpebræ ejus intérrogant fílios hóminum.

Dóminus intérrogat justum et ímpium: * qui autem díligit iniquitátem, odit ánimam suam.

Pluet super peccatóres láqueos: * ignis, et sulphur, et spíritus procellárum pars cálicis eórum.

Quóniam justus Dóminus, et justítias diléxit: * æquitátem vidit vultus ejus.

\rubric{In Adventu:} \antlabel{Ant.} In advéntu summi Regis mundéntur corda hóminum, ut digne ambulémus in occúrsum illíus: quia ecce véniet, et non tardábit.

\V\ Egrediétur Dóminus de loco sancto suo. \R\ Véniet, ut salvet pópulum suum.

\rubric{In Vigilia Nativ. Domini:} \V\ Hódie sciétis, quia véniet Dóminus. \R\ Et mane vidébitis glóriam ejus.

\rubric{Per Annum:} \antlabel{Ant.} Justus Dóminus, et justítiam diléxit.

\V\ Prævenérunt óculi mei ad te dilúculo. \R\ Ut meditárer elóquia tua, Dómine.

\absolutiolabel{Absolutio} A vínculis.

\benedlabel{Bened.} 7 Evangélica léctio. 8 Divínum auxílium. 9 Ad societátem cívium.

\sectionornament

\hourheading{Ad Laudes I}

\textit{A Dominica I Adventus usque ad Dominicam ante Septuagesimam inclusive.}

(Pater noster et Ave).

\V\ Deus, in adjutórium.

\rubric{In Adventu: Antiphonæ de Laudibus Dominicæ currentis ut in Proprio de Tempore.}

\rubric{Per Annum:} \antlabel{Ant.} Allelúja.

\rubric{Antiphonæ cum Psalmis, per Annum assignatæ, adhibentur semper, extra Octavam Nativitatis, quando fit Officium de qualibet Dominica minori.}

\psalmheading{Psalmus 92, 1}

\lettrine[lines=2]{D}{OMINUS} regnávit, decórem indútus est: * indútus est Dóminus fortitúdinem, et præcínxit se.

Etenim firmávit orbem terræ, * qui non commovébitur.

Paráta sedes tua ex tunc: * a século tu es.

Elevavérunt flúmina, Dómine: * elevavérunt flúmina vocem suam.

Elevavérunt flúmina fluctus suos, * a vócibus aquárum multárum.

Mirábiles elatiónes maris: * mirábilis in altis Dóminus.

Testimónia tua credibília facta sunt nimis: * domum tuam decet sanctitúdo, Dómine, in longitúdinem diérum.

\antlabel{Ant.} Allelúja, Dóminus regnávit, decórem índuit, allelúja, allelúja.

\antlabel{Ant.} Jubiláte. \textdagger

\psalmheading{Psalmus 99, 2}

\lettrine[lines=2]{J}{UBILATE} \textdagger\ Deo, omnis terra: * servíte Dómino in lætítia.

Introíte in conspéctu ejus, * in exsultatióne.

Scitóte quóniam Dóminus ipse est Deus: * ipse fecit nos, et non ipsi nos.

Pópulus ejus, et oves páscuæ ejus: * introíte portas ejus in confessióne, átria ejus in hymnis: confitémini illi.

Laudáte nomen ejus: quóniam suávis est Dóminus, in ætérnum misericórdia ejus, * et usque in generatiónem et generatiónem véritas ejus.

\antlabel{Ant.} Jubiláte Deo, omnis terra, allelúja.

\antlabel{Ant.} Benedícam te.

\psalmheading{Psalmus 62, 3}

\lettrine[lines=2]{D}{EUS}, Deus meus, * ad te de luce vígilo.

Sitívit in te ánima mea, * quam multiplíciter tibi caro mea.

In terra desérta, et ínvia, et inaquósa: * sic in sancto appárui tibi, ut vidérem virtútem tuam, et glóriam tuam.

Quóniam mélior est misericórdia tua super vitas: * lábia mea laudábunt te.

Sic benedícam te in vita mea: * et in nómine tuo levábo manus meas.

Sicut ádipe et pinguédine repleátur ánima mea: * et lábiis exsultatiónis laudábit os meum.

Si memor fui tui super stratum meum, in matutínis meditábor in te: * quia fuísti adjútor meus.

Et in velaménto alárum tuárum exsultábo, adhǽsit ánima mea post te: * me suscépit déxtera tua.

Ipsi vero in vanum quæsiérunt ánimam meam, introíbunt in inferióra terræ: * tradéntur in manus gládii, partes vúlpium erunt.

Rex vero lætábitur in Deo, laudabúntur omnes qui jurant in eo: * quia obstrúctum est os loquéntium iníqua.

\antlabel{Ant.} Benedícam te in vita mea, Dómine: et in nómine tuo levábo manus meas, allelúja.

\antlabel{Ant.} Tres púeri.

\canticumheading{Canticum trium Puerorum, 4}{Dan. 3, 57-88 et 56}

\lettrine[lines=2]{B}{ENEDICITE}, ómnia ópera Dómini, Dómino: * laudáte et superexaltáte eum in sécula.

Benedícite, Angeli Dómini, Dómino: * benedícite, cæli, Dómino.

Benedícite, aquæ omnes, quæ super cælos sunt, Dómino: * benedícite, omnes virtútes Dómini, Dómino.

Benedícite, sol et luna, Dómino: * benedícite, stellæ cæli, Dómino.

Benedícite, omnis imber et ros, Dómino: * benedícite, omnes spíritus Dei, Dómino.

Benedícite, ignis et æstus, Dómino: * benedícite, frigus et æstus, Dómino.

Benedícite, rores et pruína, Dómino: * benedícite, gelu et frigus, Dómino.

Benedícite, glácies et nives, Dómino: * benedícite, noctes et dies, Dómino.

Benedícite, lux et ténebræ, Dómino: * benedícite, fúlgura et nubes, Dómino.

Benedícat terra Dóminum: * laudet et superexáltet eum in sécula.

Benedícite, montes et colles, Dómino: * benedícite, univérsa germinántia in terra, Dómino.

Benedícite, fontes, Dómino: * benedícite, mária et flúmina, Dómino.

Benedícite, cete, et ómnia, quæ movéntur in aquis, Dómino: * benedícite, omnes vólucres cæli, Dómino.

Benedícite, omnes béstiæ et pécora, Dómino: * benedícite, fílii hóminum, Dómino.

Benedícat Israël Dóminum: * laudet et superexáltet eum in sécula.

Benedícite, sacerdótes Dómini, Dómino: * benedícite, servi Dómini, Dómino.

Benedícite, spíritus, et ánimæ justórum, Dómino: * benedícite, sancti, et húmiles corde, Dómino.

Benedícite, Ananía, Azaría, Mísaël, Dómino: * laudáte et superexaltáte eum in sécula.

Benedicámus Patrem et Fílium cum Sancto Spíritu: * laudémus et superexaltémus eum in sécula.

Benedíctus es, Dómine, in firmaménto cæli: * et laudábilis, et gloriósus, et superexaltátus in sécula.

\rubric{Hic non dicitur Glória Patri, neque Amen.}

\antlabel{Ant.} Tres púeri jussu regis in fornácem missi sunt, non timéntes flammam ignis, dicéntes: Benedíctus Deus, allelúja.

\antlabel{Ant.} Allelúja.

\psalmheading{Psalmus 148, 5}

\lettrine[lines=2]{L}{AUDATE} Dóminum de cælis: * laudáte eum in excélsis.

Laudáte eum, omnes Angeli ejus: * laudáte eum, omnes virtútes ejus.

Laudáte eum, sol et luna: * laudáte eum, omnes stellæ et lumen.

Laudáte eum, cæli cælórum: * et aquæ omnes, quæ super cælos sunt, laudent nomen Dómini.

Quia ipse dixit, et facta sunt: * ipse mandávit, et creáta sunt.

Státuit ea in ætérnum, et in séculum séculi: * præcéptum pósuit, et non præteríbit.

Laudáte Dóminum de terra, * dracónes, et omnes abýssi.

Ignis, grando, nix, glácies, spíritus procellárum: * quæ fáciunt verbum ejus:

Montes, et omnes colles: * ligna fructífera, et omnes cedri.

Béstiæ, et univérsa pécora: * serpéntes, et vólucres pennátæ:

Reges terræ, et omnes pópuli: * príncipes, et omnes júdices terræ.

Júvenes, et vírgines: senes cum junióribus laudent nomen Dómini: * quia exaltátum est nomen ejus solíus.

Conféssio ejus super cælum et terram: * et exaltávit cornu pópuli sui.

Hymnus ómnibus sanctis ejus: * fíliis Israël, pópulo appropinquánti sibi.

\antlabel{Ant.} Allelúja, laudáte Dóminum de cælis, allelúja, allelúja.

\variantheading{In Adventu}

\rubric{Capitulum, Hymnus et Versus ut in Proprio de Tempore.}

\variantheading{Per Annum}

\rubric{In omnibus Dominicis minoribus extra Octavam Nativitatis, dicitur}

\capitulumlabel{Capitulum}{Apoc. 7, 12}

\lettrine[lines=2]{B}{ENEDICTIO}, et cláritas, et sapiéntia, et gratiárum áctio, honor, virtus, et fortitúdo Deo nostro in sécula sæculórum. Amen. \R\ Deo grátias.

\rubric{In Dominicis post Epiphaniam a die 14 Januarii et deinceps ante Septuagesimam occurrentibus dicitur}

\hymnusheading{Hymnus}

\lettrine[lines=2]{A}{ETERNE} rerum Cónditor,\\
Noctem diémque qui regis,\\
Et témporum das témpora,\\
Ut álleves fastídium.

Noctúrna lux viántibus\\
A nocte noctem ségregans,\\
Præco diéi jam sonat,\\
Jubárque solis évocat.

Hoc excitátus Lúcifer\\
Solvit polum calígine:\\
Hoc omnis errónum cohors\\
Viam nocéndi déserit.

Hoc nauta vires cólligit,\\
Pontíque mitéscunt freta:\\
Hoc, ipsa petra Ecclésiæ,\\
Canénte, culpam díluit.

Surgámus ergo strénue:\\
Gallus jacéntes éxcitat,\\
Et somnoléntos íncrepat,\\
Gallus negántes árguit.

Gallo canénte, spes redit,\\
Ægris salus refúnditur,\\
Mucro latrónis cónditur,\\
Lapsis fides revértitur.

Jesu, labántes réspice,\\
Et nos vidéndo córrige:\\
Si réspicis, labes cadunt,\\
Fletúque culpa sólvitur.

Tu, lux, refúlge sénsibus,\\
Mentísque somnum díscute:\\
Te nostra vox primum sonet,\\
Et vota solvámus tibi.

* Deo Patri sit glória,\\
Ejúsque soli Fílio,\\
Cum Spíritu Paráclito,\\
Nunc et per omne séculum.\\
Amen.

\V\ Dóminus regnávit, decórem índuit. \R\ Induit Dóminus fortitúdinem, et præcínxit se virtúte.

\rubric{Ad Benedíctus Antiphona ut in Proprio de Tempore.}

\canticumheading{Canticum Zachariæ}{Luc. 1, 68-79}

\lettrine[lines=2]{B}{ENEDICTUS} \textdagger\ Dóminus, Deus Israël, * quia visitávit, et fecit redemptiónem plebis suæ:

Et eréxit cornu salútis nobis: * in domo David, púeri sui.

Sicut locútus est per os sanctórum, * qui a século sunt, prophetárum ejus:

Salútem ex inimícis nostris, * et de manu ómnium, qui odérunt nos:

Ad faciéndam misericórdiam cum pátribus nostris: * et memorári testaménti sui sancti.

Jusjurándum, quod jurávit ad Abraham, patrem nostrum, * datúrum se nobis:

Ut sine timóre, de manu inimicórum nostrórum liberáti, * serviámus illi.

In sanctitáte, et justítia coram ipso, * ómnibus diébus nostris.

Et tu, puer, Prophéta Altíssimi vocáberis: * præíbis enim ante fáciem Dómini paráre vias ejus:

Ad dandam sciéntiam salútis plebi ejus: * in remissiónem peccatórum eórum:

Per víscera misericórdiæ Dei nostri: * in quibus visitávit nos, óriens ex alto:

Illumináre his, qui in ténebris, et in umbra mortis sedent: * ad dirigéndos pedes nostros in viam pacis.

\rubric{Oratio convéniens.}

\rubric{Reliqua ut in Ordinario, 11.}

\rubric{Psalmi cum Cantico de Dominica superius appositi dicuntur etiam, sub Antiphonis Officio diei congruentibus, in Vigilia Nativitatis Domini et Epiphaniæ, in Dominica, etiam translata, infra Octavam Nativitatis, in Duplicibus I vel II classis, in Festis novem Lectionum Domini, B. Mariæ Virginis, Angelórum, S. Joannis Baptistæ, S. Joseph, Apostolórum, Evangelistárum et per Octavas privilegiátas; itemque in quolibet Festo novem Lectionum próprias Antiphonas ad Laudes habénti, quando de eis persolvitur Officium.}

\sectionornament

\hourheading{Ad Laudes II}

\textit{In Dominicis Septuagesimæ, Sexagesimæ et Quinquagesimæ.}

(Pater noster et Ave.)

\V\ Deus, in adjutórium.

\rubric{Antiphonæ de Laudibus Dominicæ currentis, ut in Proprio de Tempore.}

\psalmheading{Psalmus 50}

\lettrine[lines=2]{M}{ISERERE} \textdagger\ mei, Deus, * secúndum magnam misericórdiam tuam.

Et secúndum multitúdinem miseratiónum tuárum, * dele iniquitátem meam.

Amplius lava me ab iniquitáte mea: * et a peccáto meo munda me.

Quóniam iniquitátem meam ego cognósco: * et peccátum meum contra me est semper.

Tibi soli peccávi, et malum coram te feci: * ut justificéris in sermónibus tuis, et vincas cum judicáris.

Ecce enim, in iniquitátibus concéptus sum: * et in peccátis concépit me mater mea.

Ecce enim, veritátem dilexísti: * incérta et occúlta sapiéntiæ tuæ manifestásti mihi.

Aspérges me hyssópo, et mundábor: * lavábis me, et super nivem dealbábor.

Audítui meo dabis gáudium et lætítiam: * et exsultábunt ossa humiliáta.

Avérte fáciem tuam a peccátis meis: * et omnes iniquitátes meas dele.

Cor mundum crea in me, Deus: * et spíritum rectum ínnova in viscéribus meis.

Ne projícias me a fácie tua: * et spíritum sanctum tuum ne áuferas a me.

Redde mihi lætítiam salutáris tui: * et spíritu principáli confírma me.

Docébo iníquos vias tuas: * et ímpii ad te converténtur.

Líbera me de sanguínibus, Deus, Deus salútis meæ: * et exsultábit lingua mea justítiam tuam.

Dómine, lábia mea apéries: * et os meum annuntiábit laudem tuam.

Quóniam si voluísses sacrifícium, dedíssem útique: * holocáustis non delectáberis.

Sacrifícium Deo spíritus contribulátus: * cor contrítum, et humiliátum, Deus, non despícies.

Benígne fac, Dómine, in bona voluntáte tua Sion: * ut ædificéntur muri Jerúsalem.

Tunc acceptábis sacrifícium justítiæ, oblatiónes, et holocáusta: * tunc impónent super altáre tuum vítulos.

\psalmheading{Psalmus 117}

\lettrine[lines=2]{C}{ONFITEMINI} Dómino quóniam bonus: * quóniam in séculum misericórdia ejus.

Dicat nunc Israël quóniam bonus: * quóniam in séculum misericórdia ejus.

Dicat nunc domus Aaron: * quóniam in séculum misericórdia ejus.

Dicant nunc qui timent Dóminum: * quóniam in séculum misericórdia ejus.

De tribulatióne invocávi Dóminum: * et exaudívit me in latitúdine Dóminus.

Dóminus mihi adjútor: * non timébo quid fáciat mihi homo.

Dóminus mihi adjútor: * et ego despíciam inimícos meos.

Bonum est confídere in Dómino, * quam confídere in hómine:

Bonum est speráre in Dómino, * quam speráre in princípibus.

Omnes gentes circuiérunt me: * et in nómine Dómini quia ultus sum in eos.

Circumdántes circumdedérunt me: * et in nómine Dómini quia ultus sum in eos.

Circumdedérunt me sicut apes, et exarsérunt sicut ignis in spinis: * et in nómine Dómini quia ultus sum in eos.

Impúlsus evérsus sum ut cáderem: * et Dóminus suscépit me.

Fortitúdo mea, et laus mea Dóminus: * et factus est mihi in salútem.

Vox exsultatiónis, et salútis * in tabernáculis justórum.

Déxtera Dómini fecit virtútem: déxtera Dómini exaltávit me, * déxtera Dómini fecit virtútem.

Non móriar, sed vivam: * et narrábo ópera Dómini.

Castígans castigávit me Dóminus: * et morti non trádidit me.

Aperíte mihi portas justítiæ, ingréssus in eas confitébor Dómino: * hæc porta Dómini, justi intrábunt in eam.

Confitébor tibi quóniam exaudísti me: * et factus es mihi in salútem.

Lápidem, quem reprobavérunt ædificántes: * hic factus est in caput ánguli.

A Dómino factum est istud: * et est mirábile in óculis nostris.

Hæc est dies, quam fecit Dóminus: * exsultémus, et lætémur in ea.

O Dómine, salvum me fac, o Dómine, bene prosperáre: * benedíctus qui venit in nómine Dómini.

Benedíximus vobis de domo Dómini: * Deus Dóminus, et illúxit nobis.

Constitúite diem solémnem in condénsis, * usque ad cornu altáris.

Deus meus es tu, et confitébor tibi: * Deus meus es tu, et exaltábo te.

Confitébor tibi quóniam exaudísti me: * et factus es mihi in salútem.

Confitémini Dómino quóniam bonus: * quóniam in séculum misericórdia ejus.

\psalmheading{Psalmus 62}

\lettrine[lines=2]{D}{EUS}, Deus meus, \textdagger\ ad te de luce vígilo.

Sitívit in te ánima mea, * quam multiplíciter tibi caro mea.

In terra desérta, et ínvia, et inaquósa: * sic in sancto appárui tibi, ut vidérem virtútem tuam, et glóriam tuam.

Quóniam mélior est misericórdia tua super vitas: * lábia mea laudábunt te.

Sic benedícam te in vita mea: * et in nómine tuo levábo manus meas.

Sicut ádipe et pinguédine repleátur ánima mea: * et lábiis exsultatiónis laudábit os meum.

Si memor fui tui super stratum meum, in matutínis meditábor in te: * quia fuísti adjútor meus.

Et in velaménto alárum tuárum exsultábo, adhǽsit ánima mea post te: * me suscépit déxtera tua.

Ipsi vero in vanum quæsiérunt ánimam meam, introíbunt in inferióra terræ: * tradéntur in manus gládii, partes vúlpium erunt.

Rex vero lætábitur in Deo, laudabúntur omnes qui jurant in eo: * quia obstrúctum est os loquéntium iníqua.

\canticumheading{Canticum trium Puerorum}{Dan. 3, 52-57}

\lettrine[lines=2]{B}{ENEDICTUS} es, \textdagger\ Dómine, Deus patrum nostrórum: * et laudábilis, et gloriósus, et superexaltátus in sécula.

Et benedíctum nomen glóriæ tuæ sanctum: * et laudábile, et superexaltátum in ómnibus sæculis.

Benedíctus es in templo sancto glóriæ tuæ: * et superlaudábilis, et supergloriósus in sécula.

Benedíctus es in throno regni tui: * et superlaudábilis, et superexaltátus in sécula.

Benedíctus es, qui intuéris abýssos, et sedes super Chérubim: * et laudábilis, et superexaltátus in sécula.

Benedíctus es in firmaménto cæli: * et laudábilis, et gloriósus in sécula.

Benedícite, ómnia ópera Dómini, Dómino: * laudáte, et superexaltáte eum in sécula.

Glória Patri.

\psalmheading{Psalmus 148}

\lettrine[lines=2]{L}{AUDATE} \textdagger\ Dóminum de cælis: * laudáte eum in excélsis.

Laudáte eum, omnes Angeli ejus: * laudáte eum, omnes virtútes ejus.

Laudáte eum, sol et luna: * laudáte eum, omnes stellæ, et lumen.

Laudáte eum, cæli cælórum: * et aquæ omnes, quæ super cælos sunt, laudent nomen Dómini.

Quia ipse dixit, et facta sunt: * ipse mandávit, et creáta sunt.

Státuit ea in ætérnum, et in séculum séculi: * præcéptum pósuit, et non præteríbit.

Laudáte Dóminum de terra, * dracónes, et omnes abýssi.

Ignis, grando, nix, glácies, spíritus procellárum: * quæ fáciunt verbum ejus:

Montes, et omnes colles: * ligna fructífera, et omnes cedri.

Béstiæ, et univérsa pécora: * serpéntes, et vólucres pennátæ.

Reges terræ, et omnes pópuli: * príncipes, et omnes júdices terræ.

Júvenes, et vírgines: senes cum junióribus laudent nomen Dómini: * quia exaltátum est nomen ejus solíus.

Conféssio ejus super cælum, et terram: * et exaltávit cornu pópuli sui.

Hymnus ómnibus sanctis ejus: * fíliis Israël, pópulo appropinquánti sibi.

\rubric{Capitulum, Hymnus, Versus, Antiphona ad Benedíctus et Oratio ut in Proprio de Tempore.}

\rubric{Reliqua ut in Ordinario, 12.}

\sectionornament

\hourheading{Ad Primam}

Pater, Ave et Credo.

\V\ Deus, in adjutórium.

\hymnusheading{Hymnus}

\lettrine[lines=2]{J}{AM} lucis orto sídere,\\
Deum precémur súpplices,\\
Ut in diúrnis áctibus\\
Nos servet a nocéntibus.

Linguam refrénans témperet,\\
Ne litis horror ínsonet:\\
Visum fovéndo cóntegat,\\
Ne vanitátes háuriat.

Sint pura cordis íntima,\\
Absístat et vecórdia:\\
Carnis terat supérbiam\\
Potus cibíque párcitas.

Ut, cum dies abscésserit\\
Noctémque sors redúxerit,\\
Mundi per abstinéntiam\\
Ipsi canámus glóriam.

* Deo Patri sit glória,\\
Ejúsque soli Fílio,\\
Cum Spíritu Paráclito,\\
Nunc et per omne séculum.\\
Amen.

\rubric{In Adventu: Antiphona, etiam ad Tertiam, Sextam et Nonam, e Laudibus Dominicæ currentis, ut in Proprio de Tempore.}

\rubric{Per Annum:} \antlabel{Ant.} Allelúja.

\rubric{Antiphona cum Psalmis, per Annum ad Horas minores assignata, adhibetur semper, extra Octavam Nativitatis, quando fit Officium de qualibet Dominica minori.}

\psalmheading{Psalmus 117}

\lettrine[lines=2]{C}{ONFITEMINI} Dómino quóniam bonus: * quóniam in séculum misericórdia ejus.

Dicat nunc Israël quóniam bonus: * quóniam in séculum misericórdia ejus.

Dicat nunc domus Aaron: * quóniam in séculum misericórdia ejus.

Dicant nunc qui timent Dóminum: * quóniam in séculum misericórdia ejus.

De tribulatióne invocávi Dóminum: * et exaudívit me in latitúdine Dóminus.

Dóminus mihi adjútor: * non timébo quid fáciat mihi homo.

Dóminus mihi adjútor: * et ego despíciam inimícos meos.

Bonum est confídere in Dómino, * quam confídere in hómine:

Bonum est speráre in Dómino, * quam speráre in princípibus.

Omnes gentes circuiérunt me: * et in nómine Dómini quia ultus sum in eos.

Circumdántes circumdedérunt me: * et in nómine Dómini quia ultus sum in eos.

Circumdedérunt me sicut apes, et exarsérunt sicut ignis in spinis: * et in nómine Dómini quia ultus sum in eos.

Impúlsus evérsus sum ut cáderem: * et Dóminus suscépit me.

Fortitúdo mea, et laus mea Dóminus: * et factus est mihi in salútem.

Vox exsultatiónis, et salútis * in tabernáculis justórum.

Déxtera Dómini fecit virtútem: déxtera Dómini exaltávit me, * déxtera Dómini fecit virtútem.

Non móriar, sed vivam: * et narrábo ópera Dómini.

Castígans castigávit me Dóminus: * et morti non trádidit me.

Aperíte mihi portas justítiæ, ingréssus in eas confitébor Dómino: * hæc porta Dómini, justi intrábunt in eam.

Confitébor tibi quóniam exaudísti me: * et factus es mihi in salútem.

Lápidem, quem reprobavérunt ædificántes: * hic factus est in caput ánguli.

A Dómino factum est istud: * et est mirábile in óculis nostris.

Hæc est dies, quam fecit Dóminus: * exsultémus, et lætémur in ea.

O Dómine, salvum me fac, o Dómine, bene prosperáre: * benedíctus qui venit in nómine Dómini.

Benedíximus vobis de domo Dómini: * Deus Dóminus, et illúxit nobis.

Constitúite diem solémnem in condénsis, * usque ad cornu altáris.

Deus meus es tu, et confitébor tibi: * Deus meus es tu, et exaltábo te.

Confitébor tibi quóniam exaudísti me: * et factus es mihi in salútem.

Confitémini Dómino quóniam bonus: * quóniam in séculum misericórdia ejus.

\rubric{In Dominicis Septuagesimæ, Sexagesimæ et Quinquagesimæ, Antiphona, etiam ad Tertiam, Sextam et Nonam, ut in Proprio de Tempore; atque loco præcedentis Ps. 117 Confitémini, dicuntur sequentes Ps. 92 et 99, qui in Laudibus locum non habuérunt:}

\psalmheading{Psalmus 92}

\lettrine[lines=2]{D}{OMINUS} regnávit, decórem indútus est: * indútus est Dóminus fortitúdinem, et præcínxit se.

Etenim firmávit orbem terræ, * qui non commovébitur.

Paráta sedes tua ex tunc: * a século tu es.

Elevavérunt flúmina, Dómine: * elevavérunt flúmina vocem suam.

Elevavérunt flúmina fluctus suos, * a vócibus aquárum multárum.

Mirábiles elatiónes maris: * mirábilis in altis Dóminus.

Testimónia tua credibília facta sunt nimis: * domum tuam decet sanctitúdo, Dómine, in longitúdinem diérum.

\psalmheading{Psalmus 99}

\lettrine[lines=2]{J}{UBILATE} Deo, omnis terra: * servíte Dómino in lætítia.

Introíte in conspéctu ejus, * in exsultatióne.

Scitóte quóniam Dóminus ipse est Deus: * ipse fecit nos, et non ipsi nos.

Pópulus ejus, et oves páscuæ ejus: * introíte portas ejus in confessióne, átria ejus in hymnis: confitémini illi.

Laudáte nomen ejus: quóniam suávis est Dóminus, in ætérnum misericórdia ejus, * et usque in generatiónem et generatiónem véritas ejus.

\rubric{In Vigiliis Nativitatis Domini et Epiphaniæ, in Dominica, etiam translata, infra Octavam Nativitatis, in Duplicibus I vel II classis, in Festis novem Lectionum Domini, beátæ Mariæ Virginis, Angelórum, sancti Joannis Baptístæ, sancti Joseph, Apostolórum, Evangelistárum, si de eis fiat Officium, ad Primam, Tertiam, Sextam et Nonam dicuntur Psalmi de Dominica, sub Antiphonis Officio diei congruentibus, nisi aliter suo loco notetur; sed ad Primam, loco Psalmi 117 Confitémini, in ipsa etiam die Dominica, semper dicitur}
% [source brace unclosed at this point in hiemalis.md; rubric
%  editorially closed here, matching the parallel "dicitur sequens"
%  constructions used throughout the Psalterium -- flagged for review]

\psalmheading{Psalmus 53}

\lettrine[lines=2]{D}{EUS}, in nómine tuo salvum me fac: * et in virtúte tua júdica me.

Deus, exáudi oratiónem meam: * áuribus pércipe verba oris mei.

Quóniam aliéni insurrexérunt advérsum me, et fortes quæsiérunt ánimam meam: * et non proposuérunt Deum ante conspéctum suum.

Ecce enim, Deus ádjuvat me: * et Dóminus suscéptor est ánimæ meæ.

Avérte mala inimícis meis: * et in veritáte tua dispérde illos.

Voluntárie sacrificábo tibi, * et confitébor nómini tuo, Dómine: quóniam bonum est:

Quóniam ex omni tribulatióne eripuísti me: * et super inimícos meos despéxit óculus meus.

\psalmheading{Psalmus 118, i}

\lettrine[lines=2]{B}{EATI} immaculáti in via: * qui ámbulant in lege Dómini.

Beáti, qui scrutántur testimónia ejus: * in toto corde exquírunt eum.

Non enim qui operántur iniquitátem, * in viis ejus ambulavérunt.

Tu mandásti * mandáta tua custodíri nimis.

Utinam dirigántur viæ meæ, * ad custodiéndas justificatiónes tuas!

Tunc non confúndar, * cum perspéxero in ómnibus mandátis tuis.

Confitébor tibi in directióne cordis: * in eo quod dídici judícia justítiæ tuæ.

Justificatiónes tuas custódiam: * non me derelínquas usquequáque.

In quo córrigit adolescéntior viam suam? * in custodiéndo sermónes tuos.

In toto corde meo exquisívi te: * ne repéllas me a mandátis tuis.

In corde meo abscóndi elóquia tua: * ut non peccem tibi.

Benedíctus es, Dómine: * doce me justificatiónes tuas.

In lábiis meis, * pronuntiávi ómnia judícia oris tui.

In via testimoniórum tuórum delectátus sum, * sicut in ómnibus divítiis.

In mandátis tuis exercébor: * et considerábo vias tuas.

In justificatiónibus tuis meditábor: * non oblivíscar sermónes tuos.

\psalmheading{Psalmus 118, ii}

\lettrine[lines=2]{R}{ETRIBUE} servo tuo, vivífica me: * et custódiam sermónes tuos.

Revéla óculos meos: * et considerábo mirabília de lege tua.

Íncola ego sum in terra: * non abscóndas a me mandáta tua.

Concupívit ánima mea desideráre justificatiónes tuas, * in omni témpore.

Increpásti supérbos: * maledícti qui declínant a mandátis tuis.

Aufer a me oppróbrium, et contémptum: * quia testimónia tua exquisívi.

Etenim sedérunt príncipes, et advérsum me loquebántur: * servus autem tuus exercebátur in justificatiónibus tuis.

Nam et testimónia tua meditátio mea est: * et consílium meum justificatiónes tuæ.

Adhǽsit paviménto ánima mea: * vivífica me secúndum verbum tuum.

Vias meas enuntiávi, et exaudísti me: * doce me justificatiónes tuas.

Viam justificatiónum tuárum instrue me: * et exercébor in mirabílibus tuis.

Dormitávit ánima mea præ tǽdio: * confírma me in verbis tuis.

Viam iniquitátis ámove a me: * et de lege tua miserére mei.

Viam veritátis elégi: * judícia tua non sum oblítus.

Adhǽsi testimóniis tuis, Dómine: * noli me confúndere.

Viam mandatórum tuórum cucúrri, * cum dilatásti cor meum.

\rubric{Si Athanasianum Symbolum omittatur, expletis Psalmis præcedentibus, statim dicitur Antiphona, ut infra. 58}

\rubric{Sequens Symbolum dicitur in Dominicis tantum minoribus post Epiphaniam, quando fit Officium de Dominica ut in Psalterio, neque ulla Commemoratio Officii duplicis vel Octavæ recurrit.}

\symbolumheading{Symbolum Athanasianum}

QUICUMQUE vult salvus esse, * ante ómnia opus est, ut téneat cathólicam fidem:

Quam nisi quisque íntegram inviolatámque serváverit, * absque dúbio in ætérnum períbit.

Fides autem cathólica hæc est: * ut unum Deum in Trinitáte, et Trinitátem in unitáte venerémur.

Neque confundéntes persónas, * neque substántiam separántes.

Ália est enim persóna Patris, ália Fílii, * ália Spíritus Sancti:

Sed Patris, et Fílii, et Spíritus Sancti una est divínitas, * æquális glória, coætérna majéstas.

Quális Pater, talis Fílius, * talis Spíritus Sanctus.

Increátus Pater, increátus Fílius, * increátus Spíritus Sanctus.

Imménsus Pater, imménsus Fílius, * imménsus Spíritus Sanctus.

Ætérnus Pater, ætérnus Fílius, * ætérnus Spíritus Sanctus.

Et tamen non tres ætérni, * sed unus ætérnus.

Sicut non tres increáti, nec tres imménsi, * sed unus increátus, et unus imménsus.

Simíliter omnípotens Pater, omnípotens Fílius, * omnípotens Spíritus Sanctus.

Et tamen non tres omnipoténtes, * sed unus omnípotens.

Ita Deus Pater, Deus Fílius, * Deus Spíritus Sanctus.

Et tamen non tres Dii, * sed unus est Deus.

Ita Dóminus Pater, Dóminus Fílius, * Dóminus Spíritus Sanctus.

Et tamen non tres Dómini, * sed unus est Dóminus.

Quia, sicut singillátim unamquámque persónam Deum ac Dóminum confitéri christiána veritáte compéllimur: * ita tres Deos aut Dóminos dícere cathólica religióne prohibémur.

Pater a nullo est factus: * nec creátus, nec génitus.

Fílius a Patre solo est: * non factus, nec creátus, sed génitus.

Spíritus Sanctus a Patre et Fílio: * non factus, nec creátus, nec génitus, sed procédens.

Unus ergo Pater, non tres Patres: unus Fílius, non tres Fílii: * unus Spíritus Sanctus, non tres Spíritus Sancti.

Et in hac Trinitáte nihil prius aut postérius, nihil majus aut minus: * sed totæ tres persónæ coætérnæ sibi sunt et coæquáles.

Ita ut per ómnia, sicut jam supra dictum est, * et únitas in Trinitáte, et Trínitas in unitáte veneránda sit.

Qui vult ergo salvus esse, * ita de Trinitáte séntiat.

Sed necessárium est ad ætérnam salútem, * ut Incarnatiónem quoque Dómini nostri Jesu Christi fidéliter credat.

Est ergo fides recta, ut credámus et confiteámur, * quia Dóminus noster Jesus Christus, Dei Fílius, Deus et homo est.

Deus est ex substántia Patris ante sécula génitus: * et homo est ex substántia matris in século natus.

Perféctus Deus, perféctus homo: * ex ánima rationáli et humána carne subsístens.

Æquális Patri secúndum divinitátem: * minor Patre secúndum humanitátem.

Qui, licet Deus sit et homo, * non duo tamen, sed unus est Christus.

Unus autem non conversióne divinitátis in carnem, * sed assumptióne humanitátis in Deum.

Unus omníno, non confusióne substántiæ, * sed unitáte persónæ.

Nam sicut ánima ratiónális et caro unus est homo: * ita Deus et homo unus est Christus.

Qui passus est pro salúte nostra: descéndit ad ínferos: * tértia die resurréxit a mórtuis.

Ascéndit ad cælos, sedet ad déxteram Dei Patris omnipoténtis: * inde ventúrus est judicáre vivos et mórtuos.

Ad cujus advéntum omnes hómines resúrgere habent cum corpóribus suis: * et redditúri sunt de factis própriis ratiónem.

Et qui bona egérunt, ibunt in vitam ætérnam: * qui vero mala, in ignem ætérnum.

Hæc est fides cathólica, * quam nisi quisque fidéliter firmitérque credíderit, salvus esse non póterit.

Glória Patri.

\rubric{Expleto Symbolo, vel, si hoc non habeat locum, Psalmis absolutis, statim dicitur Antiphona.}

\rubric{Per Annum:} \antlabel{Ant.} Allelúja, confitémini Dómino quóniam in séculum misericórdia ejus, allelúja, allelúja.

\capitulumlabel{Capitulum}{1 Tim. 1, 17}

\lettrine[lines=2]{R}{EGI} sæculórum immortáli et invisíbili, soli Deo honor et glória in sécula sæculórum. Amen.

\R\ br. Christe, Fili Dei vivi, * Miserére nobis. Christe. \V\ \textit{Qui sedes ad déxteram Patris.} Miserére nobis. Glória Patri. Christe. \V\ Exsúrge, Christe, ádjuva nos. \R\ Et líbera nos propter nomen tuum.

\rubric{In Offic. Dom. Adventus:} \V\ Qui ventúrus es in mundum.

\variantheading{Preces (si dicendæ sint)}

KYRIE, eléison. Christe, eléison. Kýrie, eléison.

Pater noster \rubric{secreto usque ad} \V\ Et ne nos indúcas in tentatiónem. \R\ Sed líbera nos a malo.

Credo in Deum \rubric{secreto usque ad} \V\ Carnis resurrectiónem. \R\ Vitam ætérnam. Amen. \V\ Et ego ad te, Dómine, clamávi. \R\ Et mane orátio mea prævéniet te. \V\ Repleátur os meum laude. \R\ Ut cantem glóriam tuam, tota die magnitúdinem tuam. \V\ Dómine, avérte fáciem tuam a peccátis meis. \R\ Et omnes iniquitátes meas dele. \V\ Cor mundum crea in me, Deus. \R\ Et spíritum rectum ínnova in viscéribus meis. \V\ Ne projícias me a fácie tua. \R\ Et spíritum sanctum tuum ne áuferas a me. \V\ Redde mihi lætítiam salutáris tui. \R\ Et spíritu principáli confírma me. \V\ Adjutórium nostrum in nómine Dómini. \R\ Qui fecit cælum et terram.

\formulalabel{Confíteor.} Misereátur. Indulgéntiam.

\V\ Dignáre, Dómine, die isto. \R\ Sine peccáto nos custodíre. \V\ Miserére nostri, Dómine. \R\ Miserére nostri. \V\ Fiat misericórdia tua, Dómine, super nos. \R\ Quemádmodum sperávimus in te. \V\ Dómine, exáudi oratiónem meam. \R\ Et clamor meus ad te véniat.

\V\ Dóminus vobíscum.

Orémus. \oratiolabel{Oratio}

\lettrine[lines=2]{D}{OMINE}, Deus omnípotens, qui ad princípium hujus diéi nos perveníre fecísti: tua nos hódie salva virtúte; ut in hac die ad nullum declinémus peccátum, sed semper ad tuam justítiam faciéndam nostra procédant elóquia, dirigántur cogitatiónes et ópera. Per Dóminum. \R\ Amen.

\V\ Dóminus vobíscum. \V\ Benedicámus Dómino.

\rubric{Deinde in Choro legitur Martyrologium, quod laudabíliter fit etiam extra Chorum.}

\rubric{Postea Hebdomad. dicit:}

\V\ Pretiósa in conspéctu Dómini. \R\ Mors Sanctórum ejus.

\rubric{Et absolute, sine Orémus:}

\oratiolabel{Oratio}

\lettrine[lines=2]{S}{ANCTA} María, et omnes Sancti intercédant pro nobis ad Dóminum, ut nos mereámur ab eo adjuvári et salvári, qui vivit et regnat in sécula sæculórum. \R\ Amen.

\V\ Deus, in adjutórium meum inténde. \R\ Dómine, ad adjuvándum me festína. \rubric{Et dicitur ter; ultimo additur:} Glória Patri. Sicut erat.

Kýrie, eléison. Christe, eléison. Kýrie, eléison.

Pater noster \rubric{secreto usque ad} \V\ Et ne nos indúcas in tentatiónem. \R\ Sed líbera nos a malo. \V\ Réspice in servos tuos, Dómine, et in ópera tua, et dírige fílios eórum. \R\ Et sit splendor Dómini, Dei nostri, super nos, et ópera mánuum nostrárum dírige super nos, et opus mánuum nostrárum dírige. \V\ Glória Patri. \R\ Sicut erat.

Orémus. \oratiolabel{Oratio}

\lettrine[lines=2]{D}{IRIGERE} et sanctificáre, régere et gubernáre dignáre, Dómine Deus, Rex cæli et terræ, hódie corda et córpora nostra, sensus, sermónes et actus nostros in lege tua, et in opéribus mandatórum tuórum: ut hic, et in ætérnum, te auxiliánte, salvi et líberi esse mereámur, Salvátor mundi: Qui vivis et regnas in sécula sæculórum. \R\ Amen.

\V\ Jube, domne.

\benedlabel{Bened.} Dies et actus nostros in sua pace dispónat Dóminus omnípotens. \R\ Amen.

\sectionornament

\variantheading{In Adventu}

\lettrine[lines=2]{D}{OMINE}, miserére nostri: te enim exspectávimus: esto bráchium nostrum in mane, et salus nostra in témpore tribulatiónis. Tu autem.

\rubric{In Vigilia Nativitatis Domini Lectio br. Per quem, ut infra.}

\variantheading{Per Annum}

\lettrine[lines=2]{D}{OMINUS} autem dírigat corda et córpora nostra in caritáte Dei, et patiéntia Christi.

\sectionornament

\sectionheading{E Proprio de Tempore}

\feastheading{In Vigilia Nativitatis}

\lettrine[lines=2]{P}{ER} quem accépimus grátiam, et apostolátum ad obediéndum fídei in ómnibus géntibus pro nómine ejus, in quibus estis et vos vocáti Jesu Christi, Dómini nostri.

\feastheading{In Nativitate Domini et in die VI infra Octavam}

\lettrine[lines=2]{I}{PSI} períbunt, tu autem permanébis; et omnes sicut vestiméntum veteráscent: et velut amíctum mutábis eos, et mutabúntur: tu autem idem ipse es, et anni tui non defícient.

\feastheading{S. Stephani Protomartyris}

\lettrine[lines=2]{P}{OSITIS} autem génibus, clamávit voce magna, dicens: Dómine, ne státuas illis hoc peccátum. Et cum hoc dixísset, obdormívit in Dómino.

\feastheading{S. Joannis Ap. et Ev.}

\lettrine[lines=2]{I}{N} médio Ecclésiæ apéruit os ejus, et implévit illum Dóminus spíritu sapiéntiæ et intelléctus, et stolam glóriæ índuit eum.

\feastheading{Ss. Innocentium Mm.}

\lettrine[lines=2]{H}{I} empti sunt ex homínibus primítiæ Deo et Agno, et in ore ipsórum non est invéntum mendácium: sine mácula enim sunt ante thronum Dei.

\feastheading{Dom. infra Oct. Nativ.}

\lettrine[lines=2]{I}{TAQUE} jam non est servus, sed fílius: quod si fílius, et heres per Deum.

\feastheading{In Circumcisione Domini}

\lettrine[lines=2]{I}{PSI} períbunt, tu autem permanébis; et omnes sicut vestiméntum veteráscent: et velut amíctum mutábis eos, et mutabúntur: tu autem idem ipse es, et anni tui non defícient.

\feastheading{Ssmi Nominis Jesu}

\lettrine[lines=2]{O}{MNE} quodcúmque fácitis in verbo, aut in ópere: ómnia in nómine Dómini Jesu Christi fácite, grátias agéntes Deo et Patri per ipsum.

\feastheading{In Vigilia Epiphaniæ}

\lettrine[lines=2]{I}{TAQUE} jam non est servus, sed fílius: quod si fílius, et heres per Deum.

\feastheading{In Epiphania et per Oct.}

\lettrine[lines=2]{O}{MNES} de Saba vénient, aurum et thus deferéntes, et laudem Dómino annuntiántes.

\feastheading{Sanctæ Familiæ J. M. J.}

\lettrine[lines=2]{S}{EMETIPSUM} exinanívit formam servi accípiens, in similitúdinem hóminum factus, et hábitu invéntus ut homo.

\sectionornament

\sectionheading{E Proprio Sanctorum}

\feastheading{S. Andreæ Apost.}

\lettrine[lines=2]{I}{SAIAS} enim dicit: Dómine, quis crédidit audítui nostro? Ergo fides ex audítu, audítus autem per verbum Christi. Sed dico: Numquid non audiérunt? Et quidem in omnem terram exívit sonus eórum, et in fines orbis terræ verba eórum.

\feastheading{In Conceptione immacul. B. Mariæ V.}

\lettrine[lines=2]{S}{IGNUM} magnum appáruit in cælo: Múlier amícta sole, et luna sub pédibus ejus, et in cápite ejus coróna stellárum duódecim.

\feastheading{S. Agnetis et S. Agathæ}

\lettrine[lines=2]{L}{AUDABIT} usque ad mortem ánima mea Dóminum, quóniam éruis sustinéntes te, et líberas eos de manu angústiæ, Dómine, Deus noster.

\feastheading{In Conversione S. Pauli}

\lettrine[lines=2]{S}{AULUS} autem multo magis convalescébat, et confundébat Judǽos, qui habitábant Damásci, affírmans, quóniam hic est Christus.

\feastheading{In Purificatione B. M. V.}

\lettrine[lines=2]{E}{T} placébit Dómino sacrifícium Juda et Jerúsalem sicut dies séculi et sicut anni antíqui, dicit Dóminus omnípotens.

\feastheading{In Apparitione B. M. V.}

\lettrine[lines=2]{I}{N} me grátia omnis viæ et veritátis, in me omnis spes vitæ et virtútis. Transíte ad me, omnes qui concupíscitis me, et a generatiónibus meis implémini.

\sectionornament

\sectionheading{E Communi Sanctorum}

\feastheading{Pro Apostolis}

\lettrine[lines=2]{I}{BANT} Apóstoli gaudéntes a conspéctu concílii, quóniam digni habíti sunt pro nómine Jesu contuméliam pati.

\feastheading{Pro Martyre}

\lettrine[lines=2]{J}{USTUS} cor suum trádidit ad vigilándum dilúculo ad Dóminum, qui fecit illum, et in conspéctu Altíssimi deprecábitur.

\feastheading{Pro pluribus Martyribus}

\lettrine[lines=2]{F}{ULGEBUNT} justi, et tamquam scintíllæ in arundinéto discúrrent. Judicábunt natiónes, et dominabúntur pópulis: et regnábit Dóminus illórum in perpétuum.

\feastheading{Pro Conf. Pont.}

\lettrine[lines=2]{F}{UNGI} sacerdótio, et habére laudem in nómine ipsíus, et offérre illi incénsum dignum in odórem suavitátis.

\feastheading{Pro Conf. non Pont.}

\lettrine[lines=2]{J}{USTUM} dedúxit Dóminus per vias rectas, et osténdit illi regnum Dei, et dedit illi sciéntiam sanctórum; honestávit illum in labóribus, et complévit labóres illíus.

\feastheading{Pro Virgine Martyre}

\lettrine[lines=2]{D}{OMINE}, Deus meus, exaltásti super terram habitatiónem meam, et pro morte defluénte deprecáta sum. Invocávi Dóminum, Patrem Dómini mei, ut non derelínquat me in die tribulatiónis meæ, et in témpore superbórum sine adjutório.

\feastheading{Pro Virgine non Martyre}

\lettrine[lines=2]{O}{} \textsc{quam} pulchra est casta generátio cum claritáte! immortális est enim memória illíus: quóniam apud Deum nota est, et apud hómines.

\feastheading{Pro Martyre tantum}

\lettrine[lines=2]{L}{AUDABIT} usque ad mortem ánima mea Dóminum, quóniam éruis sustinéntes te, et líberas eos de manu angústiæ, Dómine, Deus noster.

\feastheading{Pro nec Virg. nec Mart.}

\lettrine[lines=2]{M}{ULTÆ} fíliæ congregavérunt divítias: tu supergréssa es univérsas. Fallax grátia, et vana est pulchritúdo: múlier timens Dóminum, ipsa laudábitur.

\feastheading{Pro Dedicatione Ecclesiæ}

\lettrine[lines=2]{E}{T} abstérget Deus omnem lácrimam ab óculis eórum: et mors ultra non erit, neque luctus, neque clamor, neque dolor erit ultra, quia prima abiérunt. Et dixit qui sedébat in throno: Ecce, nova fácio ómnia.

\feastheading{In Festis B. Mariæ V.}

\lettrine[lines=2]{I}{N} platéis sicut cinnamómum et bálsamum aromatízans odórem dedi: quasi myrrha elécta dedi suavitátem odóris.

\sectionornament

\V\ Adjutórium nostrum in nómine Dómini. \R\ Qui fecit cælum et terram.

\V\ Benedícite. \R\ Deus.

\benedictiolabel{Benedictio} Dóminus nos benedícat, et ab omni malo deféndat, et ad vitam perdúcat ætérnam. Et fidélium ánimæ per misericórdiam Dei requiéscant in pace. \R\ Amen.

\rubric{Deinde dicitur tantum Pater noster secreto, nisi sequatur alia Hora.}

\sectionornament

\hourheading{Ad Tertiam}

Pater et Ave.

\V\ Deus, in adjutórium.

\hymnusheading{Hymnus}

\lettrine[lines=2]{N}{UNC}, Sancte, nobis, Spíritus,\\
Unum Patri cum Fílio,\\
Dignáre promptus íngeri\\
Nostro refúsus péctori.

Os, lingua, mens, sensus, vigor\\
Confessiónem pérsonent,\\
Flamméscat igne cáritas,\\
Accéndat ardor próximos.

* Præsta, Pater piíssime,\\
Patríque compar Unice,\\
Cum Spíritu Paráclito\\
Regnans per omne séculum.\\
Amen.

\rubric{Per Annum:} \antlabel{Ant.} Allelúja.

\psalmheading{Psalmus 118, iii}

\lettrine[lines=2]{L}{EGEM} pone mihi, Dómine, viam justificatiónum tuárum: * et exquíram eam semper.

Da mihi intelléctum, et scrutábor legem tuam: * et custódiam illam in toto corde meo.

Deduc me in sémitam mandatórum tuórum: * quia ipsam vólui.

Inclína cor meum in testimónia tua: * et non in avarítiam.

Avérte óculos meos ne vídeant vanitátem: * in via tua vivífica me.

Státue servo tuo elóquium tuum, * in timóre tuo.

Ampúta oppróbrium meum quod suspicátus sum: * quia judícia tua jucúnda.

Ecce, concupívi mandáta tua: * in æquitáte tua vivífica me.

Et véniat super me misericórdia tua, Dómine: * salutáre tuum secúndum elóquium tuum.

Et respondébo exprobrántibus mihi verbum: * quia sperávi in sermónibus tuis.

Et ne áuferas de ore meo verbum veritátis usquequáque: * quia in judíciis tuis supersperávi.

Et custódiam legem tuam semper: * in séculum et in séculum séculi.

Et ambulábam in latitúdine: * quia mandáta tua exquisívi.

Et loquébar in testimóniis tuis in conspéctu regum: * et non confundébar.

Et meditábar in mandátis tuis, * quæ diléxi.

Et levávi manus meas ad mandáta tua, quæ diléxi: * et exercébar in justificatiónibus tuis.

\psalmheading{Psalmus 118, iv}

\lettrine[lines=2]{M}{EMOR} esto verbi tui servo tuo, * in quo mihi spem dedísti.

Hæc me consoláta est in humilitáte mea: * quia elóquium tuum vivificábit me.

Supérbi iníque agébant usquequáque: * a lege autem tua non declinávi.

Memor fui judiciórum tuórum a século, Dómine: * et consolátus sum.

Deféctio ténuit me, * pro peccatóribus derelinquéntibus legem tuam.

Cantábiles mihi erant justificatiónes tuæ, * in loco peregrinatiónis meæ.

Memor fui nocte nóminis tui, Dómine: * et custodívi legem tuam.

Hæc facta est mihi: * quia justificatiónes tuas exquisívi.

Pórtio mea, Dómine, * dixi custodíre legem tuam.

Deprecátus sum fáciem tuam in toto corde meo: * miserére mei secúndum elóquium tuum.

Cogitávi vias meas: * et convérti pedes meos in testimónia tua.

Parátus sum, et non sum turbátus: * ut custódiam mandáta tua.

Funes peccatórum circumpléxi sunt me: * et legem tuam non sum oblítus.

Média nocte surgébam ad confiténdum tibi, * super judícia justificatiónis tuæ.

Párticeps ego sum ómnium timéntium te: * et custodiéntium mandáta tua.

Misericórdia tua, Dómine, plena est terra: * justificatiónes tuas doce me.

\psalmheading{Psalmus 118, v}

\lettrine[lines=2]{B}{ONITATEM} fecísti cum servo tuo, Dómine, * secúndum verbum tuum.

Bonitátem, et disciplínam, et sciéntiam doce me: * quia mandátis tuis crédidi.

Priúsquam humiliárer ego delíqui: * proptérea elóquium tuum custodívi.

Bonus es tu: * et in bonitáte tua doce me justificatiónes tuas.

Multiplicáta est super me iníquitas superbórum: * ego autem in toto corde meo scrutábor mandáta tua.

Coagulátum est sicut lac cor eórum: * ego vero legem tuam meditátus sum.

Bonum mihi quia humiliásti me: * ut discam justificatiónes tuas.

Bonum mihi lex oris tui, * super míllia auri et argénti.

Manus tuæ fecérunt me, et plasmavérunt me: * da mihi intelléctum, et discam mandáta tua.

Qui timent te vidébunt me, et lætabúntur: * quia in verba tua supersperávi.

Cognóvi, Dómine, quia æquitas judícia tua: * et in veritáte tua humiliásti me.

Fiat misericórdia tua ut consolétur me, * secúndum elóquium tuum servo tuo.

Véniant mihi miseratiónes tuæ, et vivam: * quia lex tua meditátio mea est.

Confundántur supérbi, quia injúste iniquitátem fecérunt in me: * ego autem exercébor in mandátis tuis.

Convertántur mihi timéntes te: * et qui novérunt testimónia tua.

Fiat cor meum immaculátum in justificatiónibus tuis, * ut non confúndar.

\variantheading{Per Annum}

\antlabel{Ant.} Allelúja, deduc me, Dómine, in sémitam mandatórum tuórum, allelúja, allelúja.

\capitulumlabel{Capitulum}{1 Joann. 4, 16}

\lettrine[lines=2]{D}{EUS} cáritas est: et qui manet in caritáte, in Deo manet, et Deus in eo.

\R\ br. Inclína cor meum, Deus, * In testimónia tua. Inclína. \V\ Avérte óculos meos, ne vídeant vanitátem: in via tua vivífica me. In testimónia tua. Glória Patri. Inclína. \V\ Ego dixi: Dómine, miserére mei. \R\ Sana ánimam meam, quia peccávi tibi.

\rubric{Oratio convéniens.}

\rubric{Reliqua ut in Ordinario. 20}

\sectionornament

\hymnusheading{Hymnus}

\lettrine[lines=2]{R}{ECTOR} potens, verax Deus,\\
Qui témperas rerum vices,\\
Splendóre mane illúminas,\\
Et ígnibus merídiem:

Exstíngue flammas lítium,\\
Aufer calórem nóxium,\\
Confer salútem córporum,\\
Verámque pacem córdium.

* Præsta, Pater piíssime,\\
Patríque compar Unice,\\
Cum Spíritu Paráclito\\
Regnans per omne séculum.\\
Amen.

\rubric{Per Annum:} \antlabel{Ant.} Allelúja.

\psalmheading{Psalmus 118, vi}

\lettrine[lines=2]{D}{EFECIT} in salutáre tuum ánima mea: * et in verbum tuum supersperávi.

Defecérunt óculi mei in elóquium tuum, * dicéntes: Quando consoláberis me?

Quia factus sum sicut uter in pruína: * justificatiónes tuas non sum oblítus.

Quot sunt dies servi tui? * quando fácies de persequéntibus me judícium?

Narravérunt mihi iníqui fabulatiónes: * sed non ut lex tua.

Omnia mandáta tua véritas: * iníque persecúti sunt me, ádjuva me.

Paulo minus consummavérunt me in terra: * ego autem non derelíqui mandáta tua.

Secúndum misericórdiam tuam vivífica me: * et custódiam testimónia oris tui.

In ætérnum, Dómine, * verbum tuum pérmanet in cælo.

In generatiónem et generatiónem véritas tua: * fundásti terram, et pérmanet.

Ordinatióne tua persevérat dies: * quóniam ómnia sérviunt tibi.

Nisi quod lex tua meditátio mea est: * tunc forte periíssem in humilitáte mea.

In ætérnum non oblivíscar justificatiónes tuas: * quia in ipsis vivificásti me.

Tuus sum ego, salvum me fac: * quóniam justificatiónes tuas exquisívi.

Me exspectavérunt peccatóres ut pérderent me: * testimónia tua intelléxi.

Omnis consummatiónis vidi finem: * latum mandátum tuum nimis.

\psalmheading{Psalmus 118, vii}

\lettrine[lines=2]{Q}{UOMODO} diléxi legem tuam, Dómine? * tota die meditátio mea est.

Super inimícos meos prudéntem me fecísti mandáto tuo: * quia in ætérnum mihi est.

Super omnes docéntes me intelléxi: * quia testimónia tua meditátio mea est.

Super senes intelléxi: * quia mandáta tua quæsívi.

Ab omni via mala prohíbui pedes meos: * ut custódiam verba tua.

A judíciis tuis non declinávi: * quia tu legem posuísti mihi.

Quam dúlcia fáucibus meis elóquia tua, * super mel ori meo!

A mandátis tuis intelléxi: * proptérea odívi omnem viam iniquitátis.

Lucérna pédibus meis verbum tuum, * et lumen sémitis meis.

Jurávi, et státui * custodíre judícia justítiæ tuæ.

Humiliátus sum usquequáque, Dómine: * vivífica me secúndum verbum tuum.

Voluntária oris mei beneplácita fac, Dómine: * et judícia tua doce me.

Anima mea in mánibus meis semper: * et legem tuam non sum oblítus.

Posuérunt peccatóres láqueum mihi: * et de mandátis tuis non errávi.

Hereditáte acquisívi testimónia tua in ætérnum: * quia exsultátio cordis mei sunt.

Inclinávi cor meum ad faciéndas justificatiónes tuas in ætérnum, * propter retributiónem.

\psalmheading{Psalmus 118, viii}

\lettrine[lines=2]{I}{NIQUOS} ódio hábui: * et legem tuam diléxi.

Adjútor et suscéptor meus es tu: * et in verbum tuum supersperávi.

Declináte a me, malígni: * et scrutábor mandáta Dei mei.

Súscipe me secúndum elóquium tuum, et vivam: * et non confúndas me ab exspectatióne mea.

Adjuva me, et salvus ero: * et meditábor in justificatiónibus tuis semper.

Sprevísti omnes discedéntes a judíciis tuis: * quia injústa cogitátio eórum.

Prævaricántes reputávi omnes peccatóres terræ: * ídeo diléxi testimónia tua.

Confíge timóre tuo carnes meas: * a judíciis enim tuis tímui.

Feci judícium et justítiam: * non tradas me calumniántibus me.

Súscipe servum tuum in bonum: * non calumniéntur me supérbi.

Óculi mei defecérunt in salutáre tuum: * et in elóquium justítiæ tuæ.

Fac cum servo tuo secúndum misericórdiam tuam: * et justificatiónes tuas doce me.

Servus tuus sum ego: * da mihi intelléctum, ut sciam testimónia tua.

Tempus faciéndi, Dómine: * dissipavérunt legem tuam.

Ideo diléxi mandáta tua, * super aurum et topázion.

Proptérea ad ómnia mandáta tua dirigébar: * omnem viam iníquam ódio hábui.

\variantheading{Per Annum}

\antlabel{Ant.} Allelúja, tuus sum ego, salvum me fac, Dómine, allelúja, allelúja.

\capitulumlabel{Capitulum}{Gal. 6, 2}

\lettrine[lines=2]{A}{LTER} altérius ónera portáte, et sic adimplébitis legem Christi.

\R\ br. In ætérnum, Dómine, * Pérmanet verbum tuum. In ætérnum. \V\ In séculum séculi véritas tua. Pérmanet. Glória Patri. In ætérnum. \V\ Dóminus regit me, et nihil mihi déerit. \R\ In loco páscuæ ibi me collocávit.

\rubric{Oratio convéniens.}

\rubric{Reliqua ut in Ordinario. 21}

\sectionornament

\hourheading{Ad Nonam}

Pater noster et Ave.

\V\ Deus, in adjutórium.

\hymnusheading{Hymnus}

\lettrine[lines=2]{R}{ERUM}, Deus, tenax vigor,\\
Immótus in te pérmanens,\\
Lucis diúrnæ témpora\\
Successíbus detérminans:

Largíre lumen véspere,\\
Quo vita nusquam décidat,\\
Sed prǽmium mortis sacræ\\
Perénnis instet glória.

* Præsta, Pater piíssime,\\
Patríque compar Unice,\\
Cum Spíritu Paráclito\\
Regnans per omne séculum.\\
Amen.

\rubric{Per Annum:} \antlabel{Ant.} Allelúja.

\psalmheading{Psalmus 118, ix}

\lettrine[lines=2]{M}{IRABILIA} testimónia tua: * ídeo scrutáta est ea ánima mea.

Declarátio sermónum tuórum illúminat: * et intelléctum dat párvulis.

Os meum apérui, et attráxi spíritum: * quia mandáta tua desiderábam.

Aspice in me, et miserére mei, * secúndum judícium diligéntium nomen tuum.

Gressus meos dírige secúndum elóquium tuum: * et non dominétur mei omnis injustítia.

Rédime me a calúmniis hóminum: * ut custódiam mandáta tua.

Fáciem tuam illúmina super servum tuum: * et doce me justificatiónes tuas.

Éxitus aquárum deduxérunt óculi mei: * quia non custodiérunt legem tuam.

Justus es, Dómine: * et rectum judícium tuum.

Mandásti justítiam testimónia tua: * et veritátem tuam nimis.

Tabéscere me fecit zelus meus: * quia oblíti sunt verba tua inimíci mei.

Ignítum elóquium tuum veheménter: * et servus tuus diléxit illud.

Adolescéntulus sum ego et contémptus: * justificatiónes tuas non sum oblítus.

Justítia tua, justítia in ætérnum: * et lex tua véritas.

Tribulátio, et angústia invenérunt me: * mandáta tua meditátio mea est.

Æquitas testimónia tua in ætérnum: * intelléctum da mihi, et vivam.

\psalmheading{Psalmus 118, x}

\lettrine[lines=2]{C}{LAMAVI} in toto corde meo, exáudi me, Dómine: * justificatiónes tuas requíram.

Clamávi ad te, salvum me fac: * ut custódiam mandáta tua.

Prævéni in maturitáte, et clamávi: * quia in verba tua supersperávi.

Prævenérunt óculi mei ad te dilúculo: * ut meditárer elóquia tua.

Vocem meam audi secúndum misericórdiam tuam, Dómine: * et secúndum judícium tuum vivífica me.

Appropinquavérunt persequéntes me iniquitáti: * a lege autem tua longe facti sunt.

Prope es tu, Dómine: * et omnes viæ tuæ véritas.

Inítio cognóvi de testimóniis tuis: * quia in ætérnum fundásti ea.

Vide humilitátem meam, et éripe me: * quia legem tuam non sum oblítus.

Júdica judícium meum, et rédime me: * propter elóquium tuum vivífica me.

Longe a peccatóribus salus: * quia justificatiónes tuas non exquisiérunt.

Misericórdiæ tuæ multæ, Dómine: * secúndum judícium tuum vivífica me.

Multi qui persequúntur me, et tríbulant me: * a testimóniis tuis non declinávi.

Vidi prævaricántes, et tabescébam: * quia elóquia tua non custodiérunt.

Vide quóniam mandáta tua diléxi, Dómine: * in misericórdia tua vivífica me.

Princípium verbórum tuórum, véritas: * in ætérnum ómnia judícia justítiæ tuæ.

\psalmheading{Psalmus 118, xi}

\lettrine[lines=2]{P}{RINCIPES} persecúti sunt me gratis: * et a verbis tuis formidávit cor meum.

Lætábor ego super elóquia tua: * sicut qui invénit spólia multa.

Iniquitátem ódio hábui, et abominátus sum: * legem autem tuam diléxi.

Sépties in die laudem dixi tibi, * super judícia justítiæ tuæ.

Pax multa diligéntibus legem tuam: * et non est illis scándalum.

Exspectábam salutáre tuum, Dómine: * et mandáta tua diléxi.

Custodívit ánima mea testimónia tua: * et diléxit ea veheménter.

Serváte mandáta tua, et testimónia tua: * quia omnes viæ meæ in conspéctu tuo.

Appropínquet deprecátio mea in conspéctu tuo, Dómine: * juxta elóquium tuum da mihi intelléctum.

Intret postulátio mea in conspéctu tuo: * secúndum elóquium tuum éripe me.

Eructábunt lábia mea hymnum, * cum docúeris me justificatiónes tuas.

Pronuntiábit lingua mea elóquium tuum: * quia ómnia mandáta tua æquitas.

Fiat manus tua ut salvet me: * quóniam mandáta tua elégi.

Concupívi salutáre tuum, Dómine: * et lex tua meditátio mea est.

Vivet ánima mea, et laudábit te: * et judícia tua adjuvábunt me.

Errávi, sicut ovis, quæ périit: * quǽre servum tuum, quia mandáta tua non sum oblítus.

\variantheading{Per Annum}

\antlabel{Ant.} Allelúja, fáciem tuam, Dómine, illúmina super servum tuum, allelúja, allelúja.

\capitulumlabel{Capitulum}{1 Cor. 6, 20}

\lettrine[lines=2]{E}{MPTI} enim estis prétio magno. Glorificáte et portáte Deum in córpore vestro.

\R\ br. Clamávi in toto corde meo: * Exáudi me, Dómine. Clamávi. \V\ Justificatiónes tuas requíram. Exáudi. Glória Patri. Clamávi. \V\ Ab occúltis meis munda me, Dómine. \R\ Et ab aliénis parce servo tuo.

\rubric{Oratio convéniens.}

\rubric{Reliqua ut in Ordinario. 23}

\sectionornament

\hourheading{Ad Vesperas}

Pater noster et Ave.

\V\ Deus, in adjutórium.

\rubric{In Adventu: Antiphonæ de Laudibus Dominicæ currentis, ut in Proprio de Tempore.}

\rubric{Per Annum:} \antlabel{Ant.} Dixit Dóminus. \textdagger

\rubric{Antæ cum Psalmis per Annum assignatæ adhibentur semper, quando fit Officium de Dominica, tam in Dominicis minoribus, extra Oct. Nativit., quam in Dominicis majoribus extra Adventum.}

\psalmheading{Psalmus 109}

\lettrine[lines=2]{D}{IXIT} Dóminus \textdagger\ Dómino meo: * Sede a dextris meis:

Donec ponam inimícos tuos, * scabéllum pedum tuórum.

Virgam virtútis tuæ emíttet Dóminus ex Sion: * domináre in médio inimicórum tuórum.

Tecum princípium in die virtútis tuæ in splendóribus sanctórum: * ex útero ante lucíferum génui te.

Jurávit Dóminus, et non pænitébit eum: * Tu es sacérdos in ætérnum secúndum órdinem Melchísedech.

Dóminus a dextris tuis, * confrégit in die iræ suæ reges.

Judicábit in natiónibus, implébit ruínas: * conquassábit cápita in terra multórum.

De torrénte in via bibet: * proptérea exaltábit caput.

\antlabel{Ant.} Dixit Dóminus Dómino meo: Sede a dextris meis.

\antlabel{Ant.} Magna ópera Dómini.

\psalmheading{Psalmus 110}

\lettrine[lines=2]{C}{ONFITEBOR} tibi, Dómine, in toto corde meo: * in consílio justórum, et congregatióne.

Magna ópera Dómini: * exquisíta in omnes voluntátes ejus.

Conféssio et magnificéntia opus ejus: * et justítia ejus manet in sæculum sæculi.

Memóriam fecit mirabílium suórum, miséricors et miserátor Dóminus: * escam dedit timéntibus se.

Memor erit in sæculum testaménti sui: * virtútem óperum suórum annuntiábit pópulo suo:

Ut det illis hereditátem géntium: * ópera mánuum ejus véritas, et judícium.

Fidélia ómnia mandáta ejus: confirmáta in sæculum sæculi, * facta in veritáte et æquitáte.

Redemptiónem misit pópulo suo: * mandávit in ætérnum testaméntum suum.

Sanctum, et terríbile nomen ejus: * inítium sapiéntiæ timor Dómini.

Intelléctus bonus ómnibus faciéntibus eum: * laudátio ejus manet in sæculum sæculi.

\antlabel{Ant.} Magna ópera Dómini: exquisíta in omnes voluntátes ejus.

\psalmheading{Psalmus 111}

\lettrine[lines=2]{B}{EATUS} vir, qui timet Dóminum: * in mandátis ejus volet nimis.

Potens in terra erit semen ejus: * generátio rectórum benedicétur.

Glória, et divítiæ in domo ejus: * et justítia ejus manet in sæculum sæculi.

Exórtum est in ténebris lumen rectis: * miséricors, et miserátor, et justus.

Jucúndus homo qui miserétur et cómmodat, dispónet sermónes suos in judício: * quia in ætérnum non commovébitur.

In memória ætérna erit justus: * ab auditióne mala non timébit.

Parátum cor ejus speráre in Dómino, confirmátum est cor ejus: * non commovébitur donec despíciat inimícos suos.

Dispérsit, dedit paupéribus: justítia ejus manet in sæculum sæculi, * cornu ejus exaltábitur in glória.

Peccátor vidébit, et irascétur, déntibus suis fremet et tabéscet: * desidérium peccatórum períbit.

\antlabel{Ant.} Qui timet Dóminum, in mandátis ejus cupit nimis.

\antlabel{Ant.} Sit nomen Dómini.

\psalmheading{Psalmus 112}

\lettrine[lines=2]{L}{AUDATE}, púeri, Dóminum: * laudáte nomen Dómini.

Sit nomen Dómini benedíctum, * ex hoc nunc, et usque in séculum.

A solis ortu usque ad occásum, * laudábile nomen Dómini.

Excélsus super omnes gentes Dóminus, * et super cælos glória ejus.

Quis sicut Dóminus, Deus noster, qui in altis hábitat, * et humília réspicit in cælo et in terra?

Súscitans a terra ínopem, * et de stércore érigens páuperem:

Ut cóllocet eum cum princípibus, * cum princípibus pópuli sui.

Qui habitáre facit stérilem in domo, * matrem filiórum lætántem.

\antlabel{Ant.} Sit nomen Dómini benedíctum in sécula.

\antlabel{Ant.} Deus autem noster.

\psalmheading{Psalmus 113}

\lettrine[lines=2]{I}{N} éxitu Israël de Ægýpto, * domus Jacob de pópulo bárbaro:

Facta est Judǽa sanctificátio ejus, * Israël potéstas ejus.

Mare vidit, et fugit: * Jordánis convérsus est retrórsum.

Montes exsultavérunt ut aríetes, * et colles sicut agni óvium.

Quid est tibi, mare, quod fugísti: * et tu, Jordánis, quia convérsus es retrórsum?

Montes, exsultástis sicut aríetes, * et colles, sicut agni óvium?

A fácie Dómini mota est terra, * a fácie Dei Jacob.

Qui convértit petram in stagna aquárum, * et rupem in fontes aquárum.

Non nobis, Dómine, non nobis: * sed nómini tuo da glóriam.

Super misericórdia tua, et veritáte tua: * nequándo dicant gentes: Ubi est Deus eórum?

Deus autem noster in cælo: * ómnia quæcúmque vóluit, fecit.

Simulácra géntium argéntum, et aurum, * ópera mánuum hóminum.

Os habent, et non loquéntur: * óculos habent, et non vidébunt.

Aures habent, et non áudient: * nares habent, et non odorábunt.

Manus habent, et non palpábunt: pedes habent, et non ambulábunt: * non clamábunt in gútture suo.

Símiles illis fiant qui fáciunt ea: * et omnes qui confídunt in eis.

Domus Israël sperávit in Dómino: * adjútor eórum et protéctor eórum est.

Domus Aaron sperávit in Dómino: * adjútor eórum et protéctor eórum est.

Qui timent Dóminum, speravérunt in Dómino: * adjútor eórum et protéctor eórum est.

Dóminus memor fuit nostri: * et benedíxit nobis:

Benedíxit dómui Israël: * benedíxit dómui Aaron.

Benedíxit ómnibus, qui timent Dóminum, * pusíllis cum majóribus.

Adjíciat Dóminus super vos: * super vos, et super fílios vestros.

Benedícti vos a Dómino, * qui fecit cælum, et terram.

Cælum cæli Dómino: * terram autem dedit fíliis hóminum.

Non mórtui laudábunt te, Dómine: * neque omnes, qui descéndunt in inférnum.

Sed nos qui vívimus, benedícimus Dómino, * ex hoc nunc et usque in séculum.

\variantheading{Per Annum}

\antlabel{Ant.} Deus autem noster in cælo: ómnia quæcúmque vóluit, fecit.

\variantheading{In Adventu}

\rubric{Capitulum, Hymnus et Versus ut in Proprio de Tempore.}

\variantheading{Per Annum}

\rubric{In Dominicis Septuagesimæ, Sexagesimæ et Quinquagesimæ Capitulum ut in Proprio de Tempore; in Dominicis autem post Epiphaniam a die 14 Januarii ac deinceps ante Septuagesimam occurréntibus dícitur}

\capitulumlabel{Capitulum}{2 Cor. 1, 3-4}

\lettrine[lines=2]{B}{ENEDICTUS} Deus, et Pater Dómini nostri Jesu Christi, Pater misericordiárum, et Deus totíus consolatiónis, qui consolátur nos in omni tribulatióne nostra.

\rubric{In ómnibus vero Dominicis supérius indicátis, dícitur sequens Hymnus et Versus:}

\hymnusheading{Hymnus}

\lettrine[lines=2]{L}{UCIS} Creátor óptime,\\
Lucem diérum prófereus,\\
Primórdiis lucis novæ,\\
Mundi parans oríginem:

Qui mane junctum vésperi\\
Diem vocári præcípis:\\
Illábitur tetrum chaos,\\
Audi preces cum flétibus.

Ne mens graváta crímine,\\
Vitæ sit exsul múnere,\\
Dum nil perénne cógitat,\\
Seséque culpis illigat.

Cæléste pulset óstium:\\
Vitále tollat prǽmium:\\
Vitémus omne nóxium:\\
Purgémus omne péssimum.

* Præsta, Pater piíssime,\\
Patríque compar Unice,\\
Cum Spíritu Paráclito\\
Regnans per omne séculum.\\
Amen.

\V\ Dirigátur, Dómine, orátio mea. \R\ Sicut incénsum in conspéctu tuo.

\rubric{Ad Magníficat Antiphona ut in Proprio de Tempore.}

\canticumheading{Canticum B. Mariæ Virg.}{Luc. 1, 46-55}

\lettrine[lines=2]{M}{AGNIFICAT} * ánima mea Dóminum:

Et exsultávit spíritus meus * in Deo, salutári meo.

Quia respéxit humilitátem ancíllæ suæ: * ecce enim ex hoc beátam me dicent omnes generatiónes.

Quia fecit mihi magna qui potens est: * et sanctum nomen ejus.

Et misericórdia ejus a progénie in progénies * timéntibus eum.

Fecit poténtiam in bráchio suo: * dispérsit supérbos mente cordis sui.

Depósuit poténtes de sede * et exaltávit húmiles.

Esuriéntes implévit bonis: * et dívites dimísit inánes.

Suscépit Israël, púerum suum, * recordátus misericórdiæ suæ.

Sicut locútus est ad patres nostros, * Abraham, et sémini ejus in sécula.

\rubric{Oratio convéniens.}

\rubric{Reliqua ut in Ordinario. 27}

\sectionornament

\hourheading{Ad Completorium}

\V\ Jube, domne.

\benedlabel{Bened.} Noctem quiétam, et finem perféctum concédat nobis Dóminus omnípotens. \R\ Amen.

\lectiolabel{Lectio br.}{1 Petri 5, 8-9}

\lettrine[lines=2]{F}{RATRES}: Sóbrii estóte, et vigiláte: quia adversárius vester diábolus tamquam leo rúgiens círcuit, quærens quem dévoret: cui resístite fortes in fide. Tu autem, Dómine, miserére nobis. \R\ Deo grátias.

\V\ Adjutórium nostrum in nómine Dómini. \R\ Qui fecit cælum et terram.

Pater noster, \rubric{quod dícitur totum secréto.}

\formulalabel{Confíteor. Misereátur. Indulgéntiam.}

\V\ Convérte nos, Deus, salutáris noster. \R\ Et avérte iram tuam a nobis.

\V\ Deus, in adjutórium.

\antlabel{Ant.} Miserére.

\psalmheading{Psalmus 4}

\lettrine[lines=2]{C}{UM} invocárem exaudívit me Deus justítiæ meæ: * in tribulatióne dilatásti mihi.

Miserére mei, * et exáudi oratiónem meam.

Fílii hóminum, úsquequo gravi corde? * ut quid dilígitis vanitátem, et quǽritis mendácium?

Et scitóte quóniam mirificávit Dóminus sanctum suum: * Dóminus exáudiet me cum clamávero ad eum.

Irascímini, et nolíte peccáre: * quæ dícitis in córdibus vestris, in cubílibus vestris compungímini.

Sacrificáte sacrifícium justítiæ, et speráte in Dómino. * Multi dicunt: Quis osténdit nobis bona?

Signátum est super nos lumen vultus tui, Dómine: * dedísti lætítiam in corde meo.

A fructu fruménti, vini, et ólei sui * multiplicáti sunt.

In pace in idípsum * dórmiam, et requiéscam;

Quóniam tu, Dómine, singuláriter in spe * constituísti me.

\psalmheading{Psalmus 90}

\lettrine[lines=2]{Q}{UI} hábitat in adjutório Altíssimi, * in protectióne Dei cæli commorábitur.

Dicet Dómino: Suscéptor meus es tu, et refúgium meum: * Deus meus, sperábo in eum.

Quóniam ipse liberávit me de láqueo venántium, * et a verbo áspero.

Scápulis suis obumbrábit tibi: * et sub pennis ejus sperábis.

Scuto circúmdabit te véritas ejus: * non timébis a timóre noctúrno.

A sagítta volánte in die, a negótio perambulánte in ténebris: * ab incúrsu, et dæmónio meridiáno.

Cadent a látere tuo mille, et decem míllia a dextris tuis: * ad te autem non appropinquábit.

Verúmtamen óculis tuis considerábis: * et retributiónem peccatórum vidébis.

Quóniam tu es, Dómine, spes mea: * Altíssimum posuísti refúgium tuum.

Non accédet ad te malum: * et flagéllum non appropinquábit tabernáculo tuo.

Quóniam Angelis suis mandávit de te: * ut custódiant te in ómnibus viis tuis.

In mánibus portábunt te: * ne forte offéndas ad lápidem pedem tuum.

Super áspidem, et basilíscum ambulábis: * et conculcábis leónem et dracónem.

Quóniam in me sperávit, liberábo eum: * prótegam eum, quóniam cognóvit nomen meum.

Clamábit ad me, et ego exáudiam eum: * cum ipso sum in tribulatióne: erípiam eum et glorificábo eum.

Longitúdine diérum replébo eum: * et osténdam illi salutáre meum.

\psalmheading{Psalmus 133}

\lettrine[lines=2]{E}{CCE}, nunc benedícite Dóminum, * omnes servi Dómini:

Qui statis in domo Dómini, * in átriis domus Dei nostri.

In nóctibus extóllite manus vestras in sancta, * et benedícite Dóminum.

Benedícat te Dóminus ex Sion, * qui fecit cælum et terram.

\antlabel{Ant.} Miserére mihi, Dómine, et exáudi oratiónem meam.

\hymnusheading{Hymnus}

\lettrine[lines=2]{T}{E} lucis ante términum,\\
Rerum Creátor, póscimus,\\
Ut pro tua cleméntia\\
Sis præsul et custódia.

Procul recédant sómnia,\\
Et nóctium phantásmata;\\
Hostémque nostrum cómprime,\\
Ne polluántur córpora.

* Præsta, Pater piíssime,\\
Patríque compar Unice,\\
Cum Spíritu Paráclito\\
Regnans per omne séculum.\\
Amen.

\capitulumlabel{Capitulum}{Jer. 14, 9}

\lettrine[lines=2]{T}{U} autem in nobis es, Dómine, et nomen sanctum tuum invocátum est super nos: ne derelínquas nos, Dómine, Deus noster.

\R\ br. In manus tuas, Dómine, * Comméndo spíritum meum. In manus. \V\ Redemísti nos, Dómine, Deus veritátis. Comméndo. Glória Patri. In manus.

\V\ Custódi nos, Dómine, ut pupíllam óculi. \R\ Sub umbra alárum tuárum prótege nos.

\antlabel{Ant.} Salva nos.

\canticumheading{Canticum Simeonis}{Luc. 2, 29-32}

\lettrine[lines=2]{N}{UNC} dimíttis servum tuum, Dómine, * secúndum verbum tuum in pace:

Quia vidérunt óculi mei * salutáre tuum,

Quod parásti * ante fáciem ómnium populórum,

Lumen ad revelatiónem géntium, * et glóriam plebis tuæ Israël.

Glória Patri.

\antlabel{Ant.} Salva nos, Dómine, vigilántes, custódi nos dormiéntes; ut vigilémus cum Christo, et requiescámus in pace.

\variantheading{Preces (si dicendæ sint)}

KYRIE, eléison. Christe, eléison. Kýrie, eléison.

Pater noster secreto usque ad \V\ Et ne nos indúcas in tentatiónem. \R\ Sed líbera nos a malo.

Credo in Deum secreto usque ad \V\ Carnis resurrectiónem. \R\ Vitam ætérnam. Amen.

\V\ Benedíctus es, Dómine, Deus patrum nostrórum. \R\ Et laudábilis et gloriósus in sécula. \V\ Benedicámus Patrem et Fílium cum Sancto Spíritu. \R\ Laudémus, et superexaltémus eum in sécula. \V\ Benedíctus es, Dómine, in firmaménto cæli. \R\ Et laudábilis, et gloriósus, et superexaltátus in sécula. \V\ Benedícat et custódiat nos omnípotens et miséricors Dóminus. \R\ Amen. \V\ Dignáre, Dómine, nocte ista. \R\ Sine peccáto nos custodíre. \V\ Miserére nostri, Dómine. \R\ Miserére nostri. \V\ Fiat misericórdia tua, Dómine, super nos. \R\ Quemádmodum sperávimus in te. \V\ Dómine, exáudi oratiónem meam. \R\ Et clamor meus ad te véniat.

\V\ Dóminus vobíscum.

Orémus. \oratiolabel{Oratio}

\lettrine[lines=2]{V}{ISITA}, quǽsumus, Dómine, habitatiónem istam, et omnes insídias inimíci ab ea longe repélle: Angeli tui sancti hábitent in ea, qui nos in pace custódiant; et benedíctio tua sit super nos semper. Per Dóminum nostrum. \R\ Amen.

\V\ Dóminus vobíscum. \V\ Benedicámus Dómino.

\benedictiolabel{Benedictio} Benedícat et custódiat nos omnípotens et miséricors Dóminus, Pater, et Fílius, et Spíritus Sanctus. \R\ Amen.

\antlabel{Antiphona finalis B. M. V.}

\V\ Divínum auxílium.

Pater, Ave et Credo.

\rubric{Psalmi de Dominica supérius appósiti dicúntur ad Completórium, adhíbita simul Antíphona, étiam in Vigília Epiphaníæ, post utrásque Vésperas in Domínica, quamvis translata, infra Octávam Nativitátis Dómini, in Duplícibus I vel II clássis, in Festis novem Lectiónum Dómini, B. Mariæ Vírginis, Angelórum, S. Joánnis Baptístæ, S. Joseph, Apostolórum et Evangelistárum, et per Octávas privilegiátas; si tamen præcedéntes I Vésperæ de eis íntegræ vel saltem a Capítulo sint celebrátæ, ac II Vésperæ de ipsis íntegræ fúerint persolútæ.}

\sectionornament

\feriaheading{Feria Secunda}

\hourheading{Ad Matutinum}

Pater, Ave et Credo.

\V\ Dómine, lábia mea.

\V\ Deus, in adjutórium.

\variantheading{In Officio feriali Adventus}

\variantheading{Infra Hebdomadam I et II:}

\invitatlabel{Invit.} Regem ventúrum Dóminum, * Veníte, adorémus.

\variantheading{Infra Hebdom. III et IV:}

\invitatlabel{Invit.} Prope est jam Dóminus: * Veníte, adorémus.

\variantheading{In Vigilia Nativ. Domini:}

\invitatlabel{Invit.} Hódie sciétis quia véniet Dóminus: * Et mane vidébitis glóriam ejus.

\hymnusheading{Hymnus}

\lettrine[lines=2]{V}{ERBUM} supérnum, pródiens\\
E Patris ætérni sinu,\\
Qui natus orbi súbvenis,\\
Labénte cursu témporis:

Illúmina nunc péctora,\\
Tuóque amóre cóncrema;\\
Ut cor cadúca déserens\\
Cæli volúptas ímpleat.

Ut, cum tribúnal Júdicis\\
Damnábit igni nóxios,\\
Et vox amíca débitum\\
Vocábit ad cælum pios;

Non esca flammárum nigros\\
Volvámur inter túrbines,\\
Vultu Dei sed cómpotes\\
Cæli fruámur gáudiis.

Patri simúlque Fílio,\\
Tibíque, Sancte Spíritus,\\
Sicut fuit, sit júgiter\\
Séclum per omne glória.\\
Amen.

\variantheading{Per Annum}

\rubric{Invitatórium et Hymnus dicúntur quotídie de Féria curránti, ut in Psaltério hic et alibi appónuntur.}

\invitatlabel{Invitat.} Veníte, * Exsultémus Dómino. \psalmlabel{Ps. 94} Veníte, exsultémus.

\rubric{In primo Psalmi versu non dícitur: Veníte, exsultémus Dómino; sed post repetítum Invitatórium, statim subjúngitur: Jubilémus Deo, salutári nostro:}

\hymnusheading{Hymnus}

\lettrine[lines=2]{S}{OMNO} reféctis ártubus,\\
Spreto cubíli, súrgimus:\\
Nobis, Pater, canéntibus\\
Adésse te depóscimus.

Te lingua primum cóncinat,\\
Te mentis ardor ámbiat:\\
Ut áctuum sequéntium\\
Tu, Sancte, sis exórdium.

Cedant ténebræ lúmini,\\
Et nox diúrno síderi,\\
Ut culpa, quam nox íntulit,\\
Lucis labáscat múnere.

Precámur iídem súpplices,\\
Noxas ut omnes ámputes,\\
Et ore te canéntium\\
Laudéris omni témpore.

* Præsta, Pater piíssime,\\
Patríque compar Unice,\\
Cum Spíritu Paráclito\\
Regnans per omne séculum.\\
Amen.

\nocturnoheading{In I Nocturno}

\antlabel{Ant.} Dóminus de cælo * prospéxit super fílios hóminum.

\rubric{Antíphonæ, in hac et áliis Fériis apposítæ, dicúntur semper in Officio feriáli et in Festis sive trium sive novem Lectiónum, quando juxta Rúbricas suméndi sint Psalmi de Féria, ut in Psaltério.}

\rubric{In Officio autem trium Lectiónum, tam de Témpore quam de Sanctis, omittúntur Versus et relíqua ómnia, quæ in Officio novem Lectiónum solent recitári in fine I et II Nocturni, et continuántur Psalmi cum suis Antíphonis usque ad Versum III Noctúrni.}

\psalmheading{Psalmus 13}

\lettrine[lines=2]{D}{IXIT} insípiens in corde suo: * Non est Deus.

Corrúpti sunt, et abominábiles facti sunt in stúdiis suis: * non est qui fáciat bonum, non est usque ad unum.

Dóminus de cælo prospéxit super fílios hóminum, * ut vídeat si est intélligens, aut requírens Deum.

Omnes declinavérunt, simul inútiles facti sunt: * non est qui fáciat bonum, non est usque ad unum.

Sepúlcrum patens est guttur eórum: linguis suis dolóse agébant, * venénum áspidum sub lábiis eórum.

Quorum os maledictióne et amaritúdine plenum est: * velóces pedes eórum ad effundéndum sánguinem.

Contrítio et infelícitas in viis eórum, et viam pacis non cognovérunt: * non est timor Dei ante óculos eórum.

Nonne cognóscent omnes qui operántur iniquitátem, * qui dévorant plebem meam sicut escam panis?

Dóminum non invocavérunt, * illic trepidavérunt timóre, ubi non erat timor.

Quóniam Dóminus in generatióne justa est, consílium ínopis confudístis: * quóniam Dóminus spes ejus est.

Quis dabit ex Sion salutáre Israël? * cum avérterit Dóminus captivitátem plebis suæ, exsultábit Jacob, et lætábitur Israël.

\antlabel{Ant.} Dóminus de cælo prospéxit super fílios hóminum.

\antlabel{Ant.} Qui operátur justítiam, * requiéscet in monte sancto tuo, Dómine.

\psalmheading{Psalmus 14}

\lettrine[lines=2]{D}{OMINE}, quis habitábit in tabernáculo tuo? * aut quis requiéscet in monte sancto tuo?

Qui ingréditur sine mácula, * et operátur justítiam:

Qui lóquitur veritátem in corde suo, * qui non egit dolum in lingua sua:

Nec fecit próximo suo malum, * et oppróbrium non accépit advérsus próximos suos.

Ad níhilum dedúctus est in conspéctu ejus malígnus: * timéntes autem Dóminum gloríficat:

Qui jurat próximo suo, et non décipit, * qui pecúniam suam non dedit ad usúram, et múnera super innocéntem non accépit.

Qui facit hæc, * non movébitur in ætérnum.

\antlabel{Ant.} Qui operátur justítiam, requiéscet in monte sancto tuo, Dómine.

\antlabel{Ant.} Inclína, Dómine, * aurem tuam mihi, et exáudi verba mea.

\psalmheading{Psalmus 16}

\lettrine[lines=2]{E}{XAUDI}, Dómine, justítiam meam: * inténde deprecatiónem meam.

Áuribus pércipe oratiónem meam, * non in lábiis dolósis.

De vultu tuo judícium meum pródeat: * óculi tui vídeant æquitátes.

Probásti cor meum, et visitásti nocte: * igne me examinásti, et non est invénta in me iníquitas.

Ut non loquátur os meum ópera hóminum: * propter verba labiórum tuórum ego custodívi vias duras.

Pérfice gressus meos in sémitis tuis: * ut non moveántur vestígia mea.

Ego clamávi, quóniam exaudísti me, Deus: * inclína aurem tuam mihi, et exáudi verba mea.

Mirífica misericórdias tuas, * qui salvos facis sperántes in te.

A resisténtibus déxteræ tuæ custódi me, * ut pupíllam óculi.

Sub umbra alárum tuárum prótege me: * a fácie impiórum qui me afflixérunt.

Inimíci mei ánimam meam circumdedérunt, ádipem suum conclusérunt: * os eórum locútum est supérbiam.

Projiciéntes me nunc circumdedérunt me: * óculos suos statuérunt declináre in terram.

Suscepérunt me sicut leo parátus ad prǽdam: * et sicut cátulus leónis hábitans in abdítis.

Exsúrge, Dómine, prǽveni eum, et supplánta eum: * éripe ánimam meam ab ímpio, fráméam tuam ab inimícis manus tuæ.

Dómine, a paucis de terra dívide eos in vita eórum: * de abscónditis tuis adimplétus est venter eórum.

Saturáti sunt fíliis: * et dimisérunt relíquias suas párvulis suis.

Ego autem in justítia apparébo conspéctui tuo: * satiábor cum apparúerit glória tua.

\antlabel{Ant.} Inclína, Dómine, aurem tuam mihi, et exáudi verba mea.

\variantheading{In Festis novem Lectionum:}

\V\ Mirífica, Dómine, misericórdias tuas. \R\ Qui salvos facis sperántes in te.

\absolutiolabel{Absol.} Exáudi, Dómine.

\benedlabel{Bened.} 1 Benedictióne perpétua. 2 Unigénitus Dei Fílius. 3 Spíritus Sancti grátia.

\nocturnoheading{In II Nocturno}

\antlabel{Ant.} Díligam te, \textdagger\ Dómine, virtus mea.

\psalmheading{Psalmus 17, i}

\lettrine[lines=2]{D}{ILIGAM} te, \textdagger\ Dómine, fortitúdo mea: * Dóminus firmaméntum meum, et refúgium meum, et liberátor meus.

Deus meus adjútor meus, * et sperábo in eum.

Protéctor meus, et cornu salútis meæ, * et suscéptor meus.

Laudans invocábo Dóminum: * et ab inimícis meis salvus ero.

Circumdedérunt me dolóres mortis: * et torréntes iniquitátis conturbavérunt me.

Dolóres inférni circumdedérunt me: * præoccupavérunt me láquei mortis.

In tribulatióne mea invocávi Dóminum, * et ad Deum meum clamávi.

Et exaudívit de templo sancto suo vocem meam: * et clamor meus in conspéctu ejus, introívit in aures ejus.

Commóta est, et contrémuit terra: * fundaménta móntium conturbáta sunt, et commóta sunt, quóniam irátus est eis.

Ascéndit fumus in ira ejus: * et ignis a fácie ejus exársit: * carbónes succénsi sunt ab eo.

Inclinávit cælos, et descéndit: * et calígo sub pédibus ejus.

Et ascéndit super Chérubim, et volávit: * volávit super pennas ventórum.

Et pósuit ténebras latíbulum suum, in circúitu ejus tabernáculum ejus: * tenebrósa aqua in núbibus áëris.

Præ fulgóre in conspéctu ejus nubes transiérunt, * grando et carbónes ignis.

Et intónuit de cælo Dóminus, et Altíssimus dedit vocem suam: * grando et carbónes ignis.

Et misit sagíttas suas, et dissipávit eos: * fúlgura multiplicávit, et conturbávit eos.

Et apparuérunt fontes aquárum, * et reveláta sunt fundaménta orbis terrárum:

Ab increpatióne tua, Dómine, * ab inspiratióne spíritus iræ tuæ.

\antlabel{Ant.} Díligam te, Dómine, virtus mea.

\antlabel{Ant.} Retríbuet mihi Dóminus * secúndum justítiam meam.

\psalmheading{Psalmus 17, ii}

\lettrine[lines=2]{M}{ISIT} de summo, et accépit me: * et assúmpsit me de aquis multis.

Erípuit me de inimícis meis fortíssimis, et ab his qui odérunt me: * quóniam confortáti sunt super me.

Prævenérunt me in die afflictiónis meæ: * et factus est Dóminus protéctor meus.

Et edúxit me in latitúdinem: * salvum me fecit, quóniam vóluit me.

Et retríbuet mihi Dóminus secúndum justítiam meam, * et secúndum puritátem mánuum meárum retríbuet mihi:

Quia custodívi vias Dómini, * nec ímpie gessi a Deo meo.

Quóniam ómnia judícia ejus in conspéctu meo: * et justítias ejus non répuli a me.

Et ero immaculátus cum eo: * et observábo me ab iniquitáte mea.

Et retríbuet mihi Dóminus secúndum justítiam meam: * et secúndum puritátem mánuum meárum in conspéctu oculórum ejus.

Cum sancto sanctus eris, * et cum viro innocénte ínnocens eris:

Et cum elécto eléctus eris: * et cum pervérso pervertéris.

Quóniam tu pópulum húmilem salvum fácies: * et óculos superbórum humiliábis.

Quóniam tu illúminas lucérnam meam, Dómine: * Deus meus, illúmina ténebras meas.

Quóniam in te erípiar a tentatióne, * et in Deo meo transgrédiar murum.

Deus meus, impollúta via ejus: elóquia Dómini igne examináta: * protéctor est ómnium sperántium in se.

Quóniam quis Deus præter Dóminum? * aut quis Deus præter Deum nostrum?

Deus, qui præcínxit me virtúte: * et pósuit immaculátam viam meam.

Qui perfécit pedes meos tamquam cervórum, * et super excélsa státuens me.

Qui docet manus meas ad prǽlium: * et posuísti, ut arcum áëreum, bráchia mea.

\antlabel{Ant.} Retríbuet mihi Dóminus secúndum justítiam meam.

\antlabel{Ant.} Vivit Dóminus, * et benedíctus Deus salútis meæ.

\psalmheading{Psalmus 17, iii}

\lettrine[lines=2]{E}{T} dedísti mihi protectiónem salútis tuæ: * et déxtera tua suscépit me:

Et disciplína tua corréxit me in finem: * et disciplína tua ipsa me docébit.

Dilatásti gressus meos subtus me: * et non sunt infirmáta vestígia mea:

Pérsequar inimícos meos, et comprehéndam illos: * et non convértar, donec defíciant.

Confríngam illos, nec póterunt stare: * cadent subtus pedes meos.

Et præcinxísti me virtúte ad bellum: * et supplantásti insurgéntes in me subtus me.

Et inimícos meos dedísti mihi dorsum, * et odiéntes me disperdidísti.

Clamavérunt, nec erat qui salvos fáceret, ad Dóminum: * nec exaudívit eos.

Et commínuam eos, ut púlverem ante fáciem venti: * ut lutum plateárum delébo eos.

Erípies me de contradictiónibus pópuli: * constítues me in caput géntium.

Pópulus, quem non cognóvi, servívit mihi: * in audítu auris obedívit mihi.

Fílii aliéni mentíti sunt mihi: * fílii aliéni inveteráti sunt, et claudicavérunt a sémitis suis.

Vivit Dóminus, et benedíctus Deus meus, * et exaltétur Deus salútis meæ.

Deus, qui das vindíctas mihi, et subdis pópulos sub me, * liberátor meus de inimícis meis iracúndis.

Et ab insurgéntibus in me exaltábis me: * a viro iníquo éripies me.

Proptérea confitébor tibi in natiónibus, Dómine: * et nómini tuo psalmum dicam.

Magníficans salútes Regis ejus, et fáciens misericórdiam Christo suo David, * et sémini ejus usque in séculum.

\antlabel{Ant.} Vivit Dóminus, et benedíctus Deus salútis meæ.

\variantheading{In Festis novem Lectionum:}

\V\ Confitébor tibi in natiónibus, Dómine. \R\ Et nómini tuo psalmum dicam.

\absolutiolabel{Absol.} Ipsíus píetas et.

\benedlabel{Bened.} 4 Deus Pater omnípotens. 5 Christus perpétuæ. 6 Ignem sui amóris.

\nocturnoheading{In III Nocturno}

\antlabel{Ant.} Exáudiat te * \textdagger\ Dóminus in die tribulatiónis. \textdagger

\psalmheading{Psalmus 19}

\lettrine[lines=2]{E}{XAUDIAT} te \textdagger\ Dóminus in die tribulatiónis: * \textdagger\ prótegat te nomen Dei Jacob.

Mittat tibi auxílium de sancto: * et de Sion tueátur te.

Memor sit omnis sacrifícii tui: * et holocáustum tuum pingue fiat.

Tríbuat tibi secúndum cor tuum: * et omne consílium tuum confírmet.

Lætábimur in salutári tuo: * et in nómine Dei nostri magnificábimur.

Ímpleat Dóminus omnes petitiónes tuas: * nunc cognóvi quóniam salvum fecit Dóminus Christum suum.

Exáudiet illum de cælo sancto suo: * in potentátibus salus déxteræ ejus.

Hi in cúrribus, et hi in equis: * nos autem in nómine Dómini, Dei nostri, invocábimus.

Ipsi obligáti sunt, et cecidérunt: * nos autem surréximus et erécti sumus.

Dómine, salvum fac regem: * et exáudi nos in die, qua invocavérimus te.

\antlabel{Ant.} Dómine, * \textdagger\ in virtúte tua lætábitur rex. \textdaggerdbl

\psalmheading{Psalmus 20}

\lettrine[lines=2]{D}{OMINE}, \textdagger\ in virtúte tua lætábitur rex: * \textdaggerdbl\ et super salutáre tuum exsultábit vehementer.

Desidérium cordis ejus tribuísti ei: * et voluntáte labiórum ejus non fraudásti eum.

Quóniam prævenísti eum in benedictiónibus dulcédinis: * posuísti in cápite ejus corónam de lápide pretióso.

Vitam pétiit a te: * et tribuísti ei longitúdinem diérum in séculum, et in séculum séculi.

Magna est glória ejus in salutári tuo: * glóriam et magnum decórem impónes super eum.

Quóniam dabis eum in benedictiónem in séculum séculi: * lætificábis eum in gáudio cum vultu tuo.

Quóniam rex sperat in Dómino: * et in misericórdia Altíssimi non commovébitur.

Inveniátur manus tua ómnibus inimícis tuis: * déxtera tua invéniat omnes, qui te odérunt.

Pones eos ut clíbanum ignis in témpore vultus tui: * Dóminus in ira sua conturbábit eos, et devorábit eos ignis.

Fructum eórum de terra perdes: * et semen eórum a fíliis hóminum.

Quóniam declinavérunt in te mala: * cogitavérunt consília, quæ non potuérunt stabilíre.

Quóniam pones eos dorsum: * in relíquiis tuis præparábis vultum eórum.

Exaltáre, Dómine, in virtúte tua: * cantábimus et psallémus virtútes tuas.

\antlabel{Ant.} Dómine, in virtúte tua lætábitur rex.

\antlabel{Ant.} Exaltábo te, * \textdagger\ Dómine, quóniam suscepísti me. \textdaggerdbl

\psalmheading{Psalmus 29}

\lettrine[lines=2]{E}{XALTABO} te, \textdagger\ Dómine, quóniam suscepísti me: * \textdaggerdbl\ nec delectásti inimícos meos super me.

Dómine, Deus meus, clamávi ad te, * et sanásti me.

Dómine, eduxísti ab inférno ánimam meam: * salvásti me a descendéntibus in lacum.

Psállite Dómino, sancti ejus: * et confitémini memóriæ sanctitátis ejus.

Quóniam ira in indignatióne ejus: * et vita in voluntáte ejus.

Ad vésperum demorábitur fletus: * et ad matutínum lætítia.

Ego autem dixi in abundántia mea: * Non movébor in ætérnum.

Dómine, in voluntáte tua, * præstitísti decóri meo virtútem.

Avertísti fáciem tuam a me, * et factus sum conturbátus.

Ad te, Dómine, clamábo: * et ad Deum meum deprecábor.

Quæ utílitas in sánguine meo, * dum descéndo in corruptiónem?

Numquid confitébitur tibi pulvis, * aut annuntiábit veritátem tuam?

Audívit Dóminus, et misértus est mei: * Dóminus factus est adjútor meus.

Convertísti planctum meum in gáudium mihi: * conscidísti saccum meum, et circumdedísti me lætítia:

Ut cantet tibi glória mea, et non compúngar: * Dómine, Deus meus, in ætérnum confitébor tibi.

\antlabel{Ant.} Exaltábo te, Dómine, quóniam suscepísti me.

\variantheading{In feriali Officio per Annum et in Festis:}

\V\ Psállite Dómino, sancti ejus. \R\ Et confitémini memóriæ sanctitátis ejus.

\variantheading{In feriali Officio Adventus:}

\V\ Ex Sion spécies decóris ejus. \R\ Deus noster maniféste véniet.

\variantheading{In Vigilia Nativitatis Dom.:}

\V\ Hódie sciétis quia véniet Dóminus. \R\ Et mane vidébitis glóriam ejus.

\variantheading{In Officio novem Lectionum}

\absolutiolabel{Absolutio} A vínculis.

\benedlabel{Bened.} 7 Evangélica léctio. 8 In Festis Domini: Divínum auxílium. In Festis B. M. V.: Cujus festum cólimus, ipsa Virgo. In Festis Sanctorum: Cujus (\textit{vel} Quorum \textit{aut} Quarum) festum cólimus. 9 Ad societátem. (Per evangélica dicta.)

\variantheading{In Officio trium Lectionum}

\absolutiolabel{Absol.} Exáudi, Dómine.

\variantheading{In Feriis, cum Homilia:}

\benedlabel{Bened.} 1 Evangélica. 2 Divínum. 3 Ad societátem.

\variantheading{In Feriis, sine Homilia:}

\benedlabel{Bened.} 1 Benedictióne. 2 Unigénitus. 3 Spíritus Sancti.

\variantheading{In Festis Sanctorum:}

\benedlabel{Bened.} 1 Ille nos benedícat qui sine fine vivit et regnat. \R\ Amen. 2 Cujus (\textit{vel} Quorum \textit{aut} Quarum). 3 Ad societátem.

\sectionornament

\hourheading{Ad Laudes I}

\textit{In Feriis per Annum, extra hebdomádas Septuagésimæ, Sexagésimæ et Quinquagésimæ ac Vigílias commúnes occurréntibus, atque in Festis quólibet anni témpore celebrándis.}

(Pater noster et Ave.)

\V\ Deus, in adjutórium.

\antlabel{Ant.} Jubiláte * Deo in voce exsultatiónis.

\rubric{Antíphonæ, in hac et áliis Fériis apposítæ, dicúntur semper in Officio feriáli et in Festis sive trium sive novem Lectiónum, quando juxta Rubricas sumendi sint Psalmi de Feria; atque adhibéntur étiam in Festis, quæ Témpore Advéntus celebrántur.}

\psalmheading{Psalmus 46}

\lettrine[lines=2]{O}{MNES} gentes, plaúdite mánibus: * jubiláte Deo in voce exsultatiónis.

Quóniam Dóminus excélsus, terríbilis: * Rex magnus super omnem terram.

Subjécit pópulos nobis: * et gentes sub pédibus nostris.

Elégit nobis hereditátem suam: * spéciem Jacob, quam diléxit.

Ascéndit Deus in júbilo: * et Dóminus in voce tubæ.

Psállite Deo nostro, psállite: * psállite Regi nostro, psállite.

Quóniam Rex omnis terræ Deus: * psállite sapiénter.

Regnábit Deus super gentes: * Deus sedet super sedem sanctam suam.

Príncipes populórum congregáti sunt cum Deo Abraham: * quóniam dii fortes terræ veheménter eleváti sunt.

\antlabel{Ant.} Jubiláte Deo in voce exsultatiónis.

\antlabel{Ant.} Inténde * voci oratiónis meæ, Rex meus et Deus meus.

\psalmheading{Psalmus 5}

\lettrine[lines=2]{V}{ERBA} mea áuribus pércipe, Dómine, * intéllige clamórem meum.

Inténde voci oratiónis meæ, * Rex meus et Deus meus.

Quóniam ad te orábo: * Dómine, mane exáudies vocem meam.

Mane astábo tibi et vidébo: * quóniam non Deus volens iniquitátem tu es.

Neque habitábit juxta te malígnus: * neque permanébunt injústi ante óculos tuos.

Odísti omnes, qui operántur iniquitátem: * perdes omnes, qui loquúntur mendácium.

Virum sánguinum et dolósum abominábitur Dóminus: * ego autem in multitúdine misericórdiæ tuæ.

Introíbo in domum tuam: * adorábo ad templum sanctum tuum in timóre tuo.

Dómine, deduc me in justítia tua: * propter inimícos meos dírige in conspéctu tuo viam meam.

Quóniam non est in ore eórum véritas: * cor eórum vanum est.

Sepúlcrum patens est guttur eórum, linguis suis dolóse agébant, * judíca illos, Deus.

Décidant a cogitatiónibus suis, secúndum multitúdinem impietátum eórum expélle eos, * quóniam irritavérunt te, Dómine.

Et læténtur omnes, qui sperant in te, * in ætérnum exsultábunt: et habitábis in eis.

Et gloriabúntur in te omnes, qui díligunt nomen tuum, * quóniam tu benedíces justo.

Dómine, ut scuto bonæ voluntátis tuæ * coronásti nos.

\antlabel{Ant.} Inténde voci oratiónis meæ, Rex meus et Deus meus.

\antlabel{Ant.} Deus majestátis * intónuit: afférte glóriam nómini ejus.

\psalmheading{Psalmus 28}

\lettrine[lines=2]{A}{FFERTE} Dómino, fílii Dei: * afférte Dómino fílios aríetum.

Afférte Dómino glóriam et honórem, afférte Dómino glóriam nómini ejus: * adoráte Dóminum in átrio sancto ejus.

Vox Dómini super aquas, Deus majestátis intónuit: * Dóminus super aquas multas.

Vox Dómini in virtúte: * vox Dómini in magnificéntia.

Vox Dómini confringéntis cedros: * et confrínget Dóminus cedros Líbani:

Et commínuet eas tamquam vítulum Líbani: * et diléctus quemádmodum fílius unicórnium.

Vox Dómini intercidéntis flammam ignis: * vox Dómini concutiéntis desértum: et commovébit Dóminus desértum Cades.

Vox Dómini præparántis cervos, et revelábit condénsa: * et in templo ejus omnes dicent glóriam.

Dóminus dilúvium inhabitáre facit: * et sedébit Dóminus Rex in ætérnum.

Dóminus virtútem pópulo suo dabit: * Dóminus benedícet pópulo suo in pace.

\antlabel{Ant.} Deus majestátis intónuit: afférte glóriam nómini ejus.

\antlabel{Ant.} Laudámus nomen tuum * ínclitum, Deus noster.

\canticumheading{Canticum David}{1 Par. 29, 10-13}

\lettrine[lines=2]{B}{ENEDICTUS} es, Dómine, Deus Israël patris nostri, * ab ætérno in ætérnum.

Tua est, Dómine, magnificéntia, et poténtia, * et glória, atque victória:

Et tibi laus: * cuncta enim quæ in cælo sunt, et in terra, tua sunt:

Tuum, Dómine, regnum, * et tu es super omnes príncipes.

Tuæ divítiæ, et tua est glória: * tu domináris ómnium.

In manu tua virtus et poténtia: * in manu tua magnitúdo, et impérium ómnium.

Nunc ígitur, Deus noster, confitémur tibi, * et laudámus nomen tuum ínclitum.

\antlabel{Ant.} Laudámus nomen tuum ínclitum, Deus noster.

\antlabel{Ant.} Laudáte \textdagger\ Dóminum, omnes gentes. \textdaggerdbl

\psalmheading{Psalmus 116}

\lettrine[lines=2]{L}{AUDATE} \textdagger\ Dóminum, omnes gentes: *\textdaggerdbl\ laudáte eum, omnes pópuli:

Quóniam confirmáta est super nos misericórdia ejus: * et véritas Dómini manet in ætérnum.

\antlabel{Ant.} Laudáte Dóminum, omnes gentes.

\variantheading{In Officio feriali}

\rubric{Per Annum: Capitulum, Hymnus et Versus dicuntur quotidie de Feria currenti, ut in Psalterio hic et alibi apponuntur; et inde sumitur etiam Ant. ad Benedictus, præter quam in Dominica post Octavam Epiphaniæ anticipata, in qua dicitur Antiphona, ut in Proprio de Tempore.}

\capitulumlabel{Capit.}{Rom. 13, 12-13}

\lettrine[lines=2]{N}{OX} præcéssit, dies autem appropinquávit. Abjiciámus ergo ópera tenebrárum, et induámur arma lucis. Sicut in die honéste ambulémus.

\hymnusheading{Hymnus}

\lettrine[lines=2]{S}{PLENDOR} Patérnæ glóriæ,\\
De luce lucem próferens,\\
Lux lucis, et fons lúminis,\\
Diem dies illúminans:

Verúsque sol illábere,\\
Micans nitóre pérpeti:\\
Jubárque Sancti Spíritus\\
Infúnde nostris sénsibus.

Votis vocémus et Patrem,\\
Patrem poténtis grátiæ,\\
Patrem perénnis glóriæ:\\
Culpam reléget lúbricam.

Confírmet actus strénuos:\\
Dentes retúndat ínvidi:\\
Casus secúndet ásperos:\\
Agénda recte dírigat.

Mentem gubérnet et regat:\\
Sit pura nobis cástitas:\\
Fides calóre férveat,\\
Fraudis venéna nésciat.

Christúsque nobis sit cibus,\\
Potúsque noster sit fides:\\
Læti bibámus sóbriam\\
Profusiónem Spíritus.

Lætus dies hic tránseat:\\
Pudor sit ut dilúculum:\\
Fides velut merídies:\\
Crepúsculum mens nésciat.

Auróra lucem próvehit,\\
Cum luce nobis pródeat\\
In Patre totus Fílius,\\
Et totus in Verbo Pater.

Deo Patri sit glória,\\
Ejúsque soli Fílio,\\
Cum Spíritu Paráclito,\\
Nunc et per omne séculum. Amen.

\V\ Repléti sumus mane misericórdia tua. \R\ Exsultávimus, et delectáti sumus.

\antlabel{Ad Bened. Ant.} Benedíctus * \textdagger\ Dóminus, Deus Israël, quia visitávit et liberávit nos.

\rubric{Reliqua ut in Ordinario. 11}

\sectionornament

\hourheading{Ad Laudes II}

\textit{In Feriis Adventus, et in Feriis II post Dominicas Septuagesimæ, Sexagesimæ et Quinquagesimæ, necnon in Vigiliis communibus, quando fit Officium de Feria.}

(Pater noster et Ave.)

\V\ Deus, in adjutórium.

\antlabel{Ant.} Miserére. \textdagger

\rubric{Fer. II ante Vig. Nat. Dom.:} \antlabel{Ant.} Ecce véniet.

\psalmheading{Psalmus 50}

\lettrine[lines=2]{M}{ISERERE} \textdagger\ mei, Deus: * secúndum magnam misericórdiam tuam.

Et secúndum multitúdinem miseratiónum tuárum, * dele iniquitátem meam.

Amplius lava me ab iniquitáte mea: * et a peccáto meo munda me.

Quóniam iniquitátem meam ego cognósco: * et peccátum meum contra me est semper.

Tibi soli peccávi, et malum coram te feci: * ut justificéris in sermónibus tuis, et vincas cum judicáris.

Ecce enim, in iniquitátibus concéptus sum: * et in peccátis concépit me mater mea.

Ecce enim, veritátem dilexísti: * incérta et occúlta sapiéntiæ tuæ manifestásti mihi.

Aspérges me hyssópo, et mundábor: * lavábis me, et super nivem dealbábor.

Audítui meo dabis gáudium et lætítiam: * et exsultábunt ossa humiliáta.

Avérte fáciem tuam a peccátis meis: * et omnes iniquitátes meas dele.

Cor mundum crea in me, Deus: * et spíritum rectum ínnova in viscéribus meis.

Ne projícias me a fácie tua: * et spíritum sanctum tuum ne áuferas a me.

Redde mihi lætítiam salutáris tui: * et spíritu principáli confírma me.

Docébo iníquos vias tuas: * et ímpii ad te converténtur.

Líbera me de sanguínibus, Deus, Deus salútis meæ: * et exsultábit lingua mea justítiam tuam.

Dómine, lábia mea apéries: * et os meum annuntiábit laudem tuam.

Quóniam si voluísses sacrifícium, dedíssem útique: * holocáustis non delectáberis.

Sacrifícium Deo spíritus contribulátus: * cor contrítum, et humiliátum, Deus, non despícies.

Benígne fac, Dómine, in bona voluntáte tua Sion: * ut ædificéntur muri Jerúsalem.

Tunc acceptábis sacrifícium justítiæ, oblatiónes, et holocáusta: * tunc impónent super altáre tuum vítulos.

\antlabel{Ant.} Miserére mei, Deus, secúndum magnam misericórdiam tuam.

\antlabel{Ant.} Deduc me.

\rubric{Fer. II ante Vig. Nat. Dom.:} \antlabel{Ant.} Ecce véniet Dóminus princeps regum terræ: beáti, qui paráti sunt occúrrere illi.

\antlabel{Ant.} Cum vénerit.

\psalmheading{Psalmus 5}

\lettrine[lines=2]{V}{ERBA} mea áuribus pércipe, Dómine, * intéllige clamórem meum.

Inténde voci oratiónis meæ, * Rex meus et Deus meus.

Quóniam ad te orábo: * Dómine, mane exáudies vocem meam.

Mane astábo tibi et vidébo: * quóniam non Deus volens iniquitátem tu es.

Neque habitábit juxta te malígnus: * neque permanébunt injústi ante óculos tuos.

Odísti omnes, qui operántur iniquitátem: * perdes omnes, qui loquúntur mendácium.

Virum sánguinum et dolósum abominábitur Dóminus: * ego autem in multitúdine misericórdiæ tuæ.

Introíbo in domum tuam: * adorábo ad templum sanctum tuum in timóre tuo.

Dómine, deduc me in justítia tua: * propter inimícos meos dírige in conspéctu tuo viam meam.

Quóniam non est in ore eórum véritas: * cor eórum vanum est.

Sepúlcrum patens est guttur eórum, linguis suis dolóse agébant, * júdica illos, Deus.

Décidant a cogitatiónibus suis, secúndum multitúdinem impietátum eórum expélle eos, * quóniam irritavérunt te, Dómine.

Et læténtur omnes, qui sperant in te, * in ætérnum exsultábunt: et habitábis in eis.

Et gloriabúntur in te omnes, qui díligunt nomen tuum, * quóniam tu benedíces justo.

Dómine, ut scuto bonæ voluntátis tuæ * coronásti nos.

\antlabel{Ant.} Deduc me in justítia tua, Dómine.

\antlabel{Ant.} Dóminus dabit virtútem.

\rubric{Fer. II ante Vig. Nat. Dom.:} \antlabel{Ant.} Cum vénerit Fílius hóminis, putas invéniet fidem super terram?

\antlabel{Ant.} Ecce jam venit.

\psalmheading{Psalmus 28}

\lettrine[lines=2]{A}{FFERTE} Dómino, fílii Dei: * afférte Dómino fílios aríetum.

Afférte Dómino glóriam et honórem, afférte Dómino glóriam nómini ejus: * adoráte Dóminum in átrio sancto ejus.

Vox Dómini super aquas, Deus majestátis intónuit: * Dóminus super aquas multas.

Vox Dómini in virtúte: * vox Dómini in magnificéntia.

Vox Dómini confringéntis cedros: * et confrínget Dóminus cedros Líbani:

Et commínuet eas tamquam vítulum Líbani: * et diléctus quemádmodum fílius unicórnium.

Vox Dómini intercidéntis flammam ignis: * vox Dómini concutiéntis desértum: et commovébit Dóminus desértum Cades.

Vox Dómini præparántis cervos, et revelábit condénsa: * et in templo ejus omnes dicent glóriam.

Dóminus dilúvium inhabitáre facit: * et sedébit Dóminus Rex in ætérnum.

Dóminus virtútem pópulo suo dabit: * Dóminus benedícet pópulo suo in pace.

\antlabel{Ant.} Dóminus dabit virtútem et benedícet pópulo suo in pace.

\antlabel{Ant.} Convérsus est furor tuus.

\rubric{Fer. II ante Vig. Nat. Dom.:} \antlabel{Ant.} Ecce jam venit plenitúdo témporis, in quo misit Deus Fílium suum in terras.

\antlabel{Ant.} Hauriétis aquas.

\canticumheading{Canticum Isaiæ}{Is. 12, 1-6}

\lettrine[lines=2]{C}{ONFITEBOR} tibi, Dómine, quóniam irátus es mihi: * convérsus est furor tuus, et consolátus es me.

Ecce Deus salvátor meus, * fiduciáliter agam, et non timébo:

Quia fortitúdo mea, et laus mea Dóminus, * et factus est mihi in salútem.

Hauriétis aquas in gáudio de fóntibus Salvatóris: * et dicétis in die illa: Confitémini Dómino, et invocáte nomen ejus:

Notas fácite in pópulis adinventiónes ejus: * mementóte quóniam excélsum est nomen ejus.

Cantáte Dómino quóniam magnífice fecit: * annuntiáte hoc in univérsa terra.

Exsúlta, et lauda, habitátio Sion: * quia magnus in médio tui Sanctus Israël.

\antlabel{Ant.} Convérsus est furor tuus, Dómine, et consolátus es me.

\antlabel{Ant.} Laudáte. \textdagger

\rubric{Fer. II ante Vig. Nat. Dom.:} \antlabel{Ant.} Hauriétis aquas in gáudio de fóntibus Salvatóris.

\antlabel{Ant.} Egrediétur.

\psalmheading{Psalmus 116}

\lettrine[lines=2]{L}{AUDATE} \textdagger\ Dóminum, omnes gentes: *\textdaggerdbl\ laudáte eum, omnes pópuli:

Quóniam confirmáta est super nos misericórdia ejus: * et véritas Dómini manet in ætérnum.

\antlabel{Ant.} Laudáte Dóminum, quóniam confirmáta est super nos misericórdia ejus.

\rubric{Fer. II ante Vig. Nat. Dom.:} \antlabel{Ant.} Egrediétur Dóminus de loco sancto suo: véniet ut salvet pópulum suum.

\variantheading{In Adventu}

\capitulumlabel{Capitulum}{Is. 2, 3}

\lettrine[lines=2]{V}{ENITE}, et ascendámus ad montem Dómini, et ad domum Dei Jacob, et docébit nos vias suas, et ambulábimus in sémitis ejus: quia de Sion exíbit lex, et verbum Dómini de Jerúsalem.

\hymnusheading{Hymnus}

\lettrine[lines=2]{E}{N}, clara vox redárguit\\
Obscúra quæque, pérsonans:\\
Procul fugéntur sómnia:\\
Ab alto Jesus prómicat.

Mens jam resúrgat, tórpida\\
Non ámplius jacens humi:\\
Sidus refúlget jam novum,\\
Ut tollat omne nóxium.

En, Agnus ad nos míttitur\\
Laxáre gratis débitum:\\
Omnes simul cum lácrimis\\
Precémur indulgéntiam:

Ut, cum secúndo fúlserit\\
Metúque mundum cínxerit,\\
Non pro reátu púniat,\\
Sed nos pius tunc prótegat.

Virtus, honor, laus, glória\\
Deo Patri cum Fílio,\\
Sancto simul Paráclito,\\
In sæculórum sæcula.\\
Amen.

\V\ Vox clamántis in desérto: Paráte viam Dómini. \R\ Rectas fácite sémitas ejus.

\rubric{Antiphona ad Benedictus ut in Proprio de Tempore.}

\rubric{Expleto Cantico ac repetita Antiphona, dicuntur sequentes}

\variantheading{Preces feriales}
\label{feria2:precesferiales}

KYRIE, eléison. Christe, eléison. Kýrie, eléison.

\rubric{Pater noster, quod dicitur a solo Hebdomadario totum clara voce, usque ad} \V\ Et ne nos indúcas in tentatiónem. \R\ Sed líbera nos a malo. \V\ Ego dixi: Dómine, miserére mei. \R\ Sana ánimam meam, quia peccávi tibi. \V\ Convértere, Dómine, úsquequo? \R\ Et deprecábilis esto super servos tuos. \V\ Fiat misericórdia tua, Dómine, super nos. \R\ Quemádmodum sperávimus in te. \V\ Sacerdótes tui induántur justítiam. \R\ Et sancti tui exsúltent.

\V\ Orémus pro beatíssimo Papa nostro N. \R\ Dóminus consérvet eum, et vivíficet eum, et beátum fáciat eum in terra, et non tradat eum in ánimam inimicórum ejus.

\V\ Orémus et pro Antístite nostro N. \R\ Stet et pascat in fortitúdine tua, Dómine, in sublimitáte nóminis tui.

\V\ Dómine, salvum fac regem. \R\ Et exáudi nos in die, qua invocavérimus te. \V\ Salvum fac pópulum tuum, Dómine, et bénedic hereditáti tuæ. \R\ Et rege eos, et extólle illos usque in ætérnum. \V\ Meménto Congregatiónis tuæ. \R\ Quam possedísti ab inítio. \V\ Fiat pax in virtúte tua. \R\ Et abundántia in túrribus tuis. \V\ Orémus pro benefactóribus nostris. \R\ Retribúere dignáre, Dómine, ómnibus, nobis bona faciéntibus propter nomen tuum, vitam ætérnam. Amen. \V\ Orémus pro fidélibus defúnctis. \R\ Réquiem ætérnam dona eis, Dómine, et lux perpétua lúceat eis. \V\ Requiéscant in pace. \R\ Amen. \V\ Pro frátribus nostris abséntibus. \R\ Salvos fac servos tuos, Deus meus, sperántes in te. \V\ Pro afflíctis et captívis. \R\ Líbera eos, Deus Israël, ex ómnibus tribulatiónibus suis. \V\ Mitte eis, Dómine, auxílium de sancto. \R\ Et de Sion tuére eos. \V\ Dómine, Deus virtútum, convérte nos. \R\ Et osténde fáciem tuam, et salvi érimus. \V\ Exsúrge, Christe, ádjuva nos. \R\ Et líbera nos propter nomen tuum. \V\ Dómine, exáudi oratiónem meam. \R\ Et clamor meus ad te véniat.

\variantheading{Per Annum}

\capitulumlabel{Capit.}{Rom. 13, 12-13}

\lettrine[lines=2]{N}{OX} præcéssit, dies autem appropinquávit. Abjiciámus ergo ópera tenebrárum, et induámur arma lucis. Sicut in die honéste ambulémus.

\hymnusheading{Hymnus}

\lettrine[lines=2]{S}{PLENDOR} Patérnæ glóriæ,\\
De luce lucem próferens,\\
Lux lucis, et fons lúminis,\\
Diem dies illúminans:

Verúsque sol illábere,\\
Micans nitóre pérpeti:\\
Jubárque Sancti Spíritus\\
Infúnde nostris sénsibus.

Votis vocémus et Patrem,\\
Patrem poténtis grátiæ,\\
Patrem perénnis glóriæ:\\
Culpam reléget lúbricam.

Confírmet actus strénuos:\\
Dentes retúndat ínvidi:\\
Casus secúndet ásperos:\\
Agénda recte dírigat.

Mentem gubérnet et regat:\\
Sit pura nobis cástitas:\\
Fides calóre férveat,\\
Fraudis venéna nésciat.

Christúsque nobis sit cibus,\\
Potúsque noster sit fides:\\
Læti bibámus sóbriam\\
Profusiónem Spíritus.

Lætus dies hic tránseat:\\
Pudor sit ut dilúculum:\\
Fides velut merídies:\\
Crepúsculum mens nésciat.

Auróra lucem próvehit,\\
Cum luce nobis pródeat\\
In Patre totus Fílius,\\
Et totus in Verbo Pater.

* Deo Patri sit glória,\\
Ejúsque soli Fílio,\\
Cum Spíritu Paráclito,\\
Nunc et per omne séculum.\\
Amen.

\V\ Repléti sumus mane misericórdia tua. \R\ Exsultávimus, et delectáti sumus.

\antlabel{Ad Bened. Ant.} Benedíctus * \textdagger\ Dóminus, Deus Israël, quia visitávit et liberávit nos.

\rubric{In Vigiliis communibus Preces feriales, ut supra.} (p.~\pageref{feria2:precesferiales})

\rubric{Reliqua ut in Ordinario. 11}

\sectionornament

\hourheading{Ad Primam}

Pater, Ave et Credo.

\V\ Deus, in adjutórium.

\hymnusheading{Hymnus}

\lettrine[lines=2]{J}{AM} lucis orto sídere,\\
Deum precémur súpplices,\\
Ut in diúrnis áctibus\\
Nos servet a nocéntibus.

Linguam refrénans témperet,\\
Ne litis horror ínsonet:\\
Visum fovéndo cóntegat,\\
Ne vanitátes háuriat.

Sint pura cordis íntima,\\
Absístat et vecórdia:\\
Carnis terat supérbiam\\
Potus cibíque párcitas.

Ut, cum dies abscésserit\\
Noctémque sors redúxerit,\\
Mundi per abstinéntiam\\
Ipsi canámus glóriam.

* Deo Patri sit glória,\\
Ejúsque soli Fílio,\\
Cum Spíritu Paráclito,\\
Nunc et per omne séculum.\\
Amen.

\variantheading{Per Annum} \antlabel{Ant.} Innocens mánibus.

\rubric{¶ Antiphonæ, in hac et aliis Feriis ad Primam, Tertiam, Sextam et Nonam appositæ, dicuntur semper in Officio feriali et in Festis sive trium sive novem Lectionum, quando juxta Rubricas sumendi sint Psalmi de Feria; atque adhibéntur étiam in Festis, quæ Témpore Advéntus celebrántur.}

\variantheading{In Officio feriali Adventus}

\rubric{Feria II post Dominicam I:} \antlabel{Ant.} In illa die.

\rubric{Feria II post Dominicam II:} \antlabel{Ant.} Ecce in núbibus cæli.

\rubric{Fer. II ante Vig. Nat. Dom.:} \antlabel{Ant.} Ecce véniet.

\psalmheading{Psalmus 23}

\lettrine[lines=2]{D}{OMINI} est terra, et plenitúdo ejus: * orbis terrárum, et univérsi qui hábitant in eo.

Quia ipse super mária fundávit eum: * et super flúmina præparávit eum.

Quis ascéndet in montem Dómini? * aut quis stabit in loco sancto ejus?

Innocens mánibus et mundo corde, * qui non accépit in vano ánimam suam, nec jurávit in dolo próximo suo.

Hic accípiet benedictiónem a Dómino: * et misericórdiam a Deo, salutári suo.

Hæc est generátio quæréntium eum, * quæréntium fáciem Dei Jacob.

Attóllite portas, príncipes, vestras, et elevámini, portæ æternáles: * et introíbit Rex glóriæ.

Quis est iste Rex glóriæ? * Dóminus fortis et potens: Dóminus potens in prœlio.

Attóllite portas, príncipes, vestras, et elevámini, portæ æternáles: * et introíbit Rex glóriæ.

Quis est iste Rex glóriæ? * Dóminus virtútum ipse est Rex glóriæ.

\psalmheading{Psalmus 18, i}

\lettrine[lines=2]{C}{ÆLI} enárrant glóriam Dei, * et ópera mánuum ejus annúntiat firmaméntum.

Dies diéi erúctat verbum, * et nox nocti índicat sciéntiam.

Non sunt loquélæ, neque sermónes, * quorum non audiántur voces eórum.

In omnem terram exívit sonus eórum: * et in fines orbis terræ verba eórum.

In sole pósuit tabernáculum suum: * et ipse tamquam sponsus procédens de thálamo suo:

Exsultávit ut gigas ad curréndam viam, * a summo cælo egréssio ejus:

Et occúrsus ejus usque ad summum ejus: * nec est qui se abscóndat a calóre ejus.

\psalmheading{Psalmus 18, ii}

\lettrine[lines=2]{L}{EX} Dómini immaculáta, convértens ánimas: * testimónium Dómini fidéle, sapiéntiam præstans párvulis.

Justítiæ Dómini rectæ, lætificántes corda: * præcéptum Dómini lúcidum, illúminans óculos.

Timor Dómini sanctus, pérmanens in sæculum sæculi: * judícia Dómini vera, justificáta in semetípsa.

Desiderabília super aurum et lápidem pretiósum multum: * et dulcióra super mel et favum.

Étenim servus tuus custódit ea: * in custodiéndis illis retribútio multa.

Delícta quis intélligit? ab occúltis meis munda me: * et ab aliénis parce servo tuo.

Si mei non fúerint domináti, tunc immaculátus ero: * et emundábor a delícto máximo.

Et erunt ut compláceant elóquia oris mei: * et meditátio cordis mei in conspéctu tuo semper.

Dómine, adjútor meus, * et redémptor meus.

\rubric{Quando ad Laudes dictus fuerit Ps. 50 Miserére, hic subjungitur sequens Ps. 46, qui in Laudibus locum non habuit.}

\psalmheading{Psalmus 46}

\lettrine[lines=2]{O}{MNES} gentes, plaúdite mánibus: * jubiláte Deo in voce exsultatiónis.

Quóniam Dóminus excélsus, terríbilis: * Rex magnus super omnem terram.

Subjécit pópulos nobis: * et gentes sub pédibus nostris.

Elégit nobis hereditátem suam: * spéciem Jacob, quam diléxit.

Ascéndit Deus in júbilo: * et Dóminus in voce tubæ.

Psállite Deo nostro, psállite: * psállite Regi nostro, psállite.

Quóniam Rex omnis terræ Deus: * psállite sapiénter.

Regnábit Deus super gentes: * Deus sedet super sedem sanctam suam.

Príncipes populórum congregáti sunt cum Deo Abraham: * quóniam dii fortes terræ veheménter eleváti sunt.

\rubric{Secus vero, tribus Psalmis expletis, statim dicitur Ant.}

\variantheading{Per Annum} \antlabel{Ant.} Innocens mánibus et mundo corde ascéndet in montem Dómini.

\variantheading{In Officio feriali Adventus}

\rubric{Feria II post Dominicam I:} \antlabel{Ant.} In illa die stillábunt montes dulcédinem, et colles fluent lac et mel, allelúja.

\rubric{Feria II post Dominicam II:} \antlabel{Ant.} Ecce in núbibus cæli Dóminus véniet cum potestáte magna, allelúja.

\rubric{Fer. II ante Vig. Nat. Dom.:} \antlabel{Ant.} Ecce véniet Dóminus princeps regum terræ: beáti, qui paráti sunt occúrrere illi.

\rubric{In Officio Festi vel Octavæ:}

\capitulumlabel{Capitulum}{1 Tim. 1, 17}

\lettrine[lines=2]{R}{EGI} sæculórum immortáli et invisíbili, soli Deo honor et glória in sæcula sæculórum. Amen.

\rubric{In omnibus autem Feriis et in Vigiliis communibus:}

\capitulumlabel{Capitulum}{Zach. 8, 19}

\lettrine[lines=2]{P}{ACEM} et veritátem dilígite, ait Dóminus omnípotens.

\R\ br. Christe, Fili Dei vivi, * Miserére nobis. Christe. \V\ \textit{Qui sedes ad déxteram Patris.} Miserére nobis. Glória Patri. Christe. \V\ Exsúrge, Christe, ádjuva nos. \R\ Et líbera nos propter nomen tuum.

\rubric{In Adventu, nisi alius in Festis fuerit assumendus:} \V\ Qui ventúrus es in mundum.

\variantheading{Preces dominicales et feriales}

\textit{(si dicendæ sint).}

KYRIE, eléison. Christe, eléison. Kýrie, eléison.

Pater noster \rubric{secreto usque ad} \V\ Et ne nos indúcas in tentatiónem. \R\ Sed líbera nos a malo.

Credo in Deum \rubric{secreto usque ad} \V\ Carnis resurrectiónem. \R\ Vitam ætérnam. Amen. \V\ Et ego ad te, Dómine, clamávi. \R\ Et mane orátio mea prævéniet te. \V\ Repleátur os meum laude. \R\ Ut cantem glóriam tuam, tota die magnitúdinem tuam. \V\ Dómine, avérte fáciem tuam a peccátis meis. \R\ Et omnes iniquitátes meas dele. \V\ Cor mundum crea in me, Deus. \R\ Et spíritum rectum ínnova in viscéribus meis. \V\ Ne projícias me a fácie tua. \R\ Et spíritum sanctum tuum ne áuferas a me. \V\ Redde mihi lætítiam salutáris tui. \R\ Et spíritu principáli confírma me.

\variantheading{In Precibus ferialibus tantum}

\V\ Eripe me, Dómine, ab hómine malo. \R\ A viro iníquo éripe me. \V\ Eripe me de inimícis meis, Deus meus. \R\ Et ab insurgéntibus in me líbera me. \V\ Eripe me de operántibus iniquitátem. \R\ Et de viris sánguinum salva me. \V\ Sic psalmum dicam nómini tuo in séculum séculi. \R\ Ut reddam vota mea de die in diem. \V\ Exáudi nos, Deus, salutáris noster. \R\ Spes ómnium fínium terræ, et in mari longe. \V\ Deus, in adjutórium meum inténde. \R\ Dómine, ad adjuvándum me festína. \V\ Sanctus Deus, Sanctus fortis, Sanctus immortális. \R\ Miserére nobis. \V\ Bénedic, ánima mea, Dómino. \R\ Et ómnia, quæ intra me sunt, nómini sancto ejus. \V\ Bénedic, ánima mea, Dómino. \R\ Et noli oblivísci omnes retributiónes ejus.

\V\ Qui propitiátur ómnibus iniquitátibus tuis. \R\ Qui sanat omnes infirmitátes tuas. \V\ Qui rédimit de intéritu vitam tuam. \R\ Qui corónat te in misericórdia et miseratiónibus. \V\ Qui replet in bonis desidérium tuum. \R\ Renovábitur ut áquilæ juvéntus tua.

\variantheading{In utrisque Precibus}

\V\ Adjutórium nostrum in nómine Dómini. \R\ Qui fecit cælum et terram.

\formulalabel{Confíteor.} Misereátur. Indulgéntiam.

\V\ Dignáre, Dómine, die isto. \R\ Sine peccáto nos custodíre. \V\ Miserére nostri, Dómine. \R\ Miserére nostri. \V\ Fiat misericórdia tua, Dómine, super nos. \R\ Quemádmodum sperávimus in te. \V\ Dómine, exáudi oratiónem meam. \R\ Et clamor meus ad te véniat.

\sectionornament

\V\ Dóminus vobíscum.

Orémus. \oratiolabel{Oratio}

\lettrine[lines=2]{D}{OMINE}, Deus omnípotens, qui ad princípium hujus diéi nos perveníre fecísti: tua nos hódie salva virtúte; ut in hac die ad nullum declinémus peccátum, sed semper ad tuam justítiam faciéndam nostra procédant elóquia, dirigántur cogitatiónes et ópera. Per Dóminum. \R\ Amen.

\V\ Dóminus vobíscum. \V\ Benedicámus Dómino.

\rubric{Deinde in Choro legitur Martyrologium. Postea Hebdomadarius dicit:}

\V\ Pretiósa in conspéctu Dómini. \R\ Mors Sanctórum ejus.

\rubric{Et absolute, sine Orémus:}

\oratiolabel{Oratio}

\lettrine[lines=2]{S}{ANCTA} María, et omnes Sancti intercédant pro nobis ad Dóminum, ut nos mereámur ab eo adjuvári et salvári, qui vivit et regnat in sécula sæculórum. \R\ Amen.

\V\ Deus, in adjutórium meum inténde. \R\ Dómine, ad adjuvándum me festína.

\rubric{Et dicitur ter; ultimo additur:} Glória Patri. Sicut erat.

Kýrie, eléison. Christe, eléison. Kýrie, eléison.

Pater noster \rubric{secreto usque ad} \V\ Et ne nos indúcas in tentatiónem. \R\ Sed líbera nos a malo.

\V\ Réspice in servos tuos, Dómine, et in ópera tua, et dírige fílios eórum. \R\ Et sit splendor Dómini, Dei nostri, super nos, et ópera mánuum nostrárum dírige super nos, et opus mánuum nostrárum dírige. \V\ Glória Patri. \R\ Sicut erat.

Orémus. \oratiolabel{Oratio}

\lettrine[lines=2]{D}{IRIGERE} et sanctificáre, régere et gubernáre dignáre, Dómine Deus, Rex cæli et terræ, hódie corda et córpora nostra, sensus, sermónes et actus nostros in lege tua, et in opéribus mandatórum tuórum: ut hic, et in ætérnum, te auxiliánte, salvi et líberi esse mereámur, Salvátor mundi: Qui vivis et regnas in sécula sæculórum. \R\ Amen.

\V\ Jube, domne.

\benedlabel{Bened.} Dies et actus nostros in sua pace dispónat Dóminus omnípotens. \R\ Amen.

\variantheading{In Officio feriali}

\variantheading{Tempore Adventus}

\lettrine[lines=2]{D}{OMINE}, miserére nostri: te enim exspectávimus: esto bráchium nostrum in mane, et salus nostra in témpore tribulatiónis. Tu autem.

\variantheading{Per Annum}

\lettrine[lines=2]{D}{OMINUS} autem dírigat corda et córpora nostra in caritáte Dei, et patiéntia Christi.

\sectionheading{E Proprio de Tempore}

\feastheading{Pro Octava S. Stephani}

\lettrine[lines=2]{P}{OSITIS} autem génibus, clamávit voce magna, dicens: Dómine, ne státuas illis hoc peccátum. Et cum hoc dixísset, obdormívit in Dómino.

\feastheading{Pro Octava S. Joannis}

\lettrine[lines=2]{I}{N} médio ecclésiæ apéruit os ejus, et implévit illum Dóminus spíritu sapiéntiæ et intelléctus, et stolam glóriæ índuit eum.

\feastheading{Pro Oct. Ss. Innocentium}

\lettrine[lines=2]{H}{I} empti sunt ex homínibus primítiæ Deo et Agno, et in ore ipsórum non est invéntum mendácium: sine mácula enim sunt ante thronum Dei.

\sectionheading{E Proprio Sanctorum}

\feastheading{Pro Conceptione Immac. B. M. V.}

\lettrine[lines=2]{S}{IGNUM} magnum appáruit in cælo: Múlier amícta sole, et luna sub pédibus ejus, et in cápite ejus coróna stellárum duódecim.

\feastheading{Pro S. Agnete et S. Agatha}

\lettrine[lines=2]{L}{AUDABIT} usque ad mortem ánima mea Dóminum, quóniam éruis sustinéntes te, et líberas eos de manu angústiæ, Dómine, Deus noster.

\sectionheading{E Communi Sanctorum}

\feastheading{Pro Apostolis}

\lettrine[lines=2]{I}{BANT} Apóstoli gaudéntes a conspéctu concílii, quóniam digni hábiti sunt pro nómine Jesu contuméliam pati.

\feastheading{Pro Martyre}

\lettrine[lines=2]{J}{USTUS} cor suum trádidit ad vigilándum dilúculo ad Dóminum, qui fecit illum, et in conspéctu Altíssimi deprecábitur.

\feastheading{Pro pluribus Martyribus}

\lettrine[lines=2]{F}{ULGEBUNT} justi, et tamquam scintíllæ in arundinéto discúrrent. Judicábunt natiónes, et dominabúntur pópulis: et regnábit Dóminus illórum in perpétuum.

\feastheading{Pro Conf. Pont.}

\lettrine[lines=2]{F}{UNGI} sacerdótio, et habére laudem in nómine ipsíus, et offérre illi incénsum dignum in odórem suavitátis.

\feastheading{Pro Conf. non Pont.}

\lettrine[lines=2]{J}{USTUM} dedúxit Dóminus per vias rectas, et osténdit illi regnum Dei, et dedit illi sciéntiam sanctórum: honestávit illum in labóribus, et complévit labóres illíus.

\feastheading{Pro Virgine Martyre}

\lettrine[lines=2]{D}{OMINE}, Deus meus, exaltásti super terram habitatiónem meam, et pro morte defluénte deprecáta sum. Invocávi Dóminum, Patrem Dómini mei, ut non derelínquat me in die tribulatiónis meæ, et in témpore superbórum sine adjutório.

\feastheading{Pro Virgine non Martyre}

\lettrine[lines=2]{O}{} \textsc{quam} pulchra est casta generátio cum claritáte! immortális est enim memória illíus: quóniam apud Deum nota est, et apud hómines.

\feastheading{Pro Martyre tantum}

\lettrine[lines=2]{L}{AUDABIT} usque ad mortem ánima mea Dóminum, quóniam éruis sustinéntes te, et líberas eos de manu angústiæ, Dómine, Deus noster.

\feastheading{Pro nec Virg. nec Mart.}

\lettrine[lines=2]{M}{ULTÆ} fíliæ congregavérunt divítias: tu supergréssa es univérsas. Fallax grátia, et vana est pulchritúdo: múlier timens Dóminum, ipsa laudábitur.

\feastheading{Pro Dedicatione Ecclesiæ}

\lettrine[lines=2]{E}{T} abstérget Deus omnem lácrimam ab óculis eórum: et mors ultra non erit, neque luctus, neque clamor, neque dolor erit ultra, quia prima abiérunt. Et dixit qui sedébat in throno: Ecce nova fácio ómnia.

\feastheading{In Officiis B. Mariæ V.}

\lettrine[lines=2]{I}{N} platéis sicut cinnamómum et bálsamum aromatízans odórem dedi: quasi myrrha elécta dedi suavitátem odóris.

\sectionornament

\V\ Adjutórium nostrum in nómine Dómini. \R\ Qui fecit cælum et terram.

\V\ Benedícite. \R\ Deus.

\benedictiolabel{Benedictio} Dóminus nos benedícat, et ab omni malo deféndat, et ad vitam perdúcat ætérnam. Et fidélium ánimæ per misericórdiam Dei requiéscant in pace. \R\ Amen.

\rubric{Deinde dicitur tantum Pater noster secreto, nisi sequatur alia Hora.}

\sectionornament

\hourheading{Ad Tertiam}

Pater noster et Ave.

\V\ Deus, in adjutórium.

\hymnusheading{Hymnus}

\lettrine[lines=2]{N}{UNC}, Sancte, nobis, Spíritus,\\
Unum Patri cum Fílio,\\
Dignáre promptus íngeri\\
Nostro refúsus péctori.

Os, lingua, mens, sensus, vigor\\
Confessiónem pérsonent,\\
Flamméscat igne cáritas,\\
Accéndat ardor próximos.

* Præsta, Pater piíssime,\\
Patríque compar Unice,\\
Cum Spíritu Paráclito\\
Regnans per omne séculum.\\
Amen.

\rubric{Per Annum:} \antlabel{Ant.} Illuminátio mea.

\variantheading{In Officio feriali Adventus}

\rubric{Feria II post Dominicam I:} \antlabel{Ant.} Jucundáre.

\rubric{Feria II post Dominicam II:} \antlabel{Ant.} Urbs fortitúdinis.

\rubric{Fer. II ante Vig. Nat. Dom.:} \antlabel{Ant.} Cum vénerit.

\psalmheading{Psalmus 26, i}

\lettrine[lines=2]{D}{OMINUS} illuminátio mea, et salus mea, * quem timébo?

Dóminus protéctor vitæ meæ, * a quo trepidábo?

Dum apprópiant super me nocéntes, * ut edant carnes meas:

Qui tríbulant me inimíci mei, * ipsi infirmáti sunt, et cecidérunt.

Si consístant advérsum me castra, * non timébit cor meum.

Si exsúrgat advérsum me prǽlium, * in hoc ego sperábo.

Unam pétii a Dómino, hanc requíram, * ut inhábitem in domo Dómini ómnibus diébus vitæ meæ:

Ut vídeam voluptátem Dómini, * et vísitem templum ejus.

Quóniam abscóndit me in tabernáculo suo: * in die malórum protéxit me in abscóndito tabernáculi sui.

In petra exaltávit me: * et nunc exaltávit caput meum super inimícos meos.

Circuívi, et immolávi in tabernáculo ejus hóstiam vociferatiónis: * cantábo, et psalmum dicam Dómino.

\psalmheading{Psalmus 26, ii}

\lettrine[lines=2]{E}{XAUDI}, Dómine, vocem meam, qua clamávi ad te: * miserére mei, et exáudi me.

Tibi dixit cor meum, exquisívit te fácies mea: * fáciem tuam, Dómine, requíram.

Ne avértas fáciem tuam a me: * ne declínes in ira a servo tuo.

Adjútor meus esto: * ne derelínquas me, neque despícias me, Deus, salutáris meus.

Quóniam pater meus, et mater mea dereliquérunt me: * Dóminus autem assúmpsit me.

Legem pone mihi, Dómine, in via tua: * et dírige me in sémitam rectam propter inimícos meos.

Ne tradíderis me in ánimas tribulántium me: * quóniam insurrexérunt in me testes iníqui, et mentíta est iníquitas sibi.

Credo vidére bona Dómini * in terra vivéntium.

Exspécta Dóminum, viríliter age: * et confortétur cor tuum, et sústine Dóminum.

\psalmheading{Psalmus 27}

\lettrine[lines=2]{A}{D} te, Dómine, clamábo, Deus meus, ne síleas a me: * nequándo táceas a me, et assimilábor descendéntibus in lacum.

Exáudi, Dómine, vocem deprecatiónis meæ dum oro ad te: * dum extóllo manus meas ad templum sanctum tuum.

Ne simul trahas me cum peccatóribus: * et cum operántibus iniquitátem ne perdas me.

Qui loquúntur pacem cum próximo suo: * mala autem in córdibus eórum.

Da illis secúndum ópera eórum, * et secúndum nequítiam adinventiónum ipsórum.

Secúndum ópera mánuum eórum tríbue illis: * redde retributiónem eórum ipsis.

Quóniam non intellexérunt ópera Dómini, et in ópera mánuum ejus * déstrues illos, et non ædificábis eos.

Benedíctus Dóminus: * quóniam exaudívit vocem deprecatiónis meæ.

Dóminus adjútor meus, et protéctor meus: * in ipso sperávit cor meum, et adjútus sum.

Et reflóruit caro mea: * ex voluntáte mea confitébor ei.

Dóminus fortitúdo plebis suæ: * et protéctor salvatiónum Christi sui est.

Salvum fac pópulum tuum, Dómine, et bénedic hereditáti tuæ: * et rege eos, et extólle illos usque in ætérnum.

\rubric{Per Annum:} \antlabel{Ant.} Illuminátio mea, et salus mea Dóminus.

\variantheading{In Officio feriali Adventus}

\rubric{Feria II post Dominicam I:} \antlabel{Ant.} Jucundáre, fília Sion, et exsúlta satis, fília Jerúsalem, allelúja.

\rubric{Feria II post Dominicam II:} \antlabel{Ant.} Urbs fortitúdinis nostræ Sion, Salvátor ponétur in ea murus et antemurále: aperíte portas, quia nobíscum Deus, allelúja.

\rubric{Fer. II ante Vig. Nat. Dom.:} \antlabel{Ant.} Cum vénerit Fílius hóminis, putas invéniet fidem super terram?

\capitulumlabel{Capitulum}{Jer. 23, 5}

\lettrine[lines=2]{E}{CCE}, dies véniunt, dicit Dóminus, et suscitábo David germen justum: et regnábit rex, et sápiens erit: et fáciet judícium et justítiam in terra.

\R\ br. Veni ad liberándum nos, * Dómine, Deus virtútum. Veni. \V\ Osténde fáciem tuam, et salvi érimus. Dómine. Glória Patri. Veni. \V\ Timébunt gentes nomen tuum, Dómine. \R\ Et omnes reges terræ glóriam tuam.

\variantheading{Preces feriales}

KYRIE, eléison. Christe, eléison. Kýrie, eléison.

Pater noster \rubric{secreto usque ad} \V\ Et ne nos indúcas in tentatiónem. \R\ Sed líbera nos a malo. \V\ Dómine, Deus virtútum, convérte nos. \R\ Et osténde fáciem tuam, et salvi érimus. \V\ Exsúrge, Christe, ádjuva nos. \R\ Et líbera nos propter nomen tuum. \V\ Dómine, exáudi oratiónem meam. \R\ Et clamor meus ad te véniat.

\variantheading{Per Annum}

\capitulumlabel{Capitulum}{Jer. 17, 14}

\lettrine[lines=2]{S}{ANA} me, Dómine, et sanábor: salvum me fac, et salvus ero: quóniam laus mea tu es.

\R\ br. Sana ánimam meam, * Quia peccávi tibi. Sana. \V\ Ego dixi: Dómine, miserére mei. Quia. Glória Patri. Sana. \V\ Adjútor meus esto, ne derelínquas me. \R\ Neque despícias me, Deus, salutáris meus.

\rubric{In Vigiliis communibus Preces feriales, ut supra.}

\rubric{Reliqua ut in Ordinario. 22}

\sectionornament

\hourheading{Ad Sextam}

Pater noster et Ave.

\V\ Deus, in adjutórium.

\hymnusheading{Hymnus}

\lettrine[lines=2]{R}{ECTOR} potens, verax Deus,\\
Qui témperas rerum vices,\\
Splendóre mane illúminas,\\
Et ígnibus merídiem:

Exstíngue flammas lítium,\\
Aufer calórem nóxium,\\
Confer salútem córporum,\\
Verámque pacem córdium.

* Præsta, Pater piíssime,\\
Patríque compar Unice,\\
Cum Spíritu Paráclito\\
Regnans per omne séculum.\\
Amen.

\rubric{Per Annum:} \antlabel{Ant.} In tua justítia.

\variantheading{In Officio feriali Adventus}

\rubric{Feria II post Dominicam I:} \antlabel{Ant.} Ecce Dóminus véniet.

\rubric{Feria II post Dominicam II:} \antlabel{Ant.} Ecce apparébit.

\rubric{Fer. II ante Vig. Nat. Dom.:} \antlabel{Ant.} Ecce jam venit.

\psalmheading{Psalmus 30, i}

\lettrine[lines=2]{I}{N} te, Dómine, sperávi, non confúndar in ætérnum: * in justítia tua líbera me.

Inclína ad me aurem tuam, * accélera ut éruas me.

Esto mihi in Deum protectórem, et in domum refúgii: * ut salvum me fácias.

Quóniam fortitúdo mea, et refúgium meum es tu: * et propter nomen tuum dedúces me, et enútries me.

Edúces me de láqueo hoc, quem abscondérunt mihi: * quóniam tu es protéctor meus.

In manus tuas comméndo spíritum meum: * redemísti me, Dómine, Deus veritátis.

Odísti observántes vanitátes, * supervácue.

Ego autem in Dómino sperávi: * exsultábo, et lætábor in misericórdia tua.

Quóniam respexísti humilitátem meam, * salvásti de necessitátibus ánimam meam.

Nec conclusísti me in mánibus inimíci: * statuísti in loco spatióso pedes meos.

\psalmheading{Psalmus 30, ii}

\lettrine[lines=2]{M}{ISERERE} mei, Dómine, quóniam tríbulor: * conturbátus est in ira óculus meus, ánima mea, et venter meus:

Quóniam defécit in dolóre vita mea: * et anni mei in gemítibus.

Infirmáta est in paupertáte virtus mea: * et ossa mea conturbáta sunt.

Super omnes inimícos meos factus sum oppróbrium et vicínis meis valde: * et timor notis meis.

Qui vidébant me, foras fugérunt a me: * oblivióni datus sum, tamquam mórtuus a corde.

Factus sum tamquam vas pérditum: * quóniam audívi vituperatiónem multórum commorántium in circúitu.

In eo dum convenírent simul advérsum me, * accípere ánimam meam consiliáti sunt.

Ego autem in te sperávi, Dómine: * dixi: Deus meus es tu: in mánibus tuis sortes meæ.

Eripe me de manu inimicórum meórum, * et a persequéntibus me.

Illústra fáciem tuam super servum tuum, salvum me fac in misericórdia tua: * Dómine, non confúndar, quóniam invocávi te.

Erubéscant ímpii, et deducántur in inférnum: * muta fiant lábia dolósa.

Quæ loquúntur advérsus justum iniquitátem, * in supérbia, et in abusióne.

\psalmheading{Psalmus 30, iii}

\lettrine[lines=2]{Q}{UAM} magna multitúdo dulcédinis tuæ, Dómine, * quam abscondísti timéntibus te.

Perfecísti eis, qui sperant in te, * in conspéctu filiórum hóminum.

Abscóndes eos in abscóndito faciéi tuæ, * a conturbatióne hóminum.

Próteges eos in tabernáculo tuo * a contradictióne linguárum.

Benedíctus Dóminus: * quóniam mirificávit misericórdiam suam mihi in civitáte muníta.

Ego autem dixi in excéssu mentis meæ: * Projéctus sum a fácie oculórum tuórum.

Ideo exaudísti vocem oratiónis meæ, * dum clamárem ad te.

Dilígite Dóminum, omnes sancti ejus: * quóniam veritátem requíret Dóminus, et retríbuet abundánter faciéntibus supérbiam.

Viríliter ágite, et confortétur cor vestrum, * omnes, qui sperátis in Dómino.

\rubric{Per Annum:} \antlabel{Ant.} In tua justítia líbera me, Dómine.

\variantheading{In Officio feriali Adventus}

\rubric{Feria II post Dominicam I:} \antlabel{Ant.} Ecce Dóminus véniet, et omnes Sancti ejus cum eo: et erit in die illa lux magna, allelúja.

\rubric{Feria II post Dominicam II:} \antlabel{Ant.} Ecce apparébit Dóminus, et non mentiétur: si moram fécerit, exspécta eum, quia véniet, et non tardábit, allelúja.

\rubric{Fer. II ante Vig. Nat. Dom.:} \antlabel{Ant.} Ecce jam venit plenitúdo témporis, in quo misit Deus Fílium suum in terras.

\capitulumlabel{Capitulum}{Jer. 23, 6}

\lettrine[lines=2]{I}{N} diébus illis salvábitur Juda, et Israël habitábit confidénter: et hoc est nomen, quod vocábunt eum, Dóminus justus noster.

\R\ br. Osténde nobis, Dómine, * Misericórdiam tuam. Osténde. \V\ Et salutáre tuum da nobis. Misericórdiam tuam. Glória Patri. Osténde. \V\ Meménto nostri, Dómine, in beneplácito pópuli tui. \R\ Vísita nos in salutári tuo.

\variantheading{Preces feriales}

KYRIE, eléison. Christe, eléison. Kýrie, eléison.

Pater noster \rubric{secreto usque ad} \V\ Et ne nos indúcas in tentatiónem. \R\ Sed líbera nos a malo. \V\ Dómine, Deus virtútum, convérte nos. \R\ Et osténde fáciem tuam, et salvi érimus. \V\ Exsúrge, Christe, ádjuva nos. \R\ Et líbera nos propter nomen tuum. \V\ Dómine, exáudi oratiónem meam. \R\ Et clamor meus ad te véniat.

\variantheading{Per Annum}

\capitulumlabel{Capitulum}{Rom. 13, 8}

\lettrine[lines=2]{N}{EMINI} quidquam debeátis, nisi ut ínvicem diligátis: qui enim díligit próximum, legem implévit.

\R\ br. Benedícam Dóminum * In omni témpore. Benedícam. \V\ Semper laus ejus in ore meo. In. Glória Patri. Benedícam. \V\ Dóminus regit me, et nihil mihi déerit. \R\ In loco páscuæ ibi me collocávit.

\rubric{In Vigiliis communibus Preces feriales, ut supra.}

\rubric{Reliqua ut in Ordinario. 24}

\sectionornament

\hourheading{Ad Nonam}

Pater noster et Ave.

\V\ Deus, in adjutórium.

\hymnusheading{Hymnus}

\lettrine[lines=2]{R}{ERUM}, Deus, tenax vigor,\\
Immótus in te pérmanens,\\
Lucis diúrnæ témpora\\
Succéssibus detérminans:

Largíre lumen véspere,\\
Quo vita nusquam décidat,\\
Sed prǽmium mortis sacræ\\
Perénnis instet glória.

* Præsta, Pater piíssime,\\
Patríque compar Unice,\\
Cum Spíritu Paráclito\\
Regnans per omne séculum.\\
Amen.

\rubric{Per Annum:} \antlabel{Ant.} Exsultáte, justi.

\variantheading{In Officio feriali Adventus}

\rubric{Feria II post Dominicam I:} \antlabel{Ant.} Ecce véniet.

\rubric{Feria II post Dominicam II:} \antlabel{Ant.} Ecce Dóminus.

\rubric{Fer. II ante Vig. Nat. Dom.:} \antlabel{Ant.} Egrediétur.

\psalmheading{Psalmus 31}

\lettrine[lines=2]{B}{EATI} quorum remíssæ sunt iniquitátes: * et quorum tecta sunt peccáta.

Beátus vir, cui non imputávit Dóminus peccátum, * nec est in spíritu ejus dolus.

Quóniam tácui, inveteravérunt ossa mea, * dum clamárem tota die.

Quóniam die ac nocte graváta est super me manus tua: * convérsus sum in ærúmna mea, dum confígitur spina.

Delíctum meum cógnitum tibi feci: * et injustítiam meam non abscóndi.

Dixi: Confitébor advérsum me injustítiam meam Dómino: * et tu remisísti impietátem peccáti mei.

Pro hac orábit ad te omnis sanctus, * in témpore opportúno.

Verúmtamen in dilúvio aquárum multárum, * ad eum non approximábunt.

Tu es refúgium meum a tribulatióne, quæ circúmdedit me: * exsultátio mea, érue me a circumdántibus me.

Intelléctum tibi dabo, et ínstruam te in via hac, qua gradiéris: * firmábo super te óculos meos.

Nolíte fíeri sicut equus et mulus, * quibus non est intelléctus.

In camo et freno maxíllas eórum constrínge, * qui non appróximant ad te.

Multa flagélla peccatóris, * sperántem autem in Dómino misericórdia circúmdabit.

Lætámini in Dómino et exsultáte, justi, * et gloriámini, omnes recti corde.

\psalmheading{Psalmus 32, i}

\lettrine[lines=2]{E}{XSULTATE}, justi, in Dómino: * rectos decet collaudátio.

Confitémini Dómino in cíthara: * in psaltério decem chordárum psállite illi.

Cantáte ei cánticum novum: * bene psállite ei in vociferatióne.

Quia rectum est verbum Dómini, * et ómnia ópera ejus in fide.

Díligit misericórdiam et judícium: * misericórdia Dómini plena est terra.

Verbo Dómini cæli firmáti sunt: * et spíritu oris ejus omnis virtus eórum.

Cóngregans sicut in utre aquas maris: * ponens in thesáuris abýssos.

Tímeat Dóminum omnis terra: * ab eo autem commoveántur omnes inhabitántes orbem.

Quóniam ipse dixit, et facta sunt: * ipse mandávit, et creáta sunt.

Dóminus díssipat consília géntium: * réprobat autem cogitatiónes populórum, et réprobat consília príncipum.

Consílium autem Dómini in ætérnum manet: * cogitatiónes cordis ejus in generatióne et generatiónem.

\psalmheading{Psalmus 32, ii}

\lettrine[lines=2]{B}{EATA} gens, cujus est Dóminus, Deus ejus: * pópulus, quem elégit in hereditátem sibi.

De cælo respéxit Dóminus: * vidit omnes fílios hóminum.

De præparáto habitáculo suo * respéxit super omnes, qui hábitant terram.

Qui finxit sigillátim corda eórum: * qui intélligit ómnia ópera eórum.

Non salvábitur rex per multam virtútem: * et gigas non salvábitur in multitúdine virtútis suæ.

Fallax equus ad salútem: * in abundántia autem virtútis suæ non salvábitur.

Ecce óculi Dómini super metuéntes eum: * et in eis, qui sperant super misericórdia ejus:

Ut éruat a morte ánimas eórum: * et alat eos in fame.

Anima nostra sústinet Dóminum: * quóniam adjútor et protéctor noster est.

Quia in eo lætábitur cor nostrum: * et in nómine sancto ejus sperávimus.

Fiat misericórdia tua, Dómine, super nos: * quemádmodum sperávimus in te.

\rubric{Per Annum:} \antlabel{Ant.} Exsultáte, justi, et gloriámini, omnes recti corde.

\variantheading{In Officio feriali Adventus}

\rubric{Feria II post Dominicam I:} \antlabel{Ant.} Ecce véniet Prophéta magnus, et ipse renovábit Jerúsalem, allelúja.

\rubric{Feria II post Dominicam II:} \antlabel{Ant.} Ecce Dóminus noster cum virtúte véniet, et illuminábit óculos servórum suórum, allelúja.

\rubric{Fer. II ante Vig. Nat. Dom.:} \antlabel{Ant.} Egrediétur Dóminus de loco sancto suo: véniet, ut salvet pópulum suum.

\capitulumlabel{Capitulum}{Is. 14, 1}

\lettrine[lines=2]{P}{ROPE} est ut véniat tempus ejus, et dies ejus non elongabúntur: miserébitur enim Dóminus Jacob, et Israël salvábitur.

\R\ br. Super te, Jerúsalem, * Oriétur Dóminus. Super te. \V\ Et glória ejus in te vidébitur. Oriétur Dóminus. Glória Patri. Super te. \V\ Veni, Dómine, et noli tardáre. \R\ Reláxa facínora plebi tuæ.

\variantheading{Preces feriales}
\label{feria2:precesferialesnona}

KYRIE, eléison. Christe, eléison. Kýrie, eléison.

Pater noster \rubric{secreto usque ad} \V\ Et ne nos indúcas in tentatiónem. \R\ Sed líbera nos a malo. \V\ Dómine, Deus virtútum, convérte nos. \R\ Et osténde fáciem tuam, et salvi érimus. \V\ Exsúrge, Christe, ádjuva nos. \R\ Et líbera nos propter nomen tuum. \V\ Dómine, exáudi oratiónem meam. \R\ Et clamor meus ad te véniat.

\variantheading{Per Annum}

\capitulumlabel{Capit.}{1 Petri 1, 17-19}

\lettrine[lines=2]{I}{N} timóre incolátus vestri témpore conversámini: sciéntes quod non corruptibílibus auro vel argénto redémpti estis, sed pretióso sánguine quasi Agni immaculáti Christi.

\R\ br. Rédime me, Dómine, * Et miserére mei. Rédime. \V\ Pes enim meus stetit in via recta. Et miserére mei. Glória Patri. Rédime. \V\ Ab occúltis meis munda me, Dómine. \R\ Et ab aliénis parce servo tuo.

\rubric{In Vigiliis communibus Preces feriales, ut supra.} (p.~\pageref{feria2:precesferialesnona})

\rubric{Reliqua ut in Ordinario. 25}

\sectionornament

\hourheading{Ad Vesperas}

Pater noster et Ave.

\V\ Deus, in adjutórium.

\antlabel{Ant.} Inclinávit Dóminus * aurem suam mihi.

\rubric{¶ Antiphonæ, in hac et aliis Feriis appositæ, dicuntur semper in Officio feriali et in Festis sive trium sive novem Lectionum, quando juxta Rubricas sumendi sint Psalmi de Feria, ut in Psalterio.}

\psalmheading{Psalmus 114}

\lettrine[lines=2]{D}{ILEXI}, quóniam exáudiet Dóminus * vocem oratiónis meæ.

Quia inclinávit aurem suam mihi: * et in diébus meis invocábo.

Circumdedérunt me dolóres mortis: * et perícula inférni invenérunt me.

Tribulatiónem et dolórem invéni: * et nomen Dómini invocávi.

O Dómine, líbera ánimam meam: * miséricors Dóminus, et justus, et Deus noster miserétur.

Custódiens párvulos Dóminus: * humiliátus sum, et liberávit me.

Convértere, ánima mea, in réquiem tuam: * quia Dóminus benefécit tibi.

Quia erípuit ánimam meam de morte: * óculos meos a lácrimis, pedes meos a lapsu.

Placébo Dómino * in regióne vivórum.

\antlabel{Ant.} Inclinávit Dóminus aurem suam mihi.

\antlabel{Ant.} Vota mea * Dómino reddam coram omni pópulo ejus.

\psalmheading{Psalmus 115}

\lettrine[lines=2]{C}{REDIDI}, propter quod locútus sum: * ego autem humiliátus sum nimis.

Ego dixi in excéssu meo: * Omnis homo mendax.

Quid retríbuam Dómino, * pro ómnibus, quæ retríbuit mihi?

Cálicem salutáris accípiam: * et nomen Dómini invocábo.

Vota mea Dómino reddam coram omni pópulo ejus: * pretiósa in conspéctu Dómini mors sanctórum ejus:

O Dómine, quia ego servus tuus: * ego servus tuus, et fílius ancíllæ tuæ.

Dirupísti víncula mea: * tibi sacrificábo hóstiam laudis, et nomen Dómini invocábo.

Vota mea Dómino reddam in conspéctu omnis pópuli ejus: * in átriis domus Dómini, in médio tui, Jerúsalem.

\antlabel{Ant.} Vota mea Dómino reddam coram omni pópulo ejus.

\antlabel{Ant.} Clamávi, * et Dóminus exaudívit me.

\psalmheading{Psalmus 119}

\lettrine[lines=2]{A}{D} Dóminum cum tribulárer clamávi: * et exaudívit me.

Dómine, líbera ánimam meam a lábiis iníquis, * et a lingua dolósa.

Quid detur tibi, aut quid apponátur tibi * ad linguam dolósam?

Sagíttæ poténtis acútæ, * cum carbónibus desolatóriis.

Heu mihi, quia incolátus meus prolongátus est: * habitávi cum habitántibus Cedar: multum íncola fuit ánima mea.

Cum his, qui odérunt pacem, eram pacíficus: * cum loquébar illis, impugnábant me gratis.

\antlabel{Ant.} Clamávi, et Dóminus exaudívit me.

\antlabel{Ant.} Auxílium meum * a Dómino qui fecit cælum et terram.

\psalmheading{Psalmus 120}

\lettrine[lines=2]{L}{EVAVI} óculos meos in montes, * unde véniet auxílium mihi.

Auxílium meum a Dómino, * qui fecit cælum et terram.

Non det in commotiónem pedem tuum: * neque dormítet qui custódit te.

Ecce, non dormitábit neque dórmiet, * qui custódit Israël.

Dóminus custódit te, Dóminus protéctio tua, * super mánum déxteram tuam.

Per diem sol non uret te: * neque luna per noctem.

Dóminus custódit te ab omni malo: * custódiat ánimam tuam Dóminus.

Dóminus custódiat intróitum tuum, et éxitum tuum: * ex hoc nunc, et usque in séculum.

\antlabel{Ant.} Auxílium meum a Dómino, qui fecit cælum et terram.

\antlabel{Ant.} Lætátus sum * \textdagger\ in his, quæ dicta sunt mihi. \textdaggerdbl

\psalmheading{Psalmus 121}

\lettrine[lines=2]{L}{ÆTATUS} sum \textdagger\ in his, quæ dicta sunt mihi: * \textdaggerdbl\ In domum Dómini íbimus.

Stantes erant pedes nostri, * in átriis tuis, Jerúsalem.

Jerúsalem, quæ ædificátur ut cívitas: * cujus participátio ejus in idípsum.

Illuc enim ascendérunt tribus, tribus Dómini: * testimónium Israël ad confiténdum nómini Dómini.

Quia illic sedérunt sedes in judício, * sedes super domum David.

Rogáte quæ ad pacem sunt Jerúsalem: * et abundántia diligéntibus te:

Fiat pax in virtúte tua: * et abundántia in túrribus tuis.

Propter fratres meos, et próximos meos, * loquébar pacem de te:

Propter domum Dómini, Dei nostri, * quæsívi bona tibi.

\antlabel{Ant.} Lætátus sum in his, quæ dicta sunt mihi.

\variantheading{In Officio feriali Adventus}

\capitulumlabel{Capitulum}{Gen. 49, 10}

\lettrine[lines=2]{N}{ON} auferétur sceptrum de Juda, et dux de fémore ejus, donec véniat qui mitténdus est: et ipse erit exspectátio géntium.

\hymnusheading{Hymnus}

\lettrine[lines=2]{C}{REATOR} alme síderum,\\
Ætérna lux credéntium,\\
Jesu, Redémptor ómnium,\\
Inténde votis súpplicum.

Qui dæmonis ne fráudibus\\
Períret orbis, ímpetu\\
Amóris actus, lánguidi\\
Mundi medéla factus es.

Commúne qui mundi nefas\\
Ut expiáres, ad Crucem\\
E Vírginis sacrário\\
Intácta prodis víctima.

Cujus potéstas glóriæ,\\
Noménque cum primum sonat,\\
Et Cǽlites et ínferi\\
Treménte curvántur genu.

Te deprecámur última\\
Magnum diéi Júdicem,

Armis supérnæ grátiæ\\
Defénde nos ab hóstibus.

Virtus, honor, laus, glória\\
Deo Patri cum Fílio,\\
Sancto simul Paráclito,\\
In sæculórum sǽcula.\\
Amen.

\V\ Roráte, cæli, désuper, et nubes pluant justum. \R\ Aperiátur terra, et gérminet Salvatórem.

\rubric{Antiphona ad Magnificat ut in Proprio de Tempore.}

\rubric{Expleto Cantico ac repetita Antiphona, dicuntur sequentes}

\variantheading{Preces feriales}

KYRIE, eléison. Christe, eléison. Kýrie, eléison.

Pater noster \rubric{quod dicitur a solo Hebdomadario totum clara voce, usque ad} \V\ Et ne nos indúcas in tentatiónem. \R\ Sed líbera nos a malo. \V\ Ego dixi: Dómine, miserére mei. \R\ Sana ánimam meam, quia peccávi tibi. \V\ Convértere, Dómine, úsquequo? \R\ Et deprecábilis esto super servos tuos. \V\ Fiat misericórdia tua, Dómine, super nos. \R\ Quemádmodum sperávimus in te. \V\ Sacerdótes tui induántur justítiam. \R\ Et sancti tui exsúltent.

\V\ Orémus pro beatíssimo Papa nostro N. \R\ Dóminus consérvet eum, et vivíficet eum, et beátum fáciat eum in terra, et non tradat eum in ánimam inimicórum ejus.

\V\ Orémus et pro Antístite nostro N. \R\ Stet et pascat in fortitúdine tua, Dómine, in sublimitáte nóminis tui.

\V\ Dómine, salvum fac regem. \R\ Et exáudi nos in die, qua invocavérimus te. \V\ Salvum fac pópulum tuum, Dómine, et bénedic hereditáti tuæ. \R\ Et rege eos, et extólle illos usque in ætérnum. \V\ Meménto Congregatiónis tuæ. \R\ Quam possedísti ab inítio. \V\ Fiat pax in virtúte tua. \R\ Et abundántia in túrribus tuis. \V\ Orémus pro benefactóribus nostris. \R\ Retribúere dignáre, Dómine, ómnibus, nobis bona faciéntibus propter nomen tuum, vitam ætérnam. Amen. \V\ Orémus pro fidélibus defúnctis. \R\ Réquiem ætérnam dona eis, Dómine, et lux perpétua lúceat eis. \V\ Requiéscant in pace. \R\ Amen. \V\ Pro frátribus nostris abséntibus. \R\ Salvos fac servos tuos, Deus meus, sperántes in te. \V\ Pro afflíctis et captívis. \R\ Líbera eos, Deus Israël, ex ómnibus tribulatiónibus suis. \V\ Mitte eis, Dómine, auxílium de sancto. \R\ Et de Sion tuére eos. \V\ Dómine, Deus virtútum, convérte nos. \R\ Et osténde fáciem tuam, et salvi érimus. \V\ Exsúrge, Christe, ádjuva nos. \R\ Et líbera nos propter nomen tuum. \V\ Dómine, exáudi oratiónem meam. \R\ Et clamor meus ad te véniat.

\variantheading{Per Annum}

\rubric{Capitulum, Hymnus, Versus, et Ant. ad Magnificat dicuntur quotidie de Feria currenti, ut in Psalterio hic et alibi apponuntur.}

\capitulumlabel{Capitulum}{2 Cor. 1, 3-4}

\lettrine[lines=2]{B}{ENEDICTUS} Deus, et Pater Dómini nostri Jesu Christi, Pater misericordiárum, et Deus totíus consolatiónis, qui consolátur nos in omni tribulatióne nostra.

\hymnusheading{Hymnus}

\lettrine[lines=2]{I}{MMENSE} cæli Cónditor,\\
Qui mixta ne confúnderent,\\
Aquæ fluénta dívidens,\\
Cælum dedísti límitem.

Firmans locum cæléstibus,\\
Simúlque terræ rívulis;\\
Ut unda flammas témperet,\\
Terræ solum ne díssipent.

Infúnde nunc, piíssime,\\
Donum perénnis grátiæ:\\
Fraudis novæ ne cásibus\\
Nos error átterat vetus.

Lucem fides adáugeat:\\
Sic lúminis jubar ferat:\\
Hæc vana cuncta próterat:\\
Hanc falsa nulla cómprimant.

* Præsta, Pater piíssime,\\
Patríque compar Unice,\\
Cum Spíritu Paráclito\\
Regnans per omne séculum.\\
Amen.

\V\ Dirigátur, Dómine, orátio mea. \R\ Sicut incénsum in conspéctu tuo.

\rubric{A Feria II post Dominicam Septuagesimæ usque ad Feriam VI post Cineres inclusive, in singulis Feriis Ant. ad Magnificat, ut in Proprio de Tempore; in aliis vero Feriis per Annum, Ant. ad Magnificat dicitur quotidie de Feria currenti, ut in Psalterio hic et alibi apponitur.}

\antlabel{Ad Magnif. Ant.} Magníficat * \textdagger\ ánima mea Dóminum, quia respéxit Deus humilitátem meam.

\rubric{Reliqua ut in Ordinario. 27}

\sectionornament

\hourheading{Ad Completorium}

\V\ Jube, domne.

\benedlabel{Bened.} Noctem quiétam, et finem perféctum concédat nobis Dóminus omnípotens. \R\ Amen.

\lectiolabel{Lectio br.}{1 Petri 5, 8-9}

\lettrine[lines=2]{F}{RATRES}: Sóbrii estóte, et vigiláte: quia adversárius vester diábolus tamquam leo rúgiens círcuit, quærens quem dévoret: cui resístite fortes in fide. Tu autem, Dómine, miserére nobis. \R\ Deo grátias.

\V\ Adjutórium nostrum in nómine Dómini. \R\ Qui fecit cælum et terram.

Pater noster \rubric{quod dicitur totum secreto.}

\formulalabel{Confíteor.} Misereátur. Indulgéntiam.

\V\ Convérte nos, Deus, salutáris noster. \R\ Et avérte iram tuam a nobis.

\V\ Deus, in adjutórium.

\antlabel{Ant.} Salvum me fac.

\rubric{¶ Antiphonæ, in hac et aliis Feriis appositæ, dicuntur semper ad Completórium tum de Témpore tum de Sanctis, quando juxta Rubricas sumendi sint Psalmi de Feria.}

\psalmheading{Psalmus 6}

\lettrine[lines=2]{D}{OMINE}, ne in furóre tuo árguas me, * neque in ira tua corrípias me.

Miserére mei, Dómine, quóniam infírmus sum: * sana me, Dómine, quóniam conturbáta sunt ossa mea.

Et ánima mea turbáta est valde: * sed tu, Dómine, úsquequo?

Convértere, Dómine, et éripe ánimam meam: * salvum me fac propter misericórdiam tuam.

Quóniam non est in morte qui memor sit tui: * in inférno autem quis confitébitur tibi?

Laborávi in gémitu meo, lavábo per síngulas noctes lectum meum: * lácrimis meis stratum meum rigábo.

Turbátus est a furóre óculus meus: * inveterávi inter omnes inimícos meos.

Discédite a me, omnes, qui operámini iniquitátem: * quóniam exaudívit Dóminus vocem fletus mei.

Exaudívit Dóminus deprecatiónem meam, * Dóminus oratiónem meam suscépit.

Erubéscant, et conturbéntur veheménter omnes inimíci mei: * convertántur et erubéscant valde velóciter.

\psalmheading{Psalmus 7, i}

\lettrine[lines=2]{D}{OMINE}, Deus meus, in te sperávi: * salvum me fac ex ómnibus persequéntibus me, et líbera me.

Nequándo rápiat ut leo ánimam meam, * dum non est qui rédimat, neque qui salvum fáciat.

Dómine, Deus meus, si feci istud, * si est iníquitas in mánibus meis:

Si réddidi retribuéntibus mihi mala, * décidam mérito ab inimícis meis inánis.

Persequátur inimícus ánimam meam, et comprehéndat, et concúlcet in terra vitam meam, * et glóriam meam in púlverem dedúcat.

Exsúrge, Dómine, in ira tua: * et exaltáre in fínibus inimicórum meórum.

Et exsúrge, Dómine, Deus meus, in præcépto quod mandásti: * et synagóga populórum circúmdabit te.

Et propter hanc in altum regrédere: * Dóminus júdicat pópulos.

Júdica me, Dómine, secúndum justítiam meam, * et secúndum innocéntiam meam super me.

Consumétur nequítia peccatórum, et díriges justum, * scrutans corda et renes, Deus.

\psalmheading{Psalmus 7, ii}

\lettrine[lines=2]{J}{USTUM} adjutórium meum a Dómino, * qui salvos facit rectos corde.

Deus judex justus, fortis, et pátiens: * numquid iráscitur per síngulos dies?

Nisi convérsi fuéritis, gládium suum vibrábit: * arcum suum teténdit, et parávit illum.

Et in eo parávit vasa mortis: * sagíttas suas ardéntibus effécit.

Ecce, partúriit injustítiam: * concépit dolórem, et péperit iniquitátem.

Lacum apéruit, et effódit eum: * et íncidit in fóveam quam fecit.

Convertétur dolor ejus in caput ejus: * et in vérticem ipsíus iníquitas ejus descéndet.

Confitébor Dómino secúndum justítiam ejus: * et psallam nómini Dómini altíssimi.

\antlabel{Ant.} Salvum me fac, Dómine, propter misericórdiam tuam.

\hymnusheading{Hymnus}

\lettrine[lines=2]{T}{E} lucis ante términum,\\
Rerum Creátor, póscimus,\\
Ut pro tua cleméntia\\
Sis præsul et custódia.

Procul recédant sómnia,\\
Et nóctium phantásmata,\\
Hostémque nostrum cómprime,\\
Ne polluántur córpora.

* Præsta, Pater piíssime,\\
Patríque compar Unice,\\
Cum Spíritu Paráclito\\
Regnans per omne séculum.\\
Amen.

\capitulumlabel{Capitulum}{Jer. 14, 9}

\lettrine[lines=2]{T}{U} autem in nobis es, Dómine, et nomen sanctum tuum invocátum est super nos: ne derelínquas nos, Dómine, Deus noster.

\R\ br. In manus tuas, Dómine, * Comméndo spíritum meum. In manus. \V\ Redemísti nos, Dómine, Deus veritátis. Comméndo. Glória Patri. In manus. \V\ Custódi nos, Dómine, ut pupíllam óculi. \R\ Sub umbra alárum tuárum prótege nos.

\antlabel{Ant.} Salva nos.

\canticumheading{Canticum Simeonis}{Luc. 2, 29-32}

\lettrine[lines=2]{N}{UNC} dimíttis servum tuum, Dómine, * secúndum verbum tuum in pace:

Quia vidérunt óculi mei * salutáre tuum,

Quod parásti * ante fáciem ómnium populórum,

Lumen ad revelatiónem géntium, * et glóriam plebis tuæ Israël.

Glória Patri.

\antlabel{Ant.} Salva nos, Dómine, vigilántes, custódi nos dormiéntes; ut vigilémus cum Christo, et requiescámus in pace.

\variantheading{Preces}

\textit{(si dicendæ sint)}

KYRIE, eléison. Christe, eléison. Kýrie, eléison.

Pater noster \rubric{secreto usque ad} \V\ Et ne nos indúcas in tentatiónem. \R\ Sed líbera nos a malo.

Credo in Deum \rubric{secreto usque ad} \V\ Carnis resurrectiónem. \R\ Vitam ætérnam. Amen.

\V\ Benedíctus es, Dómine, Deus patrum nostrórum. \R\ Et laudábilis et gloriósus in sécula. \V\ Benedicámus Patrem et Fílium cum Sancto Spíritu. \R\ Laudémus, et superexaltémus eum in sécula. \V\ Benedíctus es, Dómine, in firmaménto cæli. \R\ Et laudábilis, et gloriósus, et superexaltátus in sécula. \V\ Benedícat et custódiat nos omnípotens et miséricors Dóminus. \R\ Amen. \V\ Dignáre, Dómine, nocte ista. \R\ Sine peccáto nos custodíre. \V\ Miserére nostri, Dómine. \R\ Miserére nostri. \V\ Fiat misericórdia tua, Dómine, super nos. \R\ Quemádmodum sperávimus in te. \V\ Dómine, exáudi oratiónem meam. \R\ Et clamor meus ad te véniat.

\V\ Dóminus vobíscum.

Orémus. \oratiolabel{Oratio}

\lettrine[lines=2]{V}{ISITA}, quæsumus, Dómine, habitatiónem istam, et omnes insídias inimíci ab ea longe repélle: Angeli tui sancti hábitent in ea, qui nos in pace custódiant; et benedíctio tua sit super nos semper. Per Dóminum nostrum. \R\ Amen.

\V\ Dóminus vobíscum.

\V\ Benedicámus Dómino.

\benedictiolabel{Benedictio} Benedícat et custódiat nos omnípotens et miséricors Dóminus, Pater, et Fílius, et Spíritus Sanctus. \R\ Amen.

\antlabel{Antiphona finalis B. M. V.}

\V\ Divínum auxílium.

Pater, Ave et Credo.

\sectionornament

\feriaheading{Feria Tertia}

\hourheading{Ad Matutinum}

Pater, Ave et Credo.

\V\ Dómine, lábia mea.

\V\ Deus, in adjutórium.

\variantheading{In Officio feriali Adventus}

\variantheading{Infra Hebdomadam I et II:}

\invitatlabel{Invit.} Regem ventúrum Dóminum, * Veníte, adorémus.

\variantheading{Infra Hebdom. III et IV:}

\invitatlabel{Invit.} Prope est jam Dóminus: * Veníte, adorémus.

\variantheading{In Vigilia Nativitatis Domini:}

\invitatlabel{Invit.} Hódie sciétis quia véniet Dóminus: * Et mane vidébitis glóriam ejus.

\hymnusheading{Hymnus}

\lettrine[lines=2]{V}{ERBUM} supérnum, pródiens\\
E Patris ætérni sinu,\\
Qui natus orbi súbvenis,\\
Labénte cursu témporis:

Illúmina nunc péctora,\\
Tuóque amóre cóncrema;\\
Ut cor cadúca déserens\\
Cæli volúptas ímpleat.

Ut, cum tribúnal Júdicis\\
Damnábit igni nóxios,\\
Et vox amíca débitum\\
Vocábit ad cælum pios;

Non esca flammárum nigros\\
Volvámur inter túrbines,\\
Vultu Dei sed cómpotes\\
Cæli fruámur gáudiis.

Sic, Christe, nobis ómnibus\\
Indúlgeas credéntibus,\\
Ut prosit exorántibus\\
Quod præcinéntes psállimus.

* Præsta, Pater piíssime,\\
Patríque compar Unice,\\
Cum Spíritu Paráclito\\
Regnans per omne séculum.\\
Amen.

\variantheading{Per Annum}

\invitatlabel{Invitat.} Jubilémus Deo, * Salutári nostro. \psalmlabel{Ps. 94} Veníte, exsultémus.

\hymnusheading{Hymnus}

\lettrine[lines=2]{C}{ONSORS} Patérni lúminis,\\
Lux ipse lucis, et dies,\\
Noctem canéndo rúmpimus:\\
Assíste postulántibus.

Aufer tenébras méntium,\\
Fuga catérvas dæmónum,\\
Expélle somnoléntiam,\\
Ne pigritántes óbruat.

\nocturnoheading{In I Nocturno}

\antlabel{Ant.} Expúgna, Dómine, * impugnántes me.

\psalmheading{Psalmus 34, i}

\lettrine[lines=2]{J}{UDICA}, Dómine, nocéntes me, * expúgna impugnántes me.

Apprehénde arma et scutum: * exsúrge in adjutórium mihi.

Effúnde frámeam, et conclúde advérsus eos, qui persequúntur me: * dic ánimæ meæ: Salus tua ego sum.

Confundántur et revereántur, * quæréntes ánimam meam.

Avertántur retrórsum, et confundántur * cogitántes mihi mala.

Fiant tamquam pulvis ante fáciem venti: * et Angelus Dómini coárctans eos.

Fiat via illórum ténebræ et lúbricum: * et Angelus Dómini pérsequens eos.

Quóniam gratis abscondérunt mihi intéritum láquei sui: * supervácue exprobravérunt ánimam meam.

Véniat illi láqueus, quem ignórat: et cáptio, quam abscóndit, apprehéndat eum: * et in láqueum cadat in ipsum.

Anima autem mea exsultábit in Dómino: * et delectábitur super salutári suo.

Omnia ossa mea dicent: * Dómine, quis símilis tibi?

Erípiens ínopem de manu fortiórum ejus: * egénum et páuperem a diripiéntibus eum.

\antlabel{Ant.} Expúgna, Dómine, impugnántes me.

\antlabel{Ant.} Restítue ánimam meam * a malefáctis eórum, Dómine.

\psalmheading{Psalmus 34, ii}

\lettrine[lines=2]{S}{URGENTES} testes iníqui, * quæ ignorábam interrogábant me.

Retribuébant mihi mala pro bonis: * sterilitátem ánimæ meæ.

Ego autem cum mihi molésti essent, * induébar cilício.

Humiliábam in jejúnio ánimam meam: * et orátio mea in sinu meo convertétur.

Quasi próximum, et quasi fratrem nostrum, sic complacébam: * quasi lugens et contristátus, sic humiliábar.

Et advérsum me lætáti sunt, et convenérunt: * congregáta sunt super me flagélla, et ignorávi.

Dissipáti sunt, nec compúncti, tentavérunt me, subsannavérunt me subsannatióne: * frenduérunt super me déntibus suis.

Dómine, quando respícies? * restítue ánimam meam a malignitáte eórum, a leónibus únicam meam.

\antlabel{Ant.} Restítue ánimam meam a malefáctis eórum, Dómine.

\antlabel{Ant.} Exsúrge, Dómine, * et inténde judício meo.

\psalmheading{Psalmus 34, iii}

\lettrine[lines=2]{C}{ONFITEBOR} tibi in ecclésia magna, * in pópulo gravi laudábo te.

Non supergáudeant mihi qui adversántur mihi iníque: * qui odérunt me gratis et ánnuunt óculis.

Quóniam mihi quidem pacífice loquebántur: * et in iracúndia terræ loquéntes, dolos cogitábant.

Et dilatavérunt super me os suum: * dixérunt: Euge, euge, vidérunt óculi nostri.

Vidísti, Dómine, ne síleas: * Dómine, ne discédas a me.

Exsúrge et inténde judício meo: * Deus meus, et Dóminus meus in causam meam.

Júdica me secúndum justítiam tuam, Dómine, Deus meus, * et non supergáudeant mihi.

Non dicant in córdibus suis: Euge, euge, ánimæ nostræ: * nec dicant: Devorávimus eum.

Erubéscant et revereántur simul, * qui gratulántur malis meis.

Induántur confusióne et reveréntia * qui magna loquúntur super me.

Exsúltent et læténtur qui volunt justítiam meam: * et dicant semper: Magnificétur Dóminus, qui volunt pacem servi ejus.

Et lingua mea meditábitur justítiam tuam, * tota die laudem tuam.

\antlabel{Ant.} Exsúrge, Dómine, et inténde judício meo.

\variantheading{In Festis novem Lectionum:}

\V\ Lingua mea meditábitur justítiam tuam. \R\ Tota die laudem tuam, Dómine.

\absolutiolabel{Absol.} Exáudi, Dómine.

\benedlabel{Bened.} 1 Benedictióne perpétua. 2 Unigénitus Dei Fílius. 3 Spíritus Sancti grátia.

\nocturnoheading{In II Nocturno}

\antlabel{Ant.} Noli æmulári * \textdagger\ in eo, qui prosperátur, et facit iniquitátem.

\psalmheading{Psalmus 36, i}

\lettrine[lines=2]{N}{OLI} æmulári in malignántibus: * neque zeláveris faciéntes iniquitátem.

Quóniam tamquam fœnum velóciter aréscent: * et quemádmodum ólera herbárum cito décident.

Spera in Dómino, et fac bonitátem: * et inhábita terram, et pascéris in divítiis ejus.

Delectáre in Dómino: * et dabit tibi petitiónes cordis tui.

Revéla Dómino viam tuam, et spera in eo: * et ipse fáciet.

Et edúcet quasi lumen justítiam tuam: et judícium tuum tamquam merídiem: * súbditus esto Dómino, et ora eum.

Noli æmulári in eo, qui prosperátur in via sua: * in hómine faciénte injustítias.

Désine ab ira, et derelínque furórem: * noli æmulári ut malignéris.

Quóniam qui malignántur, exterminabúntur: * sustinéntes autem Dóminum, ipsi hereditábunt terram.

Et adhuc pusíllum, et non erit peccátor: * et quæres locum ejus, et non inviénies.

Mansuéti autem hereditábunt terram: * et delectabúntur in multitúdine pacis.

Observábit peccátor justum: * et stridébit super eum déntibus suis.

Dóminus autem irridébit eum: * quóniam próspicit quod véniet dies ejus.

Gládium evaginavérunt peccatóres: * intendérunt arcum suum.

Ut dejíciant páuperem et ínopem: * ut trucídent rectos corde.

Gládius eórum intret in corda ipsórum: * et arcus eórum confringátur.

\antlabel{Ant.} Noli æmulári in eo, qui prosperátur, et facit iniquitátem.

\antlabel{Ant.} Bráchia peccatórum * conteréntur: confírmat autem justos Dóminus.

\psalmheading{Psalmus 36, ii}

\lettrine[lines=2]{M}{ELIUS} est módicum justo, * super divítias peccatórum multas.

Quóniam bráchia peccatórum conteréntur: * confírmat autem justos Dóminus.

Novit Dóminus dies immaculatórum: * et heréditas eórum in ætérnum erit.

Non confundéntur in témpore malo, et in diébus famis saturabúntur: * quia peccatóres períbunt.

Inimíci vero Dómini mox ut honorificáti fúerint et exaltáti: * deficiéntes, quemádmodum fumus defícient.

Mutuábitur peccátor, et non solvet: * justus autem miserétur et tríbuet.

Quia benedicéntes ei hereditábunt terram: * maledicéntes autem ei períbunt.

Apud Dóminum gressus hóminis dirigéntur: * et viam ejus volet.

Cum cecíderit, non collidétur: * quia Dóminus suppónit manum suam.

Júnior fui, étenim sénui: * et non vidi justum derelíctum, nec semen ejus quærens panem.

Tota die miserétur et cómmodat: * et semen illíus in benedictióne erit.

Declína a malo, et fac bonum: * et inhábita in séculum séculi.

Quia Dóminus amat judícium, et non derelínquet sanctos suos: * in ætérnum conservabúntur.

Injústi puniéntur: * et semen impiórum períbit.

Justi autem hereditábunt terram: * et inhabitábunt in séculum séculi super eam.

\antlabel{Ant.} Bráchia peccatórum conteréntur: confírmat autem justos Dóminus.

\antlabel{Ant.} Custódi innocéntiam, * et vide æquitátem.

\psalmheading{Psalmus 36, iii}

\lettrine[lines=2]{O}{S} justi meditábitur sapiéntiam, * et lingua ejus loquétur judícium.
% [hiemalis.md line 5911 repeats "et lingua ejus loquétur judícium."
%  as a stray duplicate fragment not present in the Vulgate text;
%  omitted here as a transcription artifact, flagged for review]

Lex Dei ejus in corde ipsíus: * et non supplantabúntur gressus ejus.

Consíderat peccátor justum: * et quærit mortificáre eum.

Dóminus autem non derelínquet eum in mánibus ejus: * nec damnábit eum, cum judicábitur illi.

Exspécta Dóminum, et custódi viam ejus: et exaltábit te ut hereditáte cápias terram: * cum períerint peccatóres vidébis.

Vidi ímpium superexaltátum, * et elevátum sicut cedros Líbani.

Et transívi, et ecce non erat: * et quæsívi eum, et non est invéntus locus ejus.

Custódi innocéntiam, et vide æquitátem: * quóniam sunt relíquiæ hómini pacífico.

Injústi autem disperíbunt simul: * relíquiæ impiórum interíbunt.

% ===================================================================
% END OF BATCH 11 (Psalterium, Feria Secunda Completorium conclusion;
% Feria Tertia Ad Matutinum through Ps 34-36; hiemalis.md lines
% 5428-5927). Transcription artifact at md line 5911 flagged/omitted.
% Continuation at line 5928.
% ===================================================================
